\documentclass[11pt]{article}
% ============================================================================
%  preamble.tex -- shared preamble for the math-only unified_theory_math.tex
%  Packages, theorem environments, and the notation macros. Symbols are
%  identical to the source document; only the prose is domain-neutral (pure
%  math / machine learning, no application-domain interpretation).
% ============================================================================

\usepackage[T1]{fontenc}
\usepackage[utf8]{inputenc}
\usepackage{lmodern}
\usepackage[letterpaper,margin=1in]{geometry}
\usepackage{amsmath,amssymb,amsthm,mathtools}
\usepackage{bm}
\usepackage{booktabs}
\usepackage{longtable}
\usepackage{array}
\usepackage{enumitem}
\usepackage{xcolor}
\usepackage{listings}
\usepackage{graphicx}
\usepackage[colorlinks=true,linkcolor=blue!55!black,citecolor=blue!55!black,%
            urlcolor=blue!55!black]{hyperref}
\usepackage{microtype}

% ---- theorem environments --------------------------------------------------
\theoremstyle{plain}
\newtheorem{theorem}{Theorem}[section]
\newtheorem{lemma}[theorem]{Lemma}
\newtheorem{corollary}[theorem]{Corollary}
\newtheorem{proposition}[theorem]{Proposition}
\theoremstyle{definition}
\newtheorem{definition}[theorem]{Definition}
\newtheorem{remark}[theorem]{Remark}

% ---- a light box for open-question / caveat flags --------------------------
\newcommand{\caveat}[1]{\par\noindent\textbf{\textcolor{red!60!black}{Caveat.}}\ #1\par}

% ---- listings style for short inline code / pseudocode ---------------------
\lstset{basicstyle=\ttfamily\small,breaklines=true,columns=fullflexible,
        keepspaces=true,frame=none,xleftmargin=1em}

% ============================================================================
%  NOTATION MACROS  (bind: symbol -> concept, one per concept)
% ============================================================================

% -- generic operators --
\newcommand{\E}{\mathbb{E}}
\newcommand{\Reals}{\mathbb{R}}
\newcommand{\Tr}{\operatorname{Tr}}
\newcommand{\diago}{\operatorname{diag}}
\newcommand{\KLdiv}{\mathrm{KL}}
\newcommand{\GProc}{\mathcal{GP}}
\newcommand{\Poi}{\mathrm{Poisson}}
\newcommand{\Gauss}{\mathcal{N}}
\newcommand{\ELBO}{\mathcal{L}}          % ELBO  (script L; NOT the Cholesky factor)
\newcommand{\transpose}{^{\mathsf{T}}}

% -- model, kernel, structured prior --
\newcommand{\lat}{\lambda}               % latent GP function lambda(x)
\newcommand{\latz}{\lambda_{0}}          % rate offset
\newcommand{\rate}{f}                    % rate f = exp(A lambda + lambda_0)
\newcommand{\lgr}{g}                     % log-rate g = A lambda + lambda_0
\newcommand{\gain}{A}                    % gain (NOT projection a)
\newcommand{\spk}{y}                     % count observation (synonyms r, n)
\newcommand{\Kf}{K}                      % prior kernel function K(x,x')
\newcommand{\Ktil}{\widetilde{K}}        % inducing gram  K(Ztil,Ztil)
\newcommand{\Cmat}{C}                    % structured (grid) prior covariance
\newcommand{\Csm}{C_{\mathrm{smooth}}}   % RBF smoothing factor
\newcommand{\Kmag}{\mathcal{M}}          % kernel magnitude sqrt(v_x v_x')  (NOT M)
\newcommand{\kang}{\theta}               % kernel angle
\newcommand{\Jang}{J}                    % angular term J(theta)  (never a Jacobian)
\newcommand{\vmag}{v}                    % self-magnitude v_x = x^T C x + sigma_0^2
\newcommand{\loc}{\alpha}                % localization weight alpha_i (NOT diffusion alpha_t)
\newcommand{\bwid}{\beta}                % localized-region half-width (prior; NOT diffusion beta_t)
\newcommand{\smoo}{\rho}                 % smoothness SD (prior; NOT correlation)
\newcommand{\sigz}{\sigma_{0}}           % bias std (sigma_0 = exp(raw_sigma_0))
\newcommand{\rfc}{\bm{\xi}_{0}}          % center of the localized region (eps0x, eps0y)
\newcommand{\pixc}{\bm{\xi}}             % grid coordinate
\newcommand{\Amp}{\mathrm{Amp}}          % amplitude: 6th kernel hyperparameter (softplus, default 1.0, optimized)

% -- variational inference & eigenspace --
\newcommand{\ipt}{\widetilde{z}}         % inducing point (NOT subspace coord)
\newcommand{\Ipt}{\widetilde{Z}}         % inducing-point matrix
\newcommand{\latt}{\widetilde{\lambda}}  % inducing values (= u)
\newcommand{\Mind}{M}                    % number of inducing points (NOT magnitude)
\newcommand{\Ntr}{N}                     % number of training inputs
\newcommand{\kvec}{k}                    % cross-covariance vector k(x)=K(Ztil,x)
\newcommand{\pv}{u}                      % projection vector u(x)=Ktil^{-1} k(x)
\newcommand{\vm}{m}                      % variational mean (inducing space)
\newcommand{\vV}{V}                      % variational covariance
\newcommand{\pmean}{\mu}                 % posterior mean mu(x) of the latent
\newcommand{\pvar}{\sigma^{2}}           % posterior variance sigma^2(x)
\newcommand{\Lchol}{L_{\Ktil}}           % Cholesky factor of Ktil (upright L)
\newcommand{\eb}{B}                      % eigenbasis of Ktil
\newcommand{\evals}{\Lambda}             % eigenvalues (diagonal; NOT lambda)
\newcommand{\nb}{n_{b}}                  % kept eigen-dimension
\newcommand{\mb}{m_{b}}                  % eigenspace mean
\newcommand{\Vb}{V_{b}}                  % eigenspace covariance (dense)
\newcommand{\Ktilb}{\Ktil_{b}}           % diagonal inducing gram in eigenspace
\newcommand{\amat}{a}                    % projection a = K Ktil^{-1}
\newcommand{\Kvec}{\mathrm{Kvec}}        % diagonal prior variance vector
\newcommand{\fbar}{\bar{f}}              % expected rate

% -- training-step helpers --
\newcommand{\gE}{g_{\mathrm{E}}}         % E-step gradient helper (NOT log-rate g)
\newcommand{\GE}{G_{\mathrm{E}}}         % E-step Fisher curvature

% -- utility & information --
\newcommand{\Ustd}{U_{\mathrm{std}}}     % standard utility
\newcommand{\Uda}{U_{\mathrm{DA}}}       % distribution-aware utility
\newcommand{\Hmarg}{H_{\mathrm{marg}}}   % marginal count entropy
\newcommand{\Hnoise}{H_{\mathrm{noise}}} % aleatoric Poisson noise entropy
\newcommand{\Hcond}{H_{\mathrm{cond}}}   % conditional entropy over p(x)
\newcommand{\px}{p(x)}                   % input (data) distribution
\newcommand{\mug}{\mu_{g}}               % log-rate mean = A mu + lambda_0
\newcommand{\varg}{\sigma_{g}^{2}}       % log-rate variance = A^2 sigma^2
\newcommand{\corr}{\rho_{\lambda}}       % posterior latent correlation (NOT smoothness)
\newcommand{\ccond}{c}                   % prior conditional covariance
\newcommand{\scal}{\kappa}               % input scaling factor (x* = kappa x_t)
\newcommand{\varr}{\sigma_{r}^{2}}       % residual/aleatoric variance
\newcommand{\subc}{\zeta}                % subspace coordinate
\newcommand{\LW}{W_{0}}                  % Lambert-W principal branch
\newcommand{\gmode}{\bar{g}_{r}}         % Laplace mode
\newcommand{\rmax}{r_{\max}}             % count-truncation bound
\newcommand{\satf}{\sigma_{f}^{2}}       % saturation-kernel amplitude

% -- diffusion --
\newcommand{\dt}{t}                      % diffusion timestep
\newcommand{\xz}{x_{0}}                  % clean input
\newcommand{\xt}{x_{t}}                  % noisy input
\newcommand{\xhat}{\widehat{x}_{0}}      % Tweedie estimate
\newcommand{\bsched}{\beta_{t}}          % noise-schedule variance (NOT the prior beta)
\newcommand{\aret}{\alpha_{t}}           % signal retention (NOT the localization weight alpha)
\newcommand{\abar}{\bar{\alpha}_{t}}     % cumulative signal retention
\newcommand{\epsth}{\epsilon_{\theta}}   % U-Net noise predictor
\newcommand{\score}{s_{\theta}}          % score function
\newcommand{\gw}{w}                      % guidance scale
\newcommand{\ggrad}{g_{t}}               % guidance gradient (utility gradient)
\newcommand{\greg}{\gamma_{\mathrm{reg}}}% regularization weight (input-optimization approach)


\title{Active Learning with a Poisson--Gaussian-Process Model:\\[2pt]
\large Variational Inference, Information Utility, and the Diverging-Utility Problem}
\author{Mathematical reference (domain-neutral)}
\date{\today}

\begin{document}
\maketitle

\begin{abstract}
\noindent
This is a self-contained mathematical reference for a sequential active-learning
problem built on a Poisson--Gaussian-process model. A latent function $\lat$ over
an input space (inputs are real-valued signals on a bounded two-dimensional grid)
is modeled by a Gaussian process with an arc-cosine kernel over a structured prior
covariance $\Cmat$; observations are counts with a Poisson likelihood whose rate is
an exponential link of the latent; the model is fit by sparse variational
inference; and active learning chooses the next input to maximize expected
information gain about $\lat$. The document develops, in one uniform notation: the
generative model and the arc-cosine kernel with its structured prior
$\Cmat(\Amp,\bwid,\smoo,\rfc)$ (Part~I); sparse variational inference, the
eigenspace representation of the posterior, and coordinate-ascent training
(Part~II); the information utility for active learning together with the
predictive-conditioning, subspace, divergence, and numerical results beneath it
(Part~III); and guided diffusion, which extends the loop from selecting inputs out
of a pool to synthesizing them by steering a diffusion model with the utility
gradient (Part~IV). The central object is the \emph{diverging-utility problem}: the
arc-cosine kernel is positively homogeneous in the input, so the acquisition
utility grows without bound as the input scale grows and drives the optimizer to
the boundary of the bounded input box. The divergence is proved and the attempted
kernel fixes (a normalized arc-cosine kernel, a saturating arc-sin kernel) are
analyzed; finding a kernel that keeps the utility bounded on the bounded domain
while retaining good fit is left as an open problem.
\end{abstract}

\tableofcontents
\clearpage

% ==== front matter: scope, conventions, the open problem (unnumbered) ====
\section*{Scope, conventions, and the open problem}
\addcontentsline{toc}{section}{Scope, conventions, and the open problem}

This is a self-contained mathematical reference for a sequential
\emph{active-learning} problem built on a Poisson--Gaussian-process model. A
latent function $\lat(x)$ over an input space is modeled by a Gaussian process
(GP) with an arc-cosine kernel over a structured prior covariance $\Cmat$. Each
input $x\in\Reals^{d}$ is a real-valued signal on a two-dimensional grid of $d$
grid points (e.g.\ a $108\times108$ grid, $d=11664$) confined to a bounded box.
Observations are nonnegative integer counts with a Poisson likelihood whose rate
is an exponential link of the latent, $\rate=\exp(\gain\lat+\latz)$. At each round
the active-learning rule chooses the next input so as to maximize the expected
information gain about $\lat$. Part~\ref{sec:diffusion} extends the loop from
\emph{selecting} the next input out of a fixed pool to \emph{synthesizing} it with
a guided diffusion model.

\paragraph{The open problem (the purpose of this document).} The arc-cosine kernel
is positively homogeneous in the input: its self-value $\Kf(x,x)=x\transpose\Cmat
x+\sigz^{2}$ grows with $\lVert x\rVert$. Under utility maximization the posterior
variance therefore grows like the \emph{square} of the input scale, the
acquisition utility diverges, and --- because inputs are confined to a bounded box
--- the optimizer drives them to the boundary of the box (saturation). Two remedies
are presented and analyzed: a normalized arc-cosine kernel and a saturating arc-sin
kernel (\S\ref{subsec:kernelfixes}). Neither is satisfactory in practice: the
normalized kernel is scale-invariant but fits poorly. It is an \textbf{open
problem} to find a kernel that simultaneously (i) keeps the acquisition utility
bounded on the bounded input domain and (ii) retains good predictive fit. The
divergence is proved in Part~\ref{sec:deeplayer} (\S\ref{subsec:divergence}); the
attempted fixes are \S\ref{subsec:kernelfixes}.

\paragraph{Conventions.}
\begin{itemize}[leftmargin=1.4em,itemsep=1pt,topsep=2pt]
  \item \textbf{One symbol, one concept} throughout; the full table is
        Appendix~\ref{app:notation}. A few symbols that collide across the
        literature are disambiguated there (the kernel smoothness $\smoo$ vs.\ the
        posterior correlation $\corr$; the prior width $\bwid$ vs.\ the diffusion
        noise schedule $\bsched$; the three entropies $\Hmarg,\Hnoise,\Hcond$).
  \item \textbf{Counts.} The count observation is written $\spk$; where the count
        index appears (as in the marginal entropy sum and its truncation
        $\rmax$) it is written $r$, denoting the same quantity, $\spk\equiv r$.
  \item \textbf{Scope of this rendering.} Content that is purely
        implementation-specific (automatic-vs-analytic differentiation machinery,
        the vector--Jacobian-product derivation, the low-level representation and
        online-update bookkeeping, and run/status details) has been omitted; the
        mathematics of the model, inference, utility, and the divergence problem is
        complete.
\end{itemize}

\paragraph{How to read this document.} The four parts are logically ordered.
Part~\ref{sec:genmodel}--\ref{sec:kernel} sets up the generative model and the
arc-cosine kernel with its structured prior. Part~II
(\ref{sec:varinf}--\ref{sec:training}) develops sparse variational inference, the
eigenspace representation of the posterior, and the coordinate-ascent training of
the model parameters. Part~III (\ref{sec:utility}--\ref{sec:deeplayer}) is the
active-learning core: the information utility and the proofs, subspace theory, and
numerical analysis beneath it --- including the divergence result and the attempted
kernel fixes that define the open problem. Part~IV (\ref{sec:diffusion}) is guided
diffusion for input synthesis. Cross-references connect the utility gradient in
Parts~III and IV back to the kernel input-gradient of Part~I, and the training
back to the ELBO of Part~II.


\part{Foundations: the generative model and its prior}
% ============================================================================
%  Part I -- The model and its prior
%  sections/part1_model_kernel.tex   (\input-ready fragment; no preamble)
%  Covers: (1) Introduction, (2) Generative model, (3) Arc-cosine kernel +
%  structured spatial prior.
% ============================================================================

\section{Introduction: the sequential active-learning problem}
\label{sec:intro}

\subsection{Setup}

We study a sequential active-learning problem for a Poisson--Gaussian-process
model of a latent function $\lat(x)$ over inputs $x$. An input
$x\in\Reals^{d}$ is a real-valued signal on a 2-D grid of $d$ grid points (e.g. a
$108\times108$ grid, so $d=108^2=11664$); grid point $i$ sits at grid coordinate
$\pixc_i$, and the input domain is a bounded box (e.g. $[0,1]^{d}$). For a single
input $x$ the model produces a scalar latent value $\lat(x)$, an exponential-link
rate $\rate(x)=\exp(\gain\,\lat(x)+\latz)$, and a Poisson count observation
$\spk\sim\Poi(\rate(x))$ (\S\ref{sec:genmodel}). None of the formulas below assume
the particular dimension $d=11664$.

The latent function $\lat$ is unknown a priori and must be inferred from the
counts. We place a Gaussian-process (GP) prior on $\lat$ and update its posterior
online as counts arrive. The setting is \emph{sequential}: after each observation
the model is refit and used to choose the next input, so the inputs presented
depend on the observations already collected.

\subsection{The active-learning goal}

The inputs are not presented in a fixed order. At each round the next input is
chosen to be maximally informative about $\lat$ --- this is \emph{active
learning}. A candidate input $x^{*}$ is scored by an information \emph{utility}
$\Ustd(x^{*})$ (Part~III) measuring the expected reduction in uncertainty about
$\lat$ if $x^{*}$ were observed next; the input maximizing the utility over the
available pool is chosen. A second variant $\Uda$ scores information about the rate
over the input (data) distribution $\px$. Because scoring and selecting from a
fixed pool is limited by the pool, a final stage (Part~IV) instead
\emph{synthesizes} an input that maximizes the utility, using a diffusion prior
over inputs steered by the utility gradient. All three stages consume two
derivatives of the GP prior kernel that are established in this part: the input
gradient $\nabla_x\Kf$~\eqref{eq:gradxK} (utility and diffusion guidance) and the
hyperparameter gradients $\partial\Cmat/\partial\theta$~\eqref{eq:dCtheta} (model
fitting).

\subsection{Roadmap of the document}

The document proceeds in the order the pipeline runs.

\begin{itemize}[leftmargin=1.4em]
\item \textbf{Part~I (this part) --- the model and its prior.} The generative model
  (\S\ref{sec:genmodel}): a latent GP over input space, an exponential link, and
  Poisson counts. The prior kernel (\S\ref{sec:kernel}): the order-1
  arc-cosine (ReLU) kernel built on a structured spatial prior covariance
  $\Cmat(\Amp,\bwid,\smoo,\rfc)$ that encodes a local, smooth weighting centred on
  the localized region, together with the two kernel gradients the rest of the
  pipeline needs.
\item \textbf{Part~II --- inference.} The Poisson likelihood breaks conjugacy, so
  the posterior over the latent is approximated by a sparse variational GP on $\Mind$
  inducing points with Gaussian posterior $q(\latt)=\Gauss(\vm,\vV)$, optimized
  against the evidence lower bound (ELBO); the model is trained by an EM-style
  coordinate-ascent loop (E-, F-, M-steps).
\item \textbf{Part~III --- active learning.} The information utility: the marginal
  count entropy $\Hmarg$, the production utility
  $\Ustd=\Hmarg-\E[\Hnoise]$, and the distribution-aware research variant
  $\Uda=\Hmarg-\E_{x\sim\px}[\Hcond]$; a deep layer of conditioning, subspace, and
  divergence results grounded in the kernel of this part.
\item \textbf{Part~IV --- synthesis.} Guided diffusion: a denoising diffusion prior
  whose reverse process is steered by $\nabla_x\Ustd$ (which flows through
  $\nabla_x\Kf$) to synthesize high-utility inputs.
\end{itemize}

\emph{The central open problem.} A single thread runs from the prior of this part
to the analysis of Part~III. The arc-cosine kernel is positively homogeneous in
the input: its self-value $\Kf(x,x)=x\transpose\Cmat x+\sigz^2$
(\S\ref{subsec:selfkernel}) grows with the input norm. Under utility maximization
the posterior variance then grows with the square of the input scale, so the
utility \emph{diverges} and the optimizer drives inputs to the boundary of the
bounded box (saturation). Whether a kernel exists that both (i) keeps the utility
bounded on the bounded input domain and (ii) fits well is the central \emph{open
problem}; it is set up here and analyzed in Part~III (with two candidate fixes, a
normalized arc-cosine kernel and a saturating arc-sin kernel).


\section{The generative model}
\label{sec:genmodel}

\subsection{Latent Gaussian process, exponential link, Poisson counts}

The count observation $\spk_i$ (also written $r_i$, or $n$ for a novel test input)
for input $x_i$ is a Poisson draw whose rate is an exponential link applied to a
latent Gaussian-process function $\lat$.

\begin{definition}[Generative model]
\label{eq:genmodel}
With a zero-mean GP prior over input space,
\begin{equation}
\lat(\cdot)\ \sim\ \GProc\!\left(0,\ \Kf\right),
\qquad
\rate(x)\ =\ \exp\!\big(\gain\,\lat(x)+\latz\big),
\qquad
\spk_i\ \sim\ \Poi\!\big(\rate(x_i)\big).
\end{equation}
Here $\lat(x)\in\Reals$ is the \emph{latent function} (what the GP learns;
\emph{not} the rate); $\rate(x)>0$ is the \emph{rate} (the Poisson rate
parameter); $\gain>0$ is the \emph{gain}, scaling how strongly the latent
modulates the rate; and $\latz\in\Reals$ is the \emph{bias} / log-baseline
rate, so that $\rate=e^{\latz}$ when $\lat=0$. The log-rate is
$\lgr(x)=\gain\,\lat(x)+\latz$, giving $\rate=e^{\lgr}$.
\end{definition}

For numerical stability the gain is constrained positive and bounded,
$\gain\in(0,10]$, and the bias to $\latz\in[-50,50]$.

The GP prior is defined by the property that any finite set of latent values
$\{\lat(x_1),\dots,\lat(x_n)\}$ is jointly Gaussian, $\lat(X)\sim\Gauss(0,\Kf)$
with $[\Kf]_{ij}=\Kf(x_i,x_j)=\mathrm{Cov}[\lat(x_i),\lat(x_j)]$. The mean is
zero. The kernel $\Kf$ is fixed in \emph{form} before seeing data but its entries
depend on the inputs through evaluation; its full definition is
\S\ref{sec:kernel}.

\subsection{Why the GP models the latent, not the rate}

Modelling $\lat$ as Gaussian and pushing it through $\rate=\exp(\gain\lat+\latz)$
means the rate is \emph{log-normally} distributed: we learn a scaled, shifted log
of the Poisson rate rather than the rate itself. Two properties make this the
right choice for count data.

\begin{itemize}[leftmargin=1.4em]
\item \emph{Positivity for free.} The exponential link guarantees $\rate>0$ for any
  real latent value. A Gaussian prior placed \emph{directly} on the rate would
  assign mass to negative (invalid) rates.
\item \emph{Smoothness where it belongs.} The GP prior imposes spatial structure
  (locality and smoothness, \S\ref{sec:kernel}) on the unconstrained latent
  $\lat$, not on the constrained rate. The prior belief --- that $\lat$ is a
  spatially local, smooth function of the input --- is naturally expressed on the
  unconstrained log-rate, which is where the structure lives.
\end{itemize}

\begin{remark}[Poisson breaks conjugacy]
For a Gaussian likelihood the posterior and predictive would be closed-form. With
the Poisson likelihood the marginal
$p(\spk\mid X,\theta)=\int \Poi(\spk\mid\rate(x))\,\Gauss(\lat\mid0,\Kf_\theta)\,d\lat$
has no closed form, so neither the exact posterior $p(\lat\mid\spk,X,\theta)$ nor
the exact predictive can be computed directly. This is what forces the variational
treatment of Part~II; here we fix only the model and its prior.
\end{remark}


\section{The arc-cosine kernel and the structured spatial prior}
\label{sec:kernel}

The prior covariance $\Kf$ is a first-order arc-cosine (ReLU) kernel built on a
structured input-covariance matrix $\Cmat$. This section owns the definitions of
both $\Kf$ and $\Cmat$, their parameterization and constraints, and the two
gradients ($\nabla_x\Kf$, $\partial\Cmat/\partial\theta$) that the utility, the
diffusion guidance, and the M-step consume.

\subsection{The order-1 arc-cosine kernel}
\label{subsec:acos}

\begin{definition}[Order-1 arc-cosine kernel]
For two input vectors $x,x'\in\Reals^{d}$,
\begin{equation}
\Kf(x,x')\ =\ \frac{1}{\pi}\,\Kmag\,\Jang(\kang),
\qquad
\Kmag=\sqrt{\vmag_x\,\vmag_{x'}},
\qquad
\cos\kang=\frac{x\transpose\Cmat x'+\sigz^2}{\Kmag},
\qquad
\Jang(\kang)=\sin\kang+(\pi-\kang)\cos\kang,
\label{eq:acoskernel}
\end{equation}
with self-magnitudes (prior variances) $\vmag_x=x\transpose\Cmat x+\sigz^2$ and
$\vmag_{x'}={x'}\transpose\Cmat x'+\sigz^2$, cross-term
$x\transpose\Cmat x'+\sigz^2$, magnitude $\Kmag$, angle $\kang$, and order-1
angular term $\Jang$. Equivalently, expanding the angle,
\begin{equation}
\kang=\arccos\!\left(
\frac{x\transpose\Cmat x'+\sigz^2}
     {\sqrt{x\transpose\Cmat x+\sigz^2}\ \sqrt{{x'}\transpose\Cmat x'+\sigz^2}}
\right).
\end{equation}
Here $\sigz^2$ is a constant \emph{bias variance} and $\Cmat\in\Reals^{d\times d}$
is the structured spatial prior covariance (\S\ref{subsec:Cprior}).
\end{definition}

\begin{remark}[Order and ReLU-network interpretation]
This is the order-$n=1$ member of the Cho--Saul (2009) arc-cosine family
($n=0$ is Heaviside/step, $n=1$ is ReLU); the angular term
$\Jang(\kang)=\sin\kang+(\pi-\kang)\cos\kang$ is exactly the ReLU case. The kernel
is the covariance of the output of an infinite-width, single-hidden-layer ReLU
network whose input-to-hidden weights are drawn $w\sim\Gauss(0,\Cmat)$, with the
bias contributing $\sigz^2$:
\begin{equation}
\Kf(x,x')\ =\ \E_{w\sim\Gauss(0,\Cmat)}\!\big[\,\mathrm{ReLU}(w\transpose x)\,\mathrm{ReLU}(w\transpose x')\,\big]
\quad\text{(up to the $1/\pi$ constant, with bias $\sigz^2$).}
\end{equation}
Thus $\Cmat$ is literally the prior covariance of the network weights across grid
points. Shaping $\Cmat$ (\S\ref{subsec:Cprior}) injects the prior belief that
$\lat$ depends on the input through a spatially \emph{local} and \emph{smooth}
region of grid points.
\end{remark}

\subsection{Non-stationarity and the self-kernel identity}
\label{subsec:selfkernel}

\begin{proposition}[Self-kernel identity]
The prior variance at any input is the $\Cmat$-quadratic form plus the bias
variance:
\begin{equation}
\Kf(x,x)\ =\ \vmag_x\ =\ x\transpose\Cmat x+\sigz^2.
\label{eq:selfkernel}
\end{equation}
\end{proposition}

\begin{proof}
At $x'=x$ the cross-term equals $\vmag_x=\Kmag$, so $\cos\kang=1$, $\kang=0$, and
$\Jang(0)=\sin0+(\pi-0)\cos0=\pi$. Hence
$\Kf(x,x)=\Kmag\,\Jang(0)/\pi=\Kmag=\sqrt{\vmag_x\,\vmag_x}=\vmag_x$.
\end{proof}

\begin{remark}[Non-stationarity and its consequence for active learning]
Equation~\eqref{eq:selfkernel} shows $\Kf$ depends on the actual input values
through $x\transpose\Cmat x$, not only on $x-x'$: the kernel is \emph{non-stationary}
and the prior variance $\vmag_x$ grows with the input norm (in the $\Cmat$-metric).
This has a direct downstream effect: the active-learning utility is driven by
posterior variance, so utility maximization is biased toward high-norm inputs
unless the kernel is normalized. For exactly this reason there are normalized and
saturating siblings of the base kernel --- a normalized arc-cosine kernel
($\bar{\Kf}=\Jang(\kang)/\pi$, constant $\bar{\Kf}(x,x)=1$) and a saturating
arc-sin kernel (bounded diagonal) --- analyzed in Part~III as fixes to the
norm-scaling divergence. This part keeps the focus on the base arc-cosine kernel.
\end{remark}

\subsection{The structured spatial prior covariance $\Cmat$}
\label{subsec:Cprior}

$\Cmat$ encodes that the localized support of the prior is spatially local and
smooth. It is \emph{constant with respect to the input} $x$ --- it depends only on
the hyperparameters $(\Amp,\bwid,\smoo,\rfc)$. Grid point $i$ sits at grid
coordinate $\pixc_i$ on a normalized $[-1,1]\times[-1,1]$ grid, and
$\rfc=(\varepsilon_{0x},\varepsilon_{0y})$ is the center of the localized region.

\begin{definition}[Structured spatial prior]
With an amplitude $\Amp$, a per-grid-point locality weight $\loc_i$, and a pairwise
smoothing factor $\Csm$,
\begin{equation}
\Cmat_{ij}
\ =\ \Amp\,\loc_i\,[\Csm]_{ij}\,\loc_j
\ =\ \Amp\,\exp\!\left(-\frac{\lVert\pixc_i-\rfc\rVert^2}{4\bwid^2}\right)\,
      \exp\!\left(-\frac{\lVert\pixc_i-\pixc_j\rVert^2}{2\smoo^2}\right)\,
      \exp\!\left(-\frac{\lVert\pixc_j-\rfc\rVert^2}{4\bwid^2}\right),
\label{eq:Cprior}
\end{equation}
where
\begin{equation}
\begin{aligned}
\loc_i&=\exp\!\left(-\frac{\lVert\pixc_i-\rfc\rVert^2}{4\bwid^2}\right)
&&\text{(locality weight)},\\
[\Csm]_{ij}&=\exp\!\left(-\frac{\lVert\pixc_i-\pixc_j\rVert^2}{2\smoo^2}\right)
&&\text{(smoothness kernel)},
\end{aligned}
\end{equation}
followed by a symmetrization $\Cmat\leftarrow(\Cmat+\Cmat\transpose)/2$.
\end{definition}

\emph{Interpretation of the natural parameters} (defaults $\Amp=1.0$,
$\bwid=\smoo=0.1$ on the $[-1,1]$ grid):
\begin{itemize}[leftmargin=1.4em]
\item $\Amp$ --- overall \emph{amplitude}: a positive gain (default $1.0$) on the
  structured covariance $\Cmat$, scaling the prior variance of the latent. It
  multiplies the $\Cmat$-quadratic part of the self-magnitude
  $\vmag_x=x\transpose\Cmat x+\sigz^2$ (not the bias $\sigz^2$); because it enters
  the kernel inside $\Kmag=\sqrt{\vmag_x\vmag_{x'}}$ and $\cos\kang$, it rescales
  $\Kf$ non-linearly rather than as a plain output scale. Unlike $\bwid,\smoo,\rfc$
  it sets the magnitude, not the spatial shape, of the localized weighting.
\item $\bwid$ --- \emph{half-width} of the localized region: the locality weight
  $\loc$ is a Gaussian of standard deviation $\sqrt{2}\,\bwid$ in grid-distance,
  reaching $e^{-1}$ at $\text{dist}=2\bwid$; for $\bwid=0.1$ the localized region
  spans roughly $\pm0.2$ ($\approx20\%$ of the axis).
\item $\smoo$ --- smoothness \emph{standard deviation} (correlation length) of the
  RBF-like $\Csm$ (variance $\smoo^2$): grid points closer than $\smoo$ are
  strongly correlated, grid points beyond $3\smoo$ nearly uncorrelated.
\item $\rfc$ --- the 2-D center of the localized region on $[-1,1]^2$.
\end{itemize}
Together, $\loc$ localizes weight around $\rfc$ and $\Csm$ correlates nearby grid
points --- a localized, smooth weight prior.

\subsection{Raw parameterization and constraints}
\label{subsec:rawparam}

For optimizer stability the model is parameterized by unconstrained log-space
\emph{raw} parameters rather than $\bwid$ and $\smoo$ directly:
\begin{equation}
\begin{aligned}
\text{raw\_m2log2beta}&=-2\log(2\bwid)&&\Rightarrow&&\exp(\text{raw\_m2log2beta})=\frac{1}{4\bwid^2},\\
\text{raw\_mlog2rho2}&=-\log(2\smoo^2)&&\Rightarrow&&\exp(\text{raw\_mlog2rho2})=\frac{1}{2\smoo^2},
\end{aligned}
\end{equation}
and $\sigz=\exp(\text{raw\_sigma\_0})$ under a positivity constraint (so the bias
variance added in~\eqref{eq:acoskernel} is $\sigz^2=(\exp(\text{raw\_sigma\_0}))^2$).
The exponentials of the raw parameters are used \emph{directly} as the exponent
coefficients, so that
$\loc_i=\exp(-\tfrac{1}{4\bwid^2}\lVert\pixc_i-\rfc\rVert^2)$ and
$[\Csm]_{ij}=\exp(-\tfrac{1}{2\smoo^2}\lVert\pixc_i-\pixc_j\rVert^2)$ reproduce
exactly the natural-parameter form~\eqref{eq:Cprior} (numerically,
$\bwid=\smoo=0.1\Rightarrow 1/(4\bwid^2)=25$, $1/(2\smoo^2)=50$). The raw
parameters invert back to the natural values,
$\bwid=\tfrac12\exp(-\text{raw\_m2log2beta}/2)$,
$\smoo=\tfrac{1}{\sqrt2}\exp(-\text{raw\_mlog2rho2}/2)$, and round-trip cleanly.

\begin{table}[h]
\centering
\begin{tabular}{@{}llll@{}}
\toprule
Param & Constraint mechanism & Bounds (natural) & Bounds (raw) \\
\midrule
$\sigz$ & $\sigz=\exp(\text{raw\_sigma\_0})$, positive & $\sigz>0$ & --- \\
$\bwid$ & via $\text{raw\_m2log2beta}$, clamped & $[0.01,\,0.3]$ & $[1.02,\,7.82]$ \\
$\smoo$ & via $\text{raw\_mlog2rho2}$, clamped & $[0.01,\,0.5]$ & $[0.69,\,8.52]$ \\
$\rfc$ & direct clamp & $[-1,1]$ per axis & --- \\
$\Amp$ & $\Amp=\mathrm{softplus}(\text{raw\_Amp})$, clamp & $(0,\,1000]$ & --- \\
\bottomrule
\end{tabular}
\end{table}

\begin{sloppypar}
Bounds are enforced two ways: by read-only rejection of an out-of-bounds trial
step during optimization, and by a projected clamp applied after each optimizer
step.
\end{sloppypar}

\begin{remark}[Restricting $\Cmat$ to its effective support]
The locality weight $\loc_i=\exp(-\lVert\pixc_i-\rfc\rVert^2/(4\bwid^2))$ decays as
a Gaussian away from $\rfc$, so grid points whose weight $\loc_i$ falls below a
small threshold (e.g. $10^{-3}$) contribute negligibly to $x\transpose\Cmat x$ and
may be dropped, shrinking $\Cmat$ from $(d,d)$ to $(d',d')$ with $d'\ll d$ (the
same restriction is applied to the input so only the retained grid points enter
$x\transpose\Cmat x$). With threshold $10^{-3}$ the retained support is the disk
$\lVert\pixc_i-\rfc\rVert\le2\bwid\sqrt{\ln1000}\approx5.25\,\bwid$ around $\rfc$;
the retained entries of $\Cmat$ still carry gradients, while the support itself is
computed from detached parameters and is therefore held fixed
(non-differentiable) and stable.
\end{remark}

\subsection{Kernel gradient with respect to the input, $\nabla_x\Kf$}
\label{subsec:gradxK}

Choosing the next input (and steering diffusion, Part~IV) requires the derivative
of the kernel with respect to the input. Treating $\Cmat$ as symmetric and
constant, $\sigz^2$ constant, and $x'$ (hence $\vmag_{x'}$) constant, the useful
intermediate gradients are
\begin{equation}
\nabla_x\,\vmag_x=2\,\Cmat x,
\qquad
\nabla_x\,(x\transpose\Cmat x'+\sigz^2)=\Cmat x',
\qquad
\nabla_x\,\Kmag=\sqrt{\tfrac{\vmag_{x'}}{\vmag_x}}\,\Cmat x,
\end{equation}
and the angular derivative $d\Jang/d\kang=-(\pi-\kang)\sin\kang$.

\begin{proposition}[Input gradient of the arc-cosine kernel]
\begin{equation}
\nabla_x\Kf(x,x')
\ =\ \frac{1}{\pi}\left[(\pi-\kang)\,\Cmat x'
\ +\ \sin\kang\,\sqrt{\frac{\vmag_{x'}}{\vmag_x}}\,\Cmat x\right]
\ =\ \frac{1}{\pi}\left[(\pi-\kang)\,\Cmat x'
\ +\ \sin\kang\,\sqrt{\frac{{x'}\transpose\Cmat x'+\sigz^2}{x\transpose\Cmat x+\sigz^2}}\,\Cmat x\right].
\label{eq:gradxK}
\end{equation}
By symmetry ($x\leftrightarrow x'$),
\begin{equation}
\nabla_{x'}\Kf(x,x')
=\frac{1}{\pi}\left[(\pi-\kang)\,\Cmat x
+\sin\kang\,\sqrt{\frac{x\transpose\Cmat x+\sigz^2}{{x'}\transpose\Cmat x'+\sigz^2}}\,\Cmat x'\right].
\end{equation}
\end{proposition}

\begin{proof}
Write $\Kf=\tfrac1\pi\Kmag\,\Jang(\kang)$ and apply the product rule
$\nabla_x\Kf\propto(\nabla_x\Kmag)\Jang+\Kmag(\nabla_x\Jang)$. Using
$\nabla_x\Jang=(d\Jang/d\kang)\nabla_x\kang$ with
$\nabla_x\kang=-(1/\sin\kang)\,\nabla_x(\cos\kang)$ and
$d\Jang/d\kang=-(\pi-\kang)\sin\kang$, the $\sin\kang$ cancels and
$\nabla_x\Jang=(\pi-\kang)\,\nabla_x(\cos\kang)$ with
$\nabla_x(\cos\kang)=\tfrac1\Kmag\Cmat x'-\tfrac{\cos\kang}{\Kmag}\sqrt{\vmag_{x'}/\vmag_x}\,\Cmat x$
(recall $\cos\kang=(x\transpose\Cmat x'+\sigz^2)/\Kmag$). Collecting terms,
\[
\pi\,\nabla_x\Kf
=\underbrace{\sin\kang\sqrt{\tfrac{\vmag_{x'}}{\vmag_x}}\Cmat x}_{[A]}
+\underbrace{(\pi-\kang)\cos\kang\sqrt{\tfrac{\vmag_{x'}}{\vmag_x}}\Cmat x}_{[B]}
+\underbrace{(\pi-\kang)\Cmat x'}_{[C]}
-\underbrace{(\pi-\kang)\cos\kang\sqrt{\tfrac{\vmag_{x'}}{\vmag_x}}\Cmat x}_{[D]} .
\]
Terms $[B]$ and $[D]$ are identical with opposite sign and cancel, leaving
$[A]+[C]$, which is~\eqref{eq:gradxK}.
\end{proof}

Every term carries a factor $\Cmat$, so $\nabla_x\Kf$ lives in the column space of
$\Cmat$; this is the algebraic root of the subspace analysis in Part~III (gradient
ascent on the utility already concentrates on the top $\Cmat$-eigenvectors).

\subsection{Kernel gradients with respect to the hyperparameters, $\partial\Cmat/\partial\theta$}
\label{subsec:dCtheta}

The M-step (Part~II training) maximizes the ELBO over the kernel hyperparameters and
needs $\partial\Cmat/\partial\theta$, which then chains through
$\vmag_x=x\transpose\Cmat x$ and the cross-term into $\partial\Kf/\partial\theta$
using the same quadratic-form structure as~\S\ref{subsec:gradxK} with $\Cmat$
replaced by $\partial\Cmat/\partial\theta$.

\begin{proposition}[Structured-prior hyperparameter gradients (natural parameters)]
Writing $\Cmat=\Amp\,\loc\,\Csm\,\loc\transpose$, the amplitude enters linearly and
the shape parameters through their Gaussian exponents. Each shape-parameter
derivative is $\Cmat_{ij}$ --- which already carries the $\Amp$ factor --- times the
derivative of its log-argument, while the amplitude derivative strips one factor of
$\Amp$:
\begin{equation}
\label{eq:dCtheta}
\begin{aligned}
\frac{\partial\Cmat_{ij}}{\partial\Amp}
&=\loc_i\,[\Csm]_{ij}\,\loc_j=\frac{\Cmat_{ij}}{\Amp},\\[2pt]
\frac{\partial\Cmat_{ij}}{\partial\bwid}
&=\Cmat_{ij}\,\frac{\lVert\pixc_i-\rfc\rVert^2+\lVert\pixc_j-\rfc\rVert^2}{2\bwid^3},\\[2pt]
\frac{\partial\Cmat_{ij}}{\partial\smoo}
&=\Cmat_{ij}\,\frac{\lVert\pixc_i-\pixc_j\rVert^2}{\smoo^3},\\[2pt]
\nabla_{\rfc}\Cmat_{ij}
&=\frac{\Cmat_{ij}}{2\bwid^2}\big[(\pixc_i-\rfc)+(\pixc_j-\rfc)\big]\qquad(\text{2D vector}).
\end{aligned}
\end{equation}
\end{proposition}

\begin{proof}
For $\Amp$: $\Cmat=\Amp\,(\loc\,\Csm\,\loc\transpose)$ is linear in $\Amp$, so
$\partial\Cmat_{ij}/\partial\Amp=\loc_i[\Csm]_{ij}\loc_j=\Cmat_{ij}/\Amp$. For $\bwid$:
$\partial_\bwid\big[-(\lVert\pixc_i-\rfc\rVert^2+\lVert\pixc_j-\rfc\rVert^2)/(4\bwid^2)\big]
=+(\lVert\pixc_i-\rfc\rVert^2+\lVert\pixc_j-\rfc\rVert^2)/(2\bwid^3)$. For $\smoo$:
$\partial_\smoo\big[-\lVert\pixc_i-\pixc_j\rVert^2/(2\smoo^2)\big]=+\lVert\pixc_i-\pixc_j\rVert^2/\smoo^3$.
For $\rfc$: $\nabla_{\rfc}\lVert\pixc-\rfc\rVert^2=2(\rfc-\pixc)$, giving
$\nabla_{\rfc}\log\Cmat_{ij}=(1/(2\bwid^2))[(\pixc_i-\rfc)+(\pixc_j-\rfc)]$.
\end{proof}

\begin{remark}[Gradients in the raw parameterization]
Because the model optimizes the raw parameters (\S\ref{subsec:rawparam}), the
corresponding \emph{raw}-parameter gradients follow from~\eqref{eq:dCtheta} by the
chain rule:
\begin{equation}
\begin{aligned}
\frac{\partial\Cmat_{ij}}{\partial(\text{raw\_m2log2beta})}
&=\Cmat_{ij}\,(\log\loc_i+\log\loc_j)
=-\Cmat_{ij}\,\frac{\lVert\pixc_i-\rfc\rVert^2+\lVert\pixc_j-\rfc\rVert^2}{4\bwid^2},\\
\frac{\partial\Cmat_{ij}}{\partial(\text{raw\_mlog2rho2})}
&=\Cmat_{ij}\,\log[\Csm]_{ij}
=-\Cmat_{ij}\,\frac{\lVert\pixc_i-\pixc_j\rVert^2}{2\smoo^2}.
\end{aligned}
\end{equation}
and, since $\rfc$ is a direct parameter,
$\partial\Cmat_{ij}/\partial\rfc=(\Cmat_{ij}/(2\bwid^2))[(\pixc_i-\rfc)+(\pixc_j-\rfc)]$
coincides with the natural form in~\eqref{eq:dCtheta}. The raw and natural
gradients are related by the constraint-transform Jacobian
$d(\text{raw})/d\bwid=-2/\bwid$, etc.
\end{remark}

\subsection{Numerical stability}
\label{subsec:numstab}

Three stability mechanisms act at three distinct levels; they must not be conflated.

\begin{enumerate}[leftmargin=1.6em]
\item \textbf{$\Cmat$ matrix --- symmetrization only.}
  $\Cmat\leftarrow(\Cmat+\Cmat\transpose)/2$ cancels the floating-point rounding
  asymmetry from the outer product. \emph{No diagonal jitter is added to $\Cmat$.}
\item \textbf{Kernel evaluation --- domain clamps.} With
  $\text{eps}=10^{-7}$, $\cos\kang$ is clamped into $(-1,1)$ so $\arccos$ is finite,
  and $1-\cos^2\kang$ is clamped at $\min=\text{eps}$ so $\sin\kang=\sqrt{1-\cos^2\kang}$
  never takes the square root of a negative rounding result. These keep the
  transcendental operations in-domain and are unrelated to Cholesky jitter.
\item \textbf{GP covariance / Cholesky --- jitter on the gram matrices, not on
  $\Cmat$.} Jitter is added to the \emph{output} gram matrices of the arc-cosine
  kernel --- the inducing gram $\Ktil=\Kf(\Ipt,\Ipt)$ and the predictive gram
  $\Kf(X,X)$ --- \emph{not} to $\Cmat$. A pre-Cholesky jitter (default $10^{-4}$) is
  added to $\Ktil$ and $\Kf(X,X)$; the kernel evaluation itself adds none (that
  would double-jitter). On a failed factorization the gram matrix is promoted to
  double precision (sequential Cholesky subtractions are prone to catastrophic loss
  of accuracy in single precision) and the factorization is retried with escalating
  jitter $10^{-4},10^{-3},10^{-2}$ (up to three tries).
\end{enumerate}

\caveat{Some intuition expects diagonal jitter on the prior $\Cmat$; there is none.
$\Cmat$ receives only symmetrization. Jitter lives exclusively on the GP gram
matrices $\Ktil$/$\Kf(X,X)$. Do not ``add jitter to $\Cmat$'' thinking it is
missing.}


\part{Inference: variational learning of the latent}
% ============================================================================
%  part2a_inference.tex  --  Part II, sections 4-5.
%  Sparse variational Gaussian-process inference + the eigenspace representation.
%  Self-contained \input-ready fragment; uses ONLY the shared preamble macros.
% ============================================================================

\section{Sparse variational Gaussian-process inference}\label{sec:varinf}

Part~I fixed the generative model: a zero-mean latent Gaussian process
$\lat \sim \GProc(0,\Kf)$ over input space, an exponential link to the rate
$\rate(x)=\exp\!\big(\gain\,\lat(x)+\latz\big)$, and Poisson counts
$\spk_i \sim \Poi\!\big(\rate(x_i)\big)$, where $\Kf$ is the first-order
arc-cosine kernel of \eqref{eq:acoskernel} built on the structured spatial
prior covariance $\Cmat$ of \eqref{eq:Cprior}. Here we turn that model into a
trainable posterior. The variational posterior developed below can be
represented equivalently in the full inducing space or in the eigenbasis of
\S\ref{sec:eigenspace}.

\subsection{Poisson breaks conjugacy}\label{sec:conjugacy}

Given a batch of inputs $X=(x_1,\dots,x_{\Ntr})$ and counts
$\spk=(\spk_1,\dots,\spk_{\Ntr})$, the exact posterior over the latent and the
marginal likelihood,
\begin{align}
  p(\lat\mid\spk,X,\theta)&\;\propto\;
    \Big[\textstyle\prod_i \Poi\!\big(\spk_i\mid e^{\gain\lat_i+\latz}\big)\Big]\,
    \Gauss(\lat\mid 0,\Kf_\theta), \notag\\
  p(\spk\mid X,\theta)&=\!\int\! \prod_i \Poi\!\big(\spk_i\mid e^{\gain\lat_i+\latz}\big)\,
    \Gauss(\lat\mid 0,\Kf_\theta)\,d\lat ,
\end{align}
have no closed form: the exponential-Poisson likelihood is not conjugate to the
Gaussian prior, so the integral over $\lat$ is intractable. We therefore
approximate the posterior by a tractable Gaussian and maximise a variational
lower bound (\S\ref{sec:elbo}).

\begin{remark}[The conjugate contrast]
For a \emph{Gaussian} observation $\spk=\lat+\varepsilon$,
$\varepsilon\sim\Gauss(0,\pvar_r I)$, everything would be closed-form: the
marginal is $\spk\mid X\sim\Gauss(0,\Kf+\pvar_r I)$, the predictive at a test
input $x_\ast$ is Gaussian with mean $\kvec_\ast\transpose(\Kf+\pvar_r I)^{-1}\spk$
and variance $\Kf(x_\ast,x_\ast)-\kvec_\ast\transpose(\Kf+\pvar_r I)^{-1}\kvec_\ast$,
and the log-evidence is
$-\tfrac12\spk\transpose(\Kf+\pvar_r I)^{-1}\spk-\tfrac12\log\lvert\Kf+\pvar_r I\rvert-\tfrac{\Ntr}{2}\log 2\pi$.
None of these survive the Poisson link; the variational route below recovers a
usable posterior in its place.
\end{remark}

\subsection{Inducing points and the variational posterior}\label{sec:inducing}

\begin{definition}[Sparse variational posterior]
Introduce $\Mind$ \emph{inducing points} $\Ipt=\{\ipt_1,\dots,\ipt_{\Mind}\}$
with latent values $\latt=\lat(\Ipt)$, and place a
Gaussian variational posterior on them,
\begin{equation}
  q(\latt)=\Gauss(\latt\mid\vm,\vV),\qquad
  \vm\in\Reals^{\Mind},\quad \vV\in\mathbb{S}_{++}^{\Mind}.
  \label{eq:qdef}
\end{equation}
The posterior at \emph{any} input $x$ is recovered by marginalising the
conditional prior $p(\lat(x)\mid\latt)$ against $q(\latt)$.
\end{definition}

Exact inference would need the $\Ntr\!\times\!\Ntr$ kernel over all training
inputs, an $O(\Ntr^3)$ cost; the inducing construction reduces it to
$O(\Ntr\Mind^2)$.

\begin{remark}[Active-learning invariant $\Mind=\Ntr$]\label{rem:MeqN}
In the active-learning setting the inducing set is taken equal to the training
set, $X=\Ipt$, so $\Mind=\Ntr$ and $\Kf=\Ktil$. The ``sparse'' posterior is then
\emph{exact}. The SVGP formalism is retained not for sparsity but for two other
reasons: (i)~it supplies the closed-form variational machinery of
\S\ref{sec:predmoments}--\S\ref{sec:fbar}, and (ii)~its inducing structure admits
an $O(\Mind)$ rank-one online growth that appends one input per active-learning
step.
\end{remark}

\subsection{Projected posterior moments}\label{sec:predmoments}

Marginalising $p(\lat(x)\mid\latt)\,q(\latt)$ yields a Gaussian
$q(\lat(x))=\Gauss\!\big(\pmean(x),\pvar(x)\big)$ whose moments are the
inducing-point moments projected through the kernel:
\begin{align}
  \pmean(x)  &= \kvec(x)\transpose\,\Ktil^{-1}\,\vm , \notag\\
  \pvar(x)   &= \Kf(x,x) + \kvec(x)\transpose\,\Ktil^{-1}\big(\vV-\Ktil\big)\Ktil^{-1}\,\kvec(x),
  \label{eq:predmoments}
\end{align}
with the cross-covariance vector $\kvec(x)=\Kf(\Ipt,x)\in\Reals^{\Mind}$, the
inducing gram $\Ktil=\Kf(\Ipt,\Ipt)\in\Reals^{\Mind\times\Mind}$, and the prior
self-variance $\Kf(x,x)$ given by the self-kernel identity \eqref{eq:selfkernel}.
Introducing the \emph{projection vector} $\pv(x)=\Ktil^{-1}\kvec(x)$ and the
Schur complement $s(x)=\Kf(x,x)-\kvec(x)\transpose\Ktil^{-1}\kvec(x)$, the
variance splits into a data-independent residual and an epistemic term,
\begin{equation}
  \pvar(x)=\underbrace{\big[\Kf(x,x)-\kvec(x)\transpose\Ktil^{-1}\kvec(x)\big]}_{\text{residual (Nystr\"om/Schur) }s(x)}
    \;+\;\underbrace{\kvec(x)\transpose\Ktil^{-1}\vV\Ktil^{-1}\kvec(x)}_{\text{epistemic}}
    \;=\; s(x)+\pv(x)\transpose\vV\,\pv(x),
  \label{eq:varsplit}
\end{equation}
and the mean reads $\pmean(x)=\pv(x)\transpose\vm$. The two forms of $\pvar$ are
algebraically identical since
$\kvec\transpose\Ktil^{-1}(\vV-\Ktil)\Ktil^{-1}\kvec=\kvec\transpose\Ktil^{-1}\vV\Ktil^{-1}\kvec-\kvec\transpose\Ktil^{-1}\kvec$.

\begin{remark}[Prior sanity check]
At the prior $\vV=\Ktil$ the correction in \eqref{eq:predmoments} vanishes and
$\pvar(x)=\Kf(x,x)$ (the prior variance); as data drive $\vV$ below $\Ktil$,
$\pvar(x)<\Kf(x,x)$.
\end{remark}

These moments can be evaluated directly in the $\Mind$-dimensional inducing
space or, equivalently, in the eigenbasis of $\Ktil$ (\S\ref{sec:eig-moments}).

\subsection{The ELBO}\label{sec:elbo}

All parameters are learned by maximising the evidence lower bound on
$\log p(\spk)$,
\begin{equation}
  \ELBO \;=\; \sum_i \E_{q(\lat_i)}\!\big[\log \Poi(\spk_i\mid\lat_i)\big]
    \;-\; \KLdiv\!\big(q(\latt)\,\big\|\,p(\latt)\big),
  \label{eq:elbo}
\end{equation}
which is tight when $q$ equals the true posterior. Because the likelihood
factorises over inputs and $q(\lat_i)=\Gauss(\pmean_i,\pvar_i)$ is Gaussian, each
expected log-likelihood term is closed-form:
\begin{equation}
  \E_{q(\lat_i)}\!\big[\log \Poi(\spk_i\mid\lat_i)\big]
    \;=\; \spk_i\big(\gain\pmean_i+\latz\big)
      \;-\; \exp\!\big(\gain\pmean_i+\tfrac12\gain^2\pvar_i+\latz\big)
      \;+\; \text{const},
  \label{eq:ell}
\end{equation}
with $\text{const}=-\log(\spk_i!)$ independent of all parameters. The KL term is
the standard Gaussian--Gaussian divergence between $q(\latt)=\Gauss(\vm,\vV)$ and
the prior $p(\latt)=\Gauss(0,\Ktil)$,
\begin{equation}
  \KLdiv(q\,\|\,p) \;=\; \tfrac12\Big[\,
    \log\frac{\lvert\Ktil\rvert}{\lvert\vV\rvert}
    \;+\; \Tr\!\big(\Ktil^{-1}\vV\big)
    \;+\; \vm\transpose\Ktil^{-1}\vm
    \;-\; \Mind \,\Big].
  \label{eq:kl}
\end{equation}

\begin{remark}[The $-\Mind$ constant]
The $-\Mind$ in \eqref{eq:kl} is parameter-independent (it is the dimension of
the inducing Gaussian) and carries zero gradient, so it affects only the reported
loss magnitude, never the optimum or the predictions. The eigenspace evaluation
of \S\ref{sec:eig-kl} carries the analogous $-\nb$ for the same reason.
\end{remark}

The full objective is maximised by EM-style coordinate ascent
(E-step on $(\vm,\vV)$, F-step on $(\gain,\latz)$, M-step on the kernel
hyperparameters); the update equations are derived in the training-algorithm
section. Equation~\eqref{eq:elbo} is the single objective all three blocks share.

\subsection{Expected rate via the Gaussian MGF}\label{sec:fbar}

The exponential term in \eqref{eq:ell} is the \emph{expected rate} under $q$,
\begin{equation}
  \fbar_i \;=\; \E_{q(\lat_i)}\!\big[\exp(\gain\lat_i+\latz)\big]
    \;=\; \exp\!\big(\gain\pmean_i+\tfrac12\gain^2\pvar_i+\latz\big).
  \label{eq:fbar}
\end{equation}
This is the Gaussian moment-generating function
$\E[e^{tZ}]=\exp(t\mu+\tfrac12 t^2\sigma^2)$ applied to $Z=\gain\lat_i+\latz$ at
$t=1$.

\begin{remark}[Jensen inflation]
The term $\tfrac12\gain^2\pvar_i$ is strictly positive whenever $\pvar_i>0$, so
posterior \emph{uncertainty raises} the predicted rate above the point estimate
$\exp(\gain\pmean_i+\latz)$ --- a direct consequence of Jensen's inequality for
the convex map $\exp$. Confidence and rate are thus coupled: a poorly-constrained
input is predicted to have a \emph{higher} rate, not lower.
\end{remark}

\subsection{Prediction: expected count}\label{sec:prediction}

The moments \eqref{eq:predmoments} do not depend on the observed $\spk$, so the
same formulas serve as the predictive distribution at any novel input $x_\ast$.
The expected count is the expected rate \eqref{eq:fbar} evaluated there,
\begin{equation}
  \langle \spk_\ast\rangle=\langle\rate(x_\ast)\rangle
    =\exp\!\big(\gain\pmean(x_\ast)+\tfrac12\gain^2\pvar(x_\ast)+\latz\big),
\end{equation}
with $\pmean(x_\ast),\pvar(x_\ast)$ from \eqref{eq:predmoments} using freshly
evaluated $\kvec(x_\ast)$ and $\Kf(x_\ast,x_\ast)$. The rate is
log-normal, so its variance is
\begin{equation}
  \operatorname{Var}\!\big(\rate(x)\big)
    =\langle\rate(x)\rangle^2\big(e^{\gain^2\pvar(x)}-1\big).
  \label{eq:fvar}
\end{equation}

\subsection{The posterior cross-covariance pitfall}\label{sec:crosscov}

The self-variance of \eqref{eq:predmoments} extends formally to the
\emph{cross-covariance} between two distinct inputs $x$ (a candidate sample
location) and $x_\ast$ (a query location):
\begin{equation}
  \Sigma_{x,x_\ast}=\operatorname{Cov}\!\big[\lat(x),\lat(x_\ast)\mid \mathcal D\big]
    =\Kf(x,x_\ast)+\pv(x)\transpose\big(\vV-\Ktil\big)\pv(x_\ast),
  \qquad \pv(\cdot)=\Ktil^{-1}\kvec(\cdot),
  \label{eq:crosscov}
\end{equation}
which reduces to $\pvar(x)$ at $x=x_\ast$.

\begin{remark}[Failure outside the inducing hull]
For a query $x_\ast$ \emph{outside} the inducing region, \eqref{eq:crosscov}
produces cross-covariances that violate the Cauchy--Schwarz bound
$\lvert\Sigma_{x,x_\ast}\rvert\le\sqrt{\Sigma_{x,x}\,\Sigma_{x_\ast,x_\ast}}$.
Measured example (RBF, $\ell=0.2$, $20$ inducing points in $[-1,1]$, sample
$x=0$): at $x_\ast=\pm1.5$ the manual formula gives $\Sigma\approx-4$ while
the consistent joint covariance gives $\approx 0$; with self-variances $\sim0.04$
the valid maximum is $\sim0.04$, so $-4$ is two orders of magnitude out of range.
Root cause: in the residual $\Kf(x,x_\ast)-\pv(x)\transpose\Ktil\,\pv(x_\ast)$ the
direct kernel $\Kf(x,x_\ast)\!\to\!0$ (finite lengthscale) outside the hull, but
$\pv(x_\ast)=\Ktil^{-1}\kvec(x_\ast)$ can have large entries through $\Ktil^{-1}$,
and the two pieces fail to cancel.
\end{remark}

Left unchecked, a value like $\Sigma_{x,x_\ast}\approx-4$ feeds the conditional
moments
$\pmean_{\mathrm{cond}}(x_\ast)=\pmean(x_\ast)+(\Sigma_{x,x_\ast}/\Sigma_{x,x})(\lat_{\mathrm{obs}}-\pmean(x))$
and $\pvar_{\mathrm{cond}}(x_\ast)=\Sigma_{x_\ast,x_\ast}-\Sigma_{x,x_\ast}^2/\Sigma_{x,x}$
of the distribution-aware utility (Part~III) with a regression coefficient
$\approx-100$, giving nonsensical conditional means and spurious utility peaks at
the edges of the inducing region. The remedy is to never assemble the cross-block
from the manual formula \eqref{eq:crosscov}; instead read the cross-block from the
jointly-formed covariance of $(x,x_\ast)$, which is Nystr\"om-consistent --- it
carries the conditional correction term that \eqref{eq:crosscov} drops outside the
hull. The pathology is therefore confined to hand-assembled cross-covariances; any
computation that forms a cross-covariance from \eqref{eq:crosscov} outside the
inducing hull will reintroduce it.

\subsection{Whitening}\label{sec:whitening}

An alternative reparameterisation is \emph{Cholesky whitening}, which makes the
prior a standard normal. With the Cholesky factor $\Ktil=\Lchol\Lchol\transpose$,
define the whitened inducing values
\begin{equation}
  \widetilde{\latt}=\Lchol^{-1}(\latt-\mu_z)\;\sim\;\Gauss(0,I)\quad\text{under the prior},
\end{equation}
where $\mu_z$ is the prior mean at the inducing points. Optimisation conditions
much better in these coordinates. In these coordinates one stores the
\emph{whitened} variational parameters $(\widetilde{\vm},\widetilde{\vV})$,
related to the actual ones by
\begin{equation}
  \vm-\mu_z=\Lchol\,\widetilde{\vm},\quad \vV=\Lchol\,\widetilde{\vV}\,\Lchol\transpose;
  \qquad
  \widetilde{\vm}=\Lchol^{-1}(\vm-\mu_z),\quad \widetilde{\vV}=\Lchol^{-1}\vV\,\Lchol^{-\mathsf{T}}.
  \label{eq:whiten}
\end{equation}
The whitened mean thus encodes the \emph{deviation} from the prior mean function,
not absolute function values. Recovering the moments from the whitened parameters
therefore requires three corrections: unwhiten the covariance
($\vV=\Lchol\widetilde{\vV}\Lchol\transpose$), unwhiten the mean
($\vm-\mu_z=\Lchol\widetilde{\vm}$), and add the mean function back,
$\pmean_\ast=\pmean(x_\ast)+\pv(x_\ast)\transpose\Lchol\widetilde{\vm}$ and
$\pvar_\ast=s(x_\ast)+\pv(x_\ast)\transpose\Lchol\widetilde{\vV}\Lchol\transpose\pv(x_\ast)$.

\begin{remark}[Zero prior mean]
With a zero prior mean, $\mu_z=0$ and the ``add the mean function'' step is
trivially zero; the whitening maps \eqref{eq:whiten} themselves apply for any
prior mean.
\end{remark}


\section{The eigenspace representation}\label{sec:eigenspace}

The variational posterior $q(\latt)$ of \S\ref{sec:varinf} can be represented
equivalently in the full $\Mind$-dimensional inducing space or, more
economically, in the eigenbasis of the inducing gram $\Ktil$. This section
develops the eigenbasis form: the same prior $\Gauss(0,\Ktil)$, the same
arc-cosine kernel, the same $\Poi(\gain,\latz)$ likelihood, the same expected
rate \eqref{eq:fbar}, and the same ELBO \eqref{eq:elbo}, re-expressed in a
low-dimensional basis in which $\Ktil$ is diagonal.

\subsection{Eigendecomposition of the inducing gram}\label{sec:eigdecomp}

The inducing gram $\Ktil$ ($\Mind\times\Mind$) has effective rank far below
$\Mind$ (typically $\sim10$--$50$). We eigendecompose it and keep only the
well-conditioned directions:
\begin{equation}
\begin{gathered}
  \Ktil=\eb\,\evals\,\eb\transpose,\qquad
  \evals=\diago(\text{eigvals}_b), \\
  \eb\in\Reals^{\Mind\times\nb},\qquad \eb\transpose\eb=I_{\nb},
\end{gathered}
  \label{eq:eigdecomp}
\end{equation}
where the kept set is chosen by a \emph{relative} threshold with an absolute
floor,
\begin{equation}
  \text{keep eigenpair }j \iff \evals_j > \max\!\big(\evals_{\max}\cdot\texttt{tol},\ \texttt{tol}\big),
  \qquad \texttt{tol}=10^{-4}.
  \label{eq:eigthresh}
\end{equation}
The eigenvalues are taken in ascending order; $\eb$ collects the kept eigenvectors
as columns and $\evals$ the kept eigenvalues. The kept count $\nb$ is
\emph{dynamic}: it shifts across training iterations as the kernel hyperparameters
move, and grows by $\sim1$ when a point is appended. Typical $\nb\approx10$--$11$
in the active-learning setting, up to $\sim50$ in larger problems.

Three payoffs motivate the projection: (i)~dimensionality drops $\Mind\to\nb$, so
per-step cost falls $O(\Mind^3)\to O(\nb^3)$; (ii)~$\Ktil$ becomes
\emph{diagonal} in the new basis, $\Ktilb=\evals$, so $\Ktil^{-1}$ is the
element-wise reciprocal $\diago(1/\evals)$ --- no matrix solve; (iii)~dropping the
small eigenvalues gives implicit regularisation.

\subsection{The eigenspace state}\label{sec:eigstate}

The eigenspace posterior is carried together with the precomputed kernel
projections used below:

\begin{center}\small
\begin{tabular}{@{}l l >{\raggedright\arraybackslash}p{8.6cm}@{}}
\toprule
symbol & shape & meaning \\
\midrule
$\mb$ & $(\nb)$ & variational mean in the eigenbasis, $\mb=\eb\transpose\vm$ \\
$\Vb$ & $(\nb,\nb)$ & variational covariance in the eigenbasis, $\Vb=\eb\transpose\vV\eb$ --- \textbf{dense} \\
$\eb$ & $(\Mind,\nb)$ & eigenvector matrix of $\Ktil$ (the basis) \\
$\evals$ & $(\nb)$ & kept eigenvalues of $\Ktil$ (ascending) \\
$\Ktilb$ & $(\nb,\nb)$ & inducing gram in the eigenbasis $=\diago(\evals)$ --- \textbf{diagonal} \\
$\Kf_b$ & $(\Ntr,\nb)$ & cross-kernel projected into the eigenbasis, $\Kf_b=\Kf(X,\Ipt)\,\eb$ \\
$\amat$ & $(\Ntr,\nb)$ & projection $\Kf\Ktil^{-1}$ in the eigenbasis, $\amat=\Kf_b/\evals$ (element-wise) \\
$\Kvec$ & $(\Ntr)$ & diagonal self-kernel $\Kf(x_i,x_i)$, the per-point prior variance \\
\bottomrule
\end{tabular}
\end{center}

\begin{remark}[Only $\Ktilb$ is diagonal]
$\Vb$ is a \emph{dense} $\nb\times\nb$ matrix even though $\Ktilb$ is diagonal;
never assume a diagonal $\Vb$. The projection row $\amat_i=(\Kf_b/\evals)_i$
is exactly $\kvec(x_i)\transpose\Ktil^{-1}$ expressed in the eigenbasis, obtained
by element-wise division rather than a matrix solve.
\end{remark}

Initialisation sets the variational parameters to the prior:
\begin{equation}
  \mb=0,\qquad \Vb=\Ktilb \quad(\text{i.e.\ }\vm=0,\ \vV=\Ktil,\ \text{so } \KLdiv=0\text{ at start}).
\end{equation}

\subsection{Posterior moments in the eigenbasis}\label{sec:eig-moments}

With $\amat=\Kf_b/\evals$ (element-wise), the projected moments
\eqref{eq:predmoments} at the training points become:
\begin{align}
  \pmean &= \amat\,\mb, \notag\\
  \pvar  &= \Kvec + \diago\!\big(\amat\,(\Vb-\Ktilb)\,\amat\transpose\big)
          = \Kvec + \operatorname{rowsum}\!\big(\amat\odot(\amat\,(\Vb-\Ktilb))\big), \notag\\
  \pvar  &\leftarrow \max\!\big(\pvar,\ 10^{-6}\big),
  \label{eq:eigmoments}
\end{align}
where $\pmean,\pvar$ are the per-training-point posterior-moment vectors and
$\odot$ is the element-wise product; the two variance expressions coincide
because $\operatorname{rowsum}(\amat\odot(\amat M))=\diago(\amat M\amat\transpose)$
for any $M$. Term by term this equals \eqref{eq:predmoments}: $\Kvec_i=\Kf(x_i,x_i)$
(the self-kernel identity \eqref{eq:selfkernel}) and
$\amat_i(\Vb-\Ktilb)\amat_i\transpose=\kvec(x_i)\transpose\Ktil^{-1}(\vV-\Ktil)\Ktil^{-1}\kvec(x_i)$.
For \textbf{test} points the same two expressions run on a freshly projected
$\amat=\Kf(x_\ast,\Ipt)\eb/\evals$ with $\Kvec=\Kf(x_\ast,x_\ast)$, so no
training-point kernel is recomputed; the predicted rate there is
$\rate_{\mathrm{pred}}=\exp(\gain\pmean+\tfrac12\gain^2\pvar+\latz)$, i.e.\
\eqref{eq:fbar}.

\subsection{ELBO and KL in the eigenbasis}\label{sec:eig-kl}

The eigenspace ELBO uses the log-likelihood \eqref{eq:ell} and, because
$\Ktilb=\diago(\evals)$, the KL \eqref{eq:kl} collapses to sums over the kept
dimensions:
\begin{align}
  \KLdiv(q\,\|\,p)
    &= \tfrac12\Big[\Tr(\Ktil^{-1}\vV)+\vm\transpose\Ktil^{-1}\vm-\nb
        +\log\lvert\Ktil\rvert-\log\lvert\vV\rvert\Big] \notag\\
    &= \tfrac12\Big[\ \textstyle\sum_j \dfrac{\Vb[j,j]}{\evals_j}
        \;+\; \sum_j \dfrac{\mb[j]^2}{\evals_j}
        \;-\; \nb
        \;+\; \sum_j \log\evals_j
        \;-\; \log\det\Vb\ \Big].
  \label{eq:eigkl}
\end{align}

\begin{remark}[Diagonal-trace assumption]
The trace term uses only $\diago(\Vb)$, which is valid \emph{only because}
$\Ktilb$ is diagonal in the basis $\eb$ in which the ELBO is evaluated --- i.e.\
when the kernel and the basis are mutually consistent. The $\log\lvert\vV\rvert$
term uses the full log-determinant of the dense $\Vb$. The $-\nb$ is a
zero-gradient constant (cf.\ the $-\Mind$ remark after \eqref{eq:kl}): it shifts
only the reported loss magnitude.
\end{remark}

% ============================================================================
%  Part II (training) -- Section 6: The training algorithm
%  \input-ready fragment. Uses only macros from preamble.tex.
%  Cross-part refs (defined elsewhere): eq:elbo, eq:predmoments, eq:fbar
%  (Part II inference).
% ============================================================================

\section{The training algorithm: EM-style coordinate ascent}\label{sec:training}

The model assembled in Parts~I--II has three groups of free parameters, and all
three sit on the \emph{single} evidence lower bound $\ELBO$ of
\eqref{eq:elbo}:

\begin{enumerate}[label=(\alph*),leftmargin=2.2em,itemsep=1pt]
  \item the \textbf{variational moments} $\mb,\Vb$ of the eigenspace posterior
        $q(\latt)=\Gauss(\vm,\vV)$;
  \item the \textbf{rate parameters} of the Poisson likelihood, the gain
        $\gain$ and the offset $\latz$, defining
        $\rate=\exp(\gain\lat+\latz)$;
  \item the \textbf{kernel hyperparameters}
        $\{\sigz,\Amp,\bwid,\smoo,\varepsilon_{0x},\varepsilon_{0y}\}$ of the
        structured arc-cosine prior (Part~I).
\end{enumerate}

Training maximises the one bound $\ELBO$ by \textbf{block coordinate ascent}:
each step improves one block while holding the other two fixed. The three steps
are the E-step (variational moments), the F-step (rate parameters), and the
M-step (kernel hyperparameters), run once per outer iteration in a fixed order
(\S\ref{subsec:training-overview}).

% ---------------------------------------------------------------------------
\subsection{Overview: three coordinate blocks, one ELBO, one loop}
\label{subsec:training-overview}

Each step optimises one block and freezes the rest:

\begin{center}\small
\begin{tabular}{@{}l l l >{\raggedright\arraybackslash}p{2.35in}@{}}
\toprule
Step & Free block & Held fixed & Method \\
\midrule
E & $\mb,\Vb$ & kernel; $\gain,\latz$; $\eb$ &
    closed-form Newton (no automatic differentiation) \\
F & $\gain$ ($\latz$ closed-form) & $\mb,\Vb$ (via $\pmean,\pvar$); kernel &
    LBFGS on $\log\gain$, or damped Newton \\
M & $\{\sigz,\Amp,\bwid,\smoo,\rfc\}$ & $\mb,\Vb$; $\gain,\latz$; $\eb$ &
    LBFGS $+$ gradients (analytical or automatic) \\
\bottomrule
\end{tabular}
\end{center}

\noindent(In the M-row, the centre $\rfc=(\varepsilon_{0x},\varepsilon_{0y})$ of
the localized region counts as the two coordinate hyperparameters, so the free
block has six scalars.)

\paragraph{Loop order (one outer iteration).}
Per outer iteration:

\begin{enumerate}[leftmargin=2.4em,itemsep=1pt]
  \item \textbf{Re-sync the eigenbasis.} Rebuild the eigenbasis $\eb$ from the
        hyperparameters the \emph{previous} M-step moved, then refresh
        $\pmean,\pvar,\fbar$ --- this is where the M-step's deferred $\eb$ update
        lands.
  \item \textbf{E-step} --- one or more closed-form Newton steps on
        $\mb,\Vb$ (\S\ref{subsec:estep}).
  \item \textbf{F-step} --- fit $\gain$ (and closed-form $\latz$), unless the
        F-step is interleaved into the E-loop (\S\ref{subsec:fstep}).
  \item \textbf{Metrics $+$ ELBO $+$ early stopping} --- computed
        \emph{deliberately before} the M-step.
  \item \textbf{M-step} --- update the kernel hyperparameters
        (\S\ref{subsec:mstep}); the $\eb$ update is \emph{deferred} to step~1 of
        the next iteration.
\end{enumerate}

\begin{remark}[Why metrics precede the M-step]
The M-step changes the kernel parameters \emph{without} recomputing the
eigenbasis $\eb$ (that is deferred, for cost and stability). Every logged
quantity --- the ELBO, the early-stopping metric --- must therefore be read while
the model state is still internally consistent, i.e.\ \emph{before} the M-step
perturbs the kernel. Early stopping is on the ELBO, with patience and
best-restore.
\end{remark}

% ---------------------------------------------------------------------------
\subsection{The E-step: a Newton update of the variational posterior}
\label{subsec:estep}

At \emph{fixed} kernel and fixed $\gain,\latz$, the E-step maximises the ELBO
\eqref{eq:elbo} over the variational moments. Writing the inducing-space
posterior $q(\latt)=\Gauss(\vm,\vV)$ against the prior $p(\latt)=\Gauss(0,\Ktil)$,

\begin{equation}
\ELBO(\vm,\vV)
  = \underbrace{\sum_i \E_{q(\lat_i)}\!\big[\log\Poi(\spk_i\mid\lat_i)\big]}_{\ELBO_{\mathrm{data}}}
    \;-\; \KLdiv\!\big(q(\latt)\,\|\,p(\latt)\big),
\label{eq:estep-obj}
\end{equation}
with the projected posterior moments $\pmean_i,\pvar_i$ of \eqref{eq:predmoments}
and the KL contributing, through the E-step variables,
$-\tfrac12\vm\transpose\Ktil^{-1}\vm$ (couples to $\vm$) and
$\tfrac12\log|\vV|-\tfrac12\Tr(\Ktil^{-1}\vV)$ (couples to $\vV$).

\paragraph{Poisson expected log-likelihood and its moment derivatives.}
The data term is closed-form via the Gaussian moment-generating function, with
the expected rate $\fbar_i$ of \eqref{eq:fbar}:
\begin{equation}
\ELBO_{\mathrm{data}}
  = \sum_i\big[\,\spk_i(\gain\pmean_i+\latz)-\fbar_i\,\big],
\qquad
\fbar_i=\exp\!\big(\gain\pmean_i+\tfrac12\gain^2\pvar_i+\latz\big),
\end{equation}
(the $-\log\spk_i!$ constant is dropped). The three point-derivatives that
assemble the Newton step follow by the chain rule on $\fbar_i$:
\begin{equation}
\frac{\partial\fbar_i}{\partial\pmean_i}=\gain\,\fbar_i,\qquad
\frac{\partial\fbar_i}{\partial\pvar_i}=\tfrac12\gain^2\,\fbar_i,\qquad
\frac{\partial}{\partial\pmean_i}\,\E_{q}[\log\Poi]=\gain(\spk_i-\fbar_i).
\end{equation}

\paragraph{The two Newton helpers.}
Accumulating those derivatives through the projection $\pmean_i=\kvec_i\transpose\Ktil^{-1}\vm$
defines the gradient helper $\gE$ and the (Poisson-Fisher) curvature $\GE$:
\begin{equation}
\gE \;=\; \gain\sum_i \kvec_i(\spk_i-\fbar_i)\ \in\Reals^{\Mind},
\qquad
\GE \;=\; \gain^2\sum_i \fbar_i\,\kvec_i\kvec_i\transpose\ \in\Reals^{\Mind\times\Mind}.
\label{eq:estep-gG}
\end{equation}
$\gE$ is the kernel-projected mismatch between observed and expected counts;
$\GE$ is the Poisson Fisher information and is PSD because every $\fbar_i>0$.
(Note $\gE$ carries the subscript $\mathrm{E}$ specifically to keep it distinct
from the log-rate $\lgr=\gain\lat+\latz$ of the utility part.)

\paragraph{The $\vV$-update (closed form).}
With $\pmean_i$ independent of $\vV$, setting $\nabla_{\vV}\ELBO=0$ gives
$\vV^{-1}=\Ktil^{-1}(\Ktil+\GE)\Ktil^{-1}$, hence
\begin{equation}
\boxed{\;\vV \;=\; \Ktil\,(\Ktil+\GE)^{-1}\,\Ktil\;}
\label{eq:estep-V}
\end{equation}
--- a \emph{symmetric} sandwich of two symmetric-PD factors, so $\vV$ is a valid
covariance. This symmetry is the substance of the ``corrected'' E-step: the
superseded form $\vV=(\Ktil+\GE)^{-1}\Ktil$ is a product of two non-commuting
symmetric matrices and is therefore \emph{not} symmetric (Cholesky/PSD failures).

\paragraph{The $\vm$-update.}
The exact Newton step $\vm\leftarrow\vm-H_{\vm}^{-1}\nabla_{\vm}\ELBO$, with
$\nabla_{\vm}\ELBO=\Ktil^{-1}(\gE-\vm)$ and Hessian
$H_{\vm}=-\Ktil^{-1}(\Ktil+\GE)\Ktil^{-1}$, gives the corrected form
$\vm\leftarrow\vm+\Ktil(\Ktil+\GE)^{-1}(\gE-\vm)$. The map actually used is
\emph{not} this; it is the algebraically older
\begin{equation}
\boxed{\;\vm \;\leftarrow\; \Ktil\,(\Ktil+\GE)^{-1}\big(\GE\,\Ktil^{-1}\vm+\gE\big)\;}
\label{eq:estep-m}
\end{equation}
(equal to $\vV_{\mathrm{new}}(\GE\Ktil^{-1}\vm+\gE)$ with $\vV_{\mathrm{new}}$ from
\eqref{eq:estep-V}). The two maps differ only in the first term
($\Ktil(\Ktil+\GE)^{-1}\GE\Ktil^{-1}\vm$ vs.\ $\GE(\Ktil+\GE)^{-1}\vm$); they
coincide iff $\Ktil$ and $\GE$ commute.

\caveat{The $\vm$-update map \eqref{eq:estep-m} differs \emph{per iteration} from
the corrected Newton step whenever $\Ktil$ and $\GE$ do not commute, but both
maps share the same fixed point $\vm^\star=\gE$, so the \emph{converged}
posterior is identical --- the discrepancy is a transient path difference, not a
different solution. It is left as-is (it converges in practice). The $\vV$-update
\eqref{eq:estep-V} is unambiguously correct, and it is $\vV$ that the
active-learning variance estimates depend on.}

\begin{remark}[Fixed-point equivalence]
Let $T_{\mathrm N}(\vm)=\vm+\Ktil(\Ktil+\GE)^{-1}(\gE-\vm)$ (corrected) and
$T_{\mathrm C}(\vm)=\Ktil(\Ktil+\GE)^{-1}(\GE\Ktil^{-1}\vm+\gE)$ (older map). At
$\vm=\gE$:
$T_{\mathrm N}(\gE)=\gE$, and
$T_{\mathrm C}(\gE)=\Ktil(\Ktil+\GE)^{-1}(\GE\Ktil^{-1}+I)\gE
   =\Ktil(\Ktil+\GE)^{-1}(\GE+\Ktil)\Ktil^{-1}\gE=\Ktil\Ktil^{-1}\gE=\gE$.
Both are valid fixed-point iterations onto the same stationary condition
$\vm^\star=\gE=\gain\sum_i\kvec_i(\spk_i-\fbar_i)$.
\end{remark}

\subsubsection{Why the update is cheap: the eigenspace realisation}
\label{subsubsec:estep-eigen}

The E-step never touches the $\Mind$-dimensional inducing space directly. It works
in the eigenbasis $\eb$ of $\Ktil$, keeping only the $\nb$ leading directions
(eigenvalues above $\max(\lambda_{\max}\!\cdot\!\mathrm{tol},\,\mathrm{tol})$,
$\mathrm{tol}=10^{-4}$; in practice $\nb\approx10\text{--}11$ while $\Mind$ may be
large). With $\Ktil=\eb\evals\eb\transpose$ the projected gram
$\Ktilb=\eb\transpose\Ktil\eb=\diago(\evals)$ is \emph{diagonal}, so $\Ktil^{-1}$
is trivial. Using the eigenspace projection $\amat=K\Ktil^{-1}$
(shape $\Ntr\times\nb$), one forms the \emph{transformed} helpers:
\begin{align}
g_{\mathrm E,b} &= \gain\,\amat\transpose(\spk-\fbar)\;=\;\Ktil^{-1}\gE
   &&(\text{eigenbasis coords}),\\
G_{\mathrm E,b} &= \gain^2\,\amat\transpose\big(\fbar\odot\amat\big)\;=\;\Ktil^{-1}\GE\Ktil^{-1},
\end{align}
and applies \eqref{eq:estep-V}--\eqref{eq:estep-m} entirely in the tiny
$\nb$-space:
\begin{equation}
\Vb \leftarrow \big(I+\Ktilb\,G_{\mathrm E,b}\big)^{-1}\Ktilb,\qquad
\mb \leftarrow \Vb\big(G_{\mathrm E,b}\,\mb+g_{\mathrm E,b}\big),\qquad
\Vb \leftarrow \tfrac12(\Vb+\Vb\transpose).
\end{equation}
The first line equals \eqref{eq:estep-V} exactly, via the identity
$(I+PQ^{-1})^{-1}=Q(Q+P)^{-1}$ with $P=\GE,Q=\Ktil$:
$(I+\Ktilb G_{\mathrm E,b})^{-1}\Ktilb=(I+\GE\Ktil^{-1})^{-1}\Ktil
=\Ktil(\Ktil+\GE)^{-1}\Ktil$. Prior initialisation is $\mb=0$,
$\Vb=\Ktilb=\evals$. The dimensional collapse $\Mind\!\to\!\nb\approx10$ ---
not any special-casing --- is what makes the Newton step cheap.

\subsubsection{Iteration, damping, and the divergence guard}

The E-step performs a \emph{single, undamped, closed-form} Newton step; iteration
and stabilisation are handled around it:

\begin{itemize}[leftmargin=1.4em,itemsep=1pt]
  \item \textbf{Coupling variable $\fbar$.} After each step the posterior is
        recomputed and $\fbar=\exp(\gain\pmean+\tfrac12\gain^2\pvar+\latz)$
        refreshed, then fed into the next step's $g_{\mathrm E,b},G_{\mathrm E,b}$.
        $\gain,\latz$ are held fixed during the E-step.
  \item \textbf{Revert-on-divergence guard (the effective damping).} Each iteration
        snapshots $(\mb,\Vb)$; if the post-step $\fbar$ mean or max exceeds a
        threshold, or $\fbar$ has NaNs, it \emph{reverts} $(\mb,\Vb)$ (and
        $\gain,\latz$ if interleaved), recomputes $\fbar$, and breaks the loop.
  \item \textbf{Optional interleaved F-step.} With interleaving enabled, a
        \emph{damped} Newton update of $(\gain,\latz)$ (\S\ref{subsec:fstep},
        $\alpha=0.25$) runs after every E-iteration, so the rate parameters
        track $(\mb,\Vb)$ as they move.
\end{itemize}

% ---------------------------------------------------------------------------
\subsection{The F-step: fitting the rate function}
\label{subsec:fstep}

Here ``F'' denotes the \textbf{rate function} $\rate$: the F-step fits the
Poisson-rate parameters $\gain,\latz$ that define
$\fbar=\exp(\gain\pmean+\tfrac12\gain^2\pvar+\latz)$, holding the variational
moments (through frozen $\pmean,\pvar$) and the kernel fixed. Here $\gain$ is the
only searched variable; $\latz$ is solved in closed form for the current $\gain$.

\paragraph{Closed-form $\latz$ given $\gain$.}
From $\partial\ELBO/\partial\latz=0$ on
$\sum_i[\spk_i\latz-\exp(\latz)\exp(\gain\pmean_i+\tfrac12\gain^2\pvar_i)]$:
\begin{equation}
\boxed{\;
\latz \;=\; \log\!\Big(\textstyle\sum_i\spk_i\Big)
       \;-\; \log\!\Big(\textstyle\sum_i\exp\!\big(\gain\pmean_i+\tfrac12\gain^2\pvar_i\big)\Big)
\;}
\label{eq:fstep}
\end{equation}
(requires $\sum_i\spk_i>0$; a zero-count target is a data problem). It is
re-evaluated whenever $\gain$ changes.

\paragraph{LBFGS on $\gain$ with an analytic gradient.}
The F-step runs LBFGS over $\log\gain$, re-solving $\latz$ at each evaluation and
using the closed-form derivative, with $\pmean,\pvar$ the frozen inputs:
\begin{equation}
\frac{\partial\ELBO}{\partial\gain}
  = \sum_i\big[\spk_i\pmean_i-(\pmean_i+\gain\pvar_i)\,\fbar_i\big],
\qquad
\frac{\partial\ELBO}{\partial(\log\gain)}=\gain\frac{\partial\ELBO}{\partial\gain}.
\end{equation}
Since LBFGS minimises $-\ELBO$, the gradient it consumes with respect to
$\log\gain$ is $-\gain\,\partial\ELBO/\partial\gain$. Guards reject a trial step
whose $\latz$ over/underflows or whose $\fbar$ exceeds the mean/max thresholds,
and revert $(\gain,\latz)$ if the final $\latz$ overflows.

\paragraph{Interleaved joint damped Newton on $(\gain,\latz)$.}
With interleaving enabled, a joint damped Newton step on $\psi=[\gain,\latz]$ runs
\emph{inside} the E-loop. With per-input
$\rate_i=\exp(\gain\pmean_i+\tfrac12\gain^2\pvar_i+\latz)$ and
$d_i=\pmean_i+\gain\pvar_i$:
\begin{equation}
R=\begin{bmatrix}\sum_i(\spk_i\pmean_i-d_i\rate_i)\\[2pt]\sum_i(\spk_i-\rate_i)\end{bmatrix},
\quad
H=-\begin{bmatrix}\sum_i(\pvar_i\rate_i+d_i^2\rate_i)&\sum_i d_i\rate_i\\[2pt]
                  \sum_i d_i\rate_i&\sum_i\rate_i\end{bmatrix},
\quad
\psi\leftarrow\psi-\alpha\,H^{-1}R,\ \ \alpha=0.25.
\end{equation}
$H$ is negative-definite (the log-likelihood is concave in $\psi$); the damping
$\alpha=0.25$ and the same $\fbar$-threshold rejections stabilise the step.

% ---------------------------------------------------------------------------
\subsection{The M-step: updating the kernel hyperparameters}
\label{subsec:mstep}

The M-step maximises the ELBO \eqref{eq:elbo} over the kernel/prior
hyperparameters while $\mb,\Vb$, the rate parameters $\gain,\latz$, and
\emph{the eigenbasis $\eb$} are held fixed. The free hyperparameters (all owned
by the arc-cosine prior) are the six
\begin{equation}
\theta\in\{\ \sigz,\ \Amp,\ \bwid,\ \smoo,\ \varepsilon_{0x},\ \varepsilon_{0y}\ \}.
\end{equation}
The eigenbasis $\eb$ is held fixed as an approximation (for cost and stability)
and re-synced at the start of the next outer iteration. This block is solved by
gradient-based optimisation (e.g.\ LBFGS on the raw, unconstrained parameters);
the hyperparameter gradients $\partial_\theta\ELBO$ are obtained by automatic or
analytical differentiation of the eigenspace ELBO.

\begin{remark}[The amplitude $\Amp$ is optimized by default]
The amplitude prefactor $\Amp$ of the structured prior
$\Cmat=\Amp\,\loc\,\Csm\,\loc\transpose$ (Part~I) is a genuine sixth kernel
hyperparameter, not a fixed constant, and is optimized alongside the other five
by default. The optional flag \texttt{fix\_Amp} (default off) instead pins
$\Amp=1$, recovering the amplitude-free model; left at its default, $\Amp$ trains
with the rest.
\end{remark}


\part{Active learning: the information utility}
% ============================================================================
%  math_version/sections/part3a_utility_core.tex  --  Part III, Section 7
%  "The acquisition objective" (utility core).
%  \input-ready fragment: no preamble/documentclass. Uses the shared macros
%  from preamble.tex. Domain-neutral (pure math / ML) rendering; symbols and
%  mathematics identical to the source. The deep proofs (conditioning,
%  divergence, predictive-mean equivalence) live in the sibling Part III-deep
%  section and are only cited here.
% ============================================================================

\section{The acquisition objective}\label{sec:utility}

The active-learning loop proceeds sequentially: after each count observation the
model is refit, and the next input to query is the one whose observation is
expected to teach us the most about the latent function. ``Teach us the most'' is
made precise as an \emph{information gain} --- an expected reduction in the entropy
of what we do not yet know about the latent function and its rate. This section
builds that utility on the generative model and variational posterior of
Parts~I--II: the Poisson-lognormal predictive of the count, its entropy, and the
two utilities considered here.

Throughout, the candidate (query) input is $x^{*}$; a count is the random variable
$R$ with realizations $r$ (this is the count written $\spk$ elsewhere in the
document --- we keep it as $r$ here to match the count-domain notation $\rmax$ and
the Laplace mode $\gmode$). The mean rate at input $x$ is
$\rate(x)=\exp(\gain\,\lat(x)+\latz)$ with latent $\lat(x)$, gain $\gain$, and
offset $\latz$; the count is $R\mid\rate(x)\sim\Poi(\rate(x))$. It is convenient to
work with the \emph{log-rate} $\lgr(x)=\gain\,\lat(x)+\latz$, so $\rate=e^{\lgr}$,
whose GP posterior at any input is Gaussian with the log-rate moments
%
\begin{equation}
  \mug \;=\; \gain\,\pmean(x^{*})+\latz,
  \qquad
  \varg \;=\; \gain^{2}\,\pvar(x^{*}),
  \label{eq:logratemoments}
\end{equation}
%
the affine push-forward of the latent posterior moments $\pmean(x^{*})$,
$\pvar(x^{*})$ read off the variational posterior~\eqref{eq:predmoments}.

\begin{remark}[Two utilities, one idea]\label{rem:two-utilities}
We consider \emph{two} acquisition functions. Both express the same idea ---
\emph{utility$(x^{*})=$ information the observation $R$ at $x^{*}$ gives about the
latent function $=$ a drop in entropy} --- and they differ only in \emph{which}
uncertainty they score (Table~\ref{tab:two-utilities}). The \textbf{standard
utility} $\Ustd$ scores the rate \emph{at $x^{*}$ itself}; the
\textbf{distribution-aware utility} $\Uda$ scores the rates at the inputs
$x\sim\px$ we ultimately want to predict. The standard utility depends only on the
marginal posterior at $x^{*}$; the distribution-aware utility couples $x^{*}$ to the
input distribution $\px$ through cross-covariance. $\Uda$ is the richer object, but
it carries a norm-exploitation divergence (\S\ref{subsec:landscape},
Thm.~\ref{thm:divergence}) that a finite candidate pool masks.
\end{remark}

\begin{table}[t]
  \centering
  \small
  \begin{tabular}{@{}lll@{}}
    \toprule
      & \textbf{Standard} $\Ustd$ & \textbf{Distribution-aware} $\Uda$\\
    \midrule
    Formula      & $\Hmarg(x^{*})-\E[\Hnoise(x^{*})]$
                 & $\Hmarg(x^{*})-\E_{x\sim\px}[\Hcond(x^{*})]$\\
    As MI        & $I\!\big(\rate(x^{*});R\mid x^{*},\mathcal D\big)$
                 & $I\!\big(\rate(X);R\mid X,x^{*},\mathcal D\big),\ X\sim\px$\\
    Scores       & rate \emph{at $x^{*}$ itself}
                 & rates at \emph{inputs} $x\sim\px$\\
    Needs        & marginal moments at $x^{*}$
                 & marginal moments \emph{and} cross-cov.\ to $\px$\\
    \bottomrule
  \end{tabular}
  \caption{The two utilities. Both are ``information gain $=$ entropy reduction'',
  about different targets.}
  \label{tab:two-utilities}
\end{table}

% ---------------------------------------------------------------------------
\subsection{Utility as expected entropy reduction and mutual information}
\label{subsec:mi}

Fix a candidate $x^{*}$ and the data $\mathcal D=\{(x_i,r_i)\}$. The utility is the
mutual information between the count $R$ we would observe at $x^{*}$ and the latent
function we want to learn. Precisely what we want to learn about is the difference
between the two utilities.

\paragraph{The distribution-aware objective.}
Let $Z:=(X,\rate(X))$ with $X\sim\px$ an input drawn from the input distribution
and $\rate(X)$ its rate. Choosing $x^{*}$ to maximize the information $R$ carries
about $Z$,
%
\begin{equation}
  \Uda(x^{*}) \;:=\; I\!\big(Z;R\mid x^{*},\mathcal D\big)
             \;=\; I\!\big((X,\rate(X));\,R\mid x^{*},\mathcal D\big).
  \label{eq:da-objective}
\end{equation}
%
The chain rule for mutual information splits this into a term for \emph{which} input
we imagine and a term for its \emph{rate}:
%
\begin{equation}
  I\!\big((X,\rate(X));R\mid x^{*},\mathcal D\big)
  \;=\;
  \underbrace{I(X;R\mid x^{*},\mathcal D)}_{=\,0}
  \;+\;
  I\!\big(\rate(X);R\mid X,x^{*},\mathcal D\big).
  \label{eq:mi-chain}
\end{equation}
%
The first term vanishes: \emph{which} input $X$ we intend to ask about is
independent of the count $R$ observed at $x^{*}$, given $x^{*}$ and $\mathcal D$.
Hence
%
\begin{equation}
  \Uda(x^{*})
  \;=\; I\!\big(\rate(X);R\mid X,x^{*},\mathcal D\big)
  \;=\; \E_{x\sim\px}\!\big[\,I\!\big(\rate(x);R\mid x,x^{*},\mathcal D\big)\big].
  \label{eq:da-expectation}
\end{equation}

\paragraph{Entropy decomposition.}
For a fixed input $x$, the per-input mutual information decomposes into a
marginal-entropy term minus an expected conditional-entropy term,
%
\begin{equation}
  I\!\big(\rate(x);R\mid x,x^{*},\mathcal D\big)
  \;=\;
  H(R\mid x^{*},\mathcal D)
  \;-\;
  \E_{\rate(x)}\!\big[\,H(R\mid x^{*},\rate(x),\mathcal D)\big],
  \label{eq:mi-entropy-split}
\end{equation}
%
using (i) $H(R\mid x,x^{*},\mathcal D)=H(R\mid x^{*},\mathcal D)$ --- the marginal
of $R$ at $x^{*}$ does not depend on which input we intend to predict --- and
(ii) that conditioning on the latent value $\rate(x)$ leaves the expected entropy
$\E_{\rate(x)}[H(R\mid x^{*},\rate(x),\mathcal D)]$. Averaging
\eqref{eq:mi-entropy-split} over $x\sim\px$ and naming the two pieces gives the
distribution-aware utility as a difference of entropies,
%
\begin{equation}
  \Uda(x^{*}) \;=\; \Hmarg(x^{*}) \;-\; \E_{x\sim\px}\!\big[\Hcond(x^{*})\big],
  \label{eq:Uda}
\end{equation}
%
with
%
\begin{align}
  \Hmarg(x^{*}) &:= H(R\mid x^{*},\mathcal D), \label{eq:def-Hmarg-cond}\\
  \Hcond(x^{*}) &:= \E_{x\sim\px,\ \lat(x)\sim q}\!\big[\,H(R\mid x^{*},\lat(x),\mathcal D)\big].
  \label{eq:def-Hcond}
\end{align}
%
By the symmetry of mutual information, $\Uda(x^{*})$ equals the reduction in
uncertainty about the rates $\{\rate(x):x\sim\px\}$ produced by observing $R$ at
$x^{*}$ --- the precise sense of ``how much does querying $x^{*}$ teach us about the
latent function over the input distribution.''

\begin{remark}[Two expectations inside $\Hcond$]\label{rem:double-exp}
$\Hcond$ in~\eqref{eq:def-Hcond} carries \emph{two} expectations: over inputs
$x\sim\px$ and over the posterior latent $\lat(x)\sim q$. The latent is integrated
because $\rate(x)=\exp(\gain\lat(x)+\latz)$ is itself uncertain --- only the
posterior of $\lat(x)$ is known, not its value. Replacing the inner expectation by
an evaluation at the mean $\lat(x)=\pmean(x)$ is \emph{biased}; the necessity of the
inner sampling is the subject of \S\ref{subsec:why-sample}.
\end{remark}

The \emph{standard} utility instead scores the rate at the query point itself. It
is the special case in which the ``input we predict'' is $x^{*}$, so $\Uda$'s outer
expectation collapses and $\Hcond$ becomes the aleatoric Poisson noise entropy at
$x^{*}$; we develop it in \S\ref{subsec:standard} once $\Hmarg$ is in hand.

% ---------------------------------------------------------------------------
\subsection{The marginal count entropy \texorpdfstring{$\Hmarg$}{H\_marg}}
\label{subsec:hmarg}

Both utilities share the marginal predictive entropy of the count at $x^{*}$,
%
\begin{equation}
  \Hmarg(x^{*})
  \;=\; -\sum_{r=0}^{\infty} p(r\mid x^{*},\mathcal D)\,\log p(r\mid x^{*},\mathcal D),
  \label{eq:Hmarg}
\end{equation}
%
where the marginal predictive integrates the Poisson likelihood against the
Gaussian posterior of the log-rate $\lgr(x^{*})\sim\Gauss(\mug,\varg)$
(cf.~\eqref{eq:logratemoments}):
%
\begin{equation}
  p(r\mid x^{*},\mathcal D)
  \;=\; \int \Poi(r\mid e^{\lgr})\,\Gauss(\lgr\mid\mug,\varg)\,d\lgr
  \;=\; \frac{1}{r!}\int \exp\!\big(r\lgr - e^{\lgr}\big)\,\Gauss(\lgr\mid\mug,\varg)\,d\lgr .
  \label{eq:marg-predictive}
\end{equation}
%
This integral has no closed form. It is evaluated by a \textbf{Laplace expansion} of
the integrand about its mode $\gmode$ (the mode of the log-integrand
$r\lgr - e^{\lgr} - (\lgr-\mug)^{2}/2\varg$):
%
\begin{equation}
  \log p(r\mid x^{*},\mathcal D)
  \;\approx\;
  \gmode\,r - e^{\gmode}
  - \frac{(\gmode-\mug)^{2}}{2\varg}
  - \tfrac12\log\!\big(1+\varg\,e^{\gmode}\big)
  - \log r! .
  \label{eq:laplace-logp}
\end{equation}

\begin{lemma}[Closed-form Laplace mode via Lambert~$W$]\label{lem:mode}
The stationarity condition $r-e^{\lgr}-(\lgr-\mug)/\varg=0$ rearranges to
$\lgr=r\varg+\mug-\varg e^{\lgr}$. Substituting $w=\varg e^{\lgr}$ gives
$w+\log w=\log\varg+r\varg+\mug$, i.e.\ $w=\LW\!\big(\varg\,e^{\,r\varg+\mug}\big)$
with $\LW$ the principal branch of the Lambert~$W$ function. Hence the mode is
closed-form,
%
\begin{equation}
  \gmode \;=\; r\varg+\mug-\LW\!\big(\varg\,e^{\,r\varg+\mug}\big).
  \label{eq:laplace-mode}
\end{equation}
\end{lemma}

The infinite sum in~\eqref{eq:Hmarg} decays fast (rates are low) and is truncated at
$\rmax$ (\S\ref{subsec:landscape} handles the truncation artifact). At very small
variance the Laplace form is unstable, and one falls back to the exact Poisson
log-pmf $\log p(r)=r\mug-e^{\mug}-\log r!$ (used for $\varg<10^{-6}$). The mode is
computed in log-space --- $\LW(e^{y})$ with $y=\log\varg+r\varg+\mug$,
Newton-iterated with the derivative $dW/dy=W/(1+W)$ --- so no $e^{r\varg+\mug}$ is
ever formed (overflow-safe and differentiable in $x^{*}$).

% ---------------------------------------------------------------------------
\subsection{The standard utility}
\label{subsec:standard}

The standard acquisition is the mutual information between the count and the latent
rate \emph{at the query point itself},
%
\begin{equation}
  \Ustd(x^{*})
  \;=\; \Hmarg(x^{*}) - \E[\Hnoise(x^{*})]
  \;=\; H(R\mid x^{*},\mathcal D) - \E_{\rate\sim\text{post.}}\!\big[H(R\mid\rate)\big]
  \;=\; I\!\big(\rate(x^{*});R\mid x^{*},\mathcal D\big).
  \label{eq:Ustd}
\end{equation}
%
It scores ``how much would a count at $x^{*}$ pin down the rate at $x^{*}$'' --- it
depends on the marginal GP moments at $x^{*}$ only, with no $\px$ and no
cross-covariance. It is exactly $\Uda$ with the predicted input taken to be the
query point, so the second term is the \emph{aleatoric} Poisson noise entropy rather
than an epistemic conditional entropy.

\subsubsection{The aleatoric noise entropy \texorpdfstring{$\E[\Hnoise]$}{E[H\_noise]} in closed form}
\label{subsubsec:hnoise}

Given a rate $\rate$, the entropy of $\Poi(\rate)$ is
$\Hnoise(\rate)=\rate\,(1-\log\rate)+\E_{r\sim\Poi(\rate)}[\log r!]$. Writing
$\lgr=\log\rate\sim\Gauss(\mug,\varg)$ and using the lognormal moment identities
%
\begin{equation}
  \E[e^{\lgr}] = \exp\!\big(\mug+\tfrac12\varg\big) = \fbar(x^{*}),
  \qquad
  \E[\lgr\,e^{\lgr}] = \exp\!\big(\mug+\tfrac12\varg\big)\,(\mug+\varg),
  \label{eq:lognormal-identities}
\end{equation}
%
(the first is the expected rate $\fbar$ of~\eqref{eq:fbar}), the first piece
integrates in closed form:
%
\begin{equation}
  \E_{\lgr}\!\big[\rate(1-\log\rate)\big]
  = \E_{\lgr}\!\big[e^{\lgr}(1-\lgr)\big]
  = \exp\!\big(\mug+\tfrac12\varg\big)\big(1-(\mug+\varg)\big)
  = -\exp\!\big(\mug+\tfrac12\varg\big)(\mug+\varg-1).
  \label{eq:hnoise-firstpiece}
\end{equation}
%
The remaining $\E[\log r!]$ is taken under the marginal predictive
$p(r\mid x^{*},\mathcal D)$ of~\eqref{eq:marg-predictive}. Hence
%
\begin{equation}
  \E[\Hnoise(x^{*})]
  \;=\; -\exp\!\big(\mug+\tfrac12\varg\big)(\mug+\varg-1)
        \;+\; \sum_{r} p(r\mid x^{*},\mathcal D)\,\log r! .
  \label{eq:Hnoise}
\end{equation}
%
This is the closed form for $\E[\Hnoise]$.
Assembling $\Ustd$ then reads the marginal latent moments
$\pmean(x^{*}),\pvar(x^{*})$, transforms them to $\mug,\varg$
by~\eqref{eq:logratemoments}, and returns $\Ustd=\Hmarg-\E[\Hnoise]$.

% ---------------------------------------------------------------------------
\subsection{The distribution-aware utility}
\label{subsec:da}

The distribution-aware utility~\eqref{eq:Uda} reduces uncertainty about the rates at
the inputs we actually want to predict. It weights a candidate $x^{*}$ by how
informative its observation is about the rates at typical inputs --- i.e.\ by the
cross-covariance between $x^{*}$ and inputs $x\sim\px$ (``distribution-aware'' $=$
aware of the input distribution $\px$). Evaluating $\Hcond$ needs the predictive
moments of $\lat(x^{*})$ \emph{after a hypothetical observation} $\lat(x)=\lat_i$ at
an input $x$.

\subsubsection{Conditional GP moments via Gaussian conditioning}
\label{subsubsec:condmoments}

Let $\pv(x)=\Ktil^{-1}\kvec(x)$ be the SVGP projection and $q(\latt)=\Gauss(\vm,\vV)$
the variational posterior on the inducing values. The marginal latent posterior
moments at $x$ and $x^{*}$ are the projected moments of~\eqref{eq:predmoments},
$\pmean(x)=\pv(x)\transpose\vm$, $\pvar(x)=s(x)+\pv(x)\transpose\vV\pv(x)$
(and likewise at $x^{*}$), where $s(x)=\kvec(x,x)-\pv(x)\transpose\kvec(x)$ is
the prior Schur complement. The object that couples the two points is the
\textbf{prior conditional covariance} of $\lat(x),\lat(x^{*})$ given the inducing
values:
%
\begin{equation}
  \ccond \;=\; \kvec(x,x^{*}) - \pv(x)\transpose\Ktil\,\pv(x^{*}) .
  \label{eq:util-ccond}
\end{equation}
%
By the law of total covariance the full posterior cross-covariance is
%
\begin{equation}
  \Sigma_{x^{*}} \;=\; \ccond + \pv(x)\transpose\vV\,\pv(x^{*})
             \;=\; \kvec(x,x^{*}) + \pv(x)\transpose(\vV-\Ktil)\,\pv(x^{*}) ,
  \label{eq:util-crosscov}
\end{equation}
%
and standard Gaussian conditioning on the joint of
$(\lat(x),\lat(x^{*}))$ gives the conditional latent moments at $x^{*}$ after
observing $\lat(x)=\lat_i$,
%
\begin{align}
  \mu_{\mathrm{cond}}(x^{*})     &= \pmean(x^{*}) + \frac{\Sigma_{x^{*}}}{\pvar(x)}\,\big(\lat_i-\pmean(x)\big),
    \label{eq:util-mucond}\\[2pt]
  \sigma^{2}_{\mathrm{cond}}(x^{*}) &= \pvar(x^{*}) - \frac{\Sigma_{x^{*}}^{2}}{\pvar(x)} .
    \label{eq:util-sigcond}
\end{align}
%
These are transformed to log-rate space, $\gain\,\mu_{\mathrm{cond}}+\latz$ and
$\gain^{2}\sigma^{2}_{\mathrm{cond}}$, and fed to the same entropy
functional~\eqref{eq:Hmarg} to give the per-input conditional entropy
$H(R\mid x^{*},\lat(x),\mathcal D)$. The full Gaussian-conditioning derivation
(augmented inducing matrix, block inverse, and the equivalence to the rank-1
posterior update) is the deep-layer result~\eqref{eq:condmoments}; and the
conditional mean~\eqref{eq:util-mucond} obtained by conditioning on the joint
posterior coincides with the mean obtained by augmenting the inducing set with the
hypothetical observation and re-integrating (Thm.~\ref{thm:predmean-equiv}). Neither
is re-derived here.

\begin{remark}[The cross-covariance pitfall: $\ccond$ must be included]\label{rem:cond-correct}
Forming the joint over $[x;x^{*}]$ and reading the \emph{full} posterior covariance
between the two points already yields $\ccond+\pv\transpose\vV\pv^{*}=\Sigma_{x^{*}}$
--- the $\ccond$ term of~\eqref{eq:util-ccond} included. The pitfall is to instead
assemble the $\pv\transpose\vV\pv^{*}$-only cross-covariance by hand, dropping
$\ccond$: this is biased outside the inducing hull. Taking
$\mu_{\mathrm{cond}},\sigma^{2}_{\mathrm{cond}}$ from the blocks of the complete
joint covariance (with a floor $\sigma^{2}_{\mathrm{cond}}\ge10^{-8}$) avoids it.
\end{remark}

\paragraph{Monte-Carlo estimator (one $\lat$ per input).}
Given pre-drawn samples $\{x_i\}_{i=1}^{N_{\mathrm{mc}}}\sim\px$:
%
\begin{enumerate}\itemsep2pt
  \item $\Hmarg(x^{*})=$ entropy~\eqref{eq:Hmarg} at the marginal moments
        $\mug,\varg$ of $x^{*}$.
  \item For $i=1,\dots,N_{\mathrm{mc}}$: read $(\pmean(x_i),\pvar(x_i))$; draw
        $\lat_i=\pmean(x_i)+\sqrt{\pvar(x_i)}\,\varepsilon,\ \varepsilon\sim\Gauss(0,1)$
        (sampled variant; the deterministic variant sets $\lat_i=\pmean(x_i)$); form
        the conditional moments \eqref{eq:util-mucond}--\eqref{eq:util-sigcond} at
        $x^{*}$; set $\Hcond^{(i)}=$ entropy at $(\gain\mu_{\mathrm{cond}}+\latz,\ \gain^{2}\sigma^{2}_{\mathrm{cond}})$.
  \item $\Hcond=\tfrac{1}{N_{\mathrm{mc}}}\sum_i\Hcond^{(i)}$;\quad
        $\Uda=\Hmarg-\Hcond$.
\end{enumerate}
%
The input posteriors and the reparameterized draws $\lat_i$ do not depend on
$x^{*}$, so they are held fixed (detached from the gradient) when differentiating in
$x^{*}$, avoiding $O(N_{\mathrm{mc}})$ graph growth; the conditional moments at
$x^{*}$ remain differentiable in $x^{*}$. The two variants: the deterministic one
takes $\lat_i=\pmean(x_i)$ (a reproducible sanity check ---
\S\ref{subsubsec:det}); the sampled one is the correct unbiased estimator
(\S\ref{subsec:why-sample}).

\subsubsection{The deterministic special case}
\label{subsubsec:det}

Setting $\lat_i=\pmean(x)$ (the posterior mean) makes the innovation
$\lat_i-\pmean(x)=0$, so~\eqref{eq:util-mucond}--\eqref{eq:util-sigcond} collapse
to $\mu_{\mathrm{cond}}=\pmean(x^{*})$ (mean unchanged) and
%
\begin{equation}
  \sigma^{2}_{\mathrm{cond}}(x^{*}) \;=\; \pvar(x^{*})\,(1-\corr^{2}),
  \qquad
  \corr^{2} \;=\; \frac{\Sigma_{x^{*}}^{2}}{\pvar(x)\,\pvar(x^{*})},
  \label{eq:varreduce}
\end{equation}
%
where $\corr$ is the posterior latent correlation between $\lat(x)$ and
$\lat(x^{*})$. In log-rate coordinates the conditioning map is
$(\mug,\varg)\mapsto(\mug,\varg(1-\corr^{2}))$: same mean, variance scaled by
$1-\corr^{2}$. Writing $\Hmarg(\cdot,\cdot)$ for the count
entropy~\eqref{eq:Hmarg} as a function of the log-rate moments, the utility is the
entropy difference
%
\begin{equation}
  \Uda(x^{*}) \;=\; \Hmarg(\mug,\varg) \;-\; \Hmarg\!\big(\mug,\ \varg(1-\corr^{2})\big).
  \label{eq:da-deterministic}
\end{equation}
%
So the deterministic $\Uda$ depends on exactly three quantities: the operating
point $\mug$, the marginal log-rate variance $\varg$ (how much entropy to start
with), and $\corr^{2}$ (the fraction of that variance conditioning removes).
$\corr^{2}$ is governed by the angle between $x^{*}$ and the conditioning input in
$\Cmat$-space: a small angle gives $\corr^{2}\!\to\!1$ (near-total variance
removal), a large angle $\corr^{2}\!\to\!0$ (no conditioning). This same $\corr^{2}$
sets the sampling bias of \S\ref{subsec:why-sample}.

% ---------------------------------------------------------------------------
\subsection{Why we sample the latent instead of using its mean}
\label{subsec:why-sample}

$\Hcond$ is a double expectation (Remark~\ref{rem:double-exp}), over $x\sim\px$ and
over $\lat(x)\sim q$. The tempting shortcut replaces the inner expectation by
evaluation at the mean,
%
\begin{equation}
  \widetilde{H}_{\mathrm{cond}}(x^{*})
  \;=\; \E_{x\sim\px}\!\big[\,H\big(R\mid x^{*},\lat(x)=\pmean(x),\mathcal D\big)\big].
  \label{eq:mean-shortcut}
\end{equation}

\begin{proposition}[The mean shortcut is biased]\label{prop:bias}
$\E_{\lat(x)}[H(R\mid x^{*},\lat(x),\mathcal D)]\neq H(R\mid x^{*},\lat(x)=\pmean(x),\mathcal D)$
in general. The predictive mean of $\lat(x^{*})$ is linear in the conditioning
value, $\E[\lat(x^{*})\mid\lat(x),\mathcal D]=\pmean(x^{*})+\alpha\,(\lat(x)-\pmean(x))$
with regression coefficient $\alpha=\mathrm{Cov}(\lat(x^{*}),\lat(x)\mid\mathcal D)/\pvar(x)$,
while the predictive \emph{variance} does not depend on the value of $\lat(x)$.
Entropy is nonlinear in the predictive mean, so a second-order Taylor expansion of
$H$ about $\lat(x)=\pmean(x)$, with $\E[\lat(x)-\pmean(x)]=0$, leaves the leading
bias
%
\begin{equation}
  \E_{\lat(x)}[H] - H\big|_{\lat(x)=\pmean(x)}
  \;\approx\; \tfrac12\,\alpha^{2}\,\pvar(x)\;\frac{\partial^{2}\Hmarg}{\partial\pmean^{2}}\Big|_{x^{*}} .
  \label{eq:bias-raw}
\end{equation}
%
Since $\mathrm{Cov}(\lat(x^{*}),\lat(x))=\corr\,\sqrt{\pvar(x^{*})\pvar(x)}$, we
have $\alpha^{2}\,\pvar(x)=\corr^{2}\,\pvar(x^{*})$, and the $\pvar(x)$ cancels:
%
\begin{equation}
  \mathrm{Bias} \;\approx\; \frac{\corr^{2}}{2}\;\pvar(x^{*})\;\frac{\partial^{2}\Hmarg}{\partial\pmean^{2}}\Big|_{x^{*}} .
  \label{eq:bias}
\end{equation}
\end{proposition}

The bias vanishes \emph{only} when $\corr=0$ ($x$ uninformative about $x^{*}$ ---
but then the utility contribution is $\approx0$ anyway), or $\pvar(x^{*})=0$
(nothing left to learn), or $\partial^{2}\Hmarg/\partial\pmean^{2}=0$ (entropy
linear in the mean --- false for the Poisson-lognormal). Notably $\pvar(x)=0$ is
\emph{not} a zero-bias condition: it cancels in~\eqref{eq:bias}.

\begin{remark}[The bias is largest exactly where it matters]\label{rem:bias-matters}
Bias $\propto\corr^{2}\,\pvar(x^{*})$ is maximal when the sampled input $x$ is most
correlated with $x^{*}$ (most relevant to the utility) \emph{and} when $x^{*}$ is
most uncertain (where the acquisition should be steering us) --- i.e.\ largest as
$\corr\to1$. A systematic bias does not shrink with more Monte-Carlo samples: it
converges to the \emph{wrong} value and corrupts the optimizer's direction.
Sampling $\lat$ is unbiased and its variance falls as $O(1/N)$. Hence: sample
$\lat$.
\end{remark}

\paragraph{One $\lat$ per input is optimal (nested MC).}
With $g(x,\lat)=H(R\mid x^{*},\lat,\mathcal D)$, let $\tau^{2}=\mathrm{Var}_x[\E_{\lat}g]$
(between-input) and $\E_x[\sigma^{2}_x]$ (within-input). For a fixed budget of $M$
entropy evaluations, using $N_{\lat}$ latent samples per input gives variance
$(N_{\lat}\tau^{2}+\E_x[\sigma^{2}_x])/M$, strictly worse for $N_{\lat}>1$.
Spending the budget on \emph{more inputs}, one $\lat$ each, is optimal.

\caveat{The magnitude of this bias is a provisional estimate. Pushed to log-rate
space the curvature carries a gain factor,
$\partial^{2}\Hmarg/\partial\pmean^{2}=\gain^{2}\,\partial^{2}\Hmarg/\partial\lgr^{2}$,
so $\mathrm{Bias}\propto\gain^{2}$; for a small gain (e.g.\ $\gain\approx0.03$) the
bias is far below the MC standard error, and the deterministic and sampled $\Uda$
nearly coincide. The regime guide $\gain\ll1$ negligible / $\gain\sim1$ moderate /
$\gain\gg1$ essential, and the specific numbers, are plausibility arguments, not
settled results. The core claim --- bias~\eqref{eq:bias} does not vanish with $N$
--- is solid; the magnitude estimate is provisional. The \emph{correct} estimator
remains the sampled one regardless.}

% ---------------------------------------------------------------------------
\subsection{The entropy landscape over input space}
\label{subsec:landscape}

\paragraph{Monotone shape --- utility is not pure epistemic uncertainty.}
The count entropy $\Hmarg(\mug,\varg)$ increases monotonically in \emph{both}
arguments: with $\mug$ (higher predicted rate) and with $\varg$ (more
GP/epistemic uncertainty). A candidate therefore scores high for a \emph{mixture}
of reasons --- it is predicted to have a high rate \emph{and} it is epistemically
uncertain. The utility is thus not a pure epistemic-uncertainty acquisition;
predicted rate is baked in. This is a deliberate modelling choice: a normalized
arc-cosine kernel would force $\Kf(x,x)=1$, removing the norm channel so that
utility tracked only angular alignment (pure epistemic), but the held-out
predictive correlation then drops by $\sim$25\% (a fitted example, $\Mind=100$:
$0.79\!\to\!0.59$) --- input norm is genuinely informative, so the kernel keeps the
norm dependence.

\paragraph{The $\rmax$ truncation artifact.}
With a fixed $\rmax=100$ the Laplace sum~\eqref{eq:Hmarg} misses the probability
mass above $\rmax$ once the Poisson peak climbs past it, and $\Hmarg$
\emph{collapses to $\approx0$} for $\mug\gtrsim\log100\approx4.6$ (and at lower
$\mug$ for large $\varg$). This is a numerical artifact, not real entropy: the
computable region is roughly triangular in $(\mug,\varg)$. A cheap pre-check
without running the Laplace sum is $z_{\mathrm{safe}}=(\log\rmax-\mug)/\sqrt{\varg}$
--- $z_{\mathrm{safe}}>2$ safe, $<1.5$ unreliable, $<0.5$ broken. Inputs from the
pool sit at $z_{\mathrm{safe}}\gg10$, always safe; trouble arises only for
artificially amplified inputs. An adaptive choice of $\rmax$, sized from the moments
so the truncation always covers the Poisson mass, removes the collapse.

\paragraph{Norm-exploitation divergence.}
Because the arc-cosine kernel is positively $1$-homogeneous
($\Kf(\scal x,y)=\scal\,\Kf(x,y)$ when the bias $\sigz=0$), scaling
$x^{*}\to\scal x^{*}$ scales $\pmean(x^{*})\sim\scal$ and $\pvar(x^{*})\sim\scal^{2}$,
so $\Hmarg$ grows while a directionally aligned conditioning input keeps
$\corr\approx1$ and $\sigma^{2}_{\mathrm{cond}}\approx0$. Under free-space gradient
ascent $\Uda$ then grows $\sim\scal^{1.9}$ (numerically) --- the optimizer
``diverges'' by amplifying input norm rather than finding angularly informative
inputs. This is intrinsic to the homogeneous kernel, \emph{not} a bug; the full
statement and the $\sigz^{2}>0$ damping (alignment peak strictly at $\scal=1$) are
Thm.~\ref{thm:divergence} in the deep layer. Over a finite \emph{pool} of candidate
inputs (discrete selection) norms are bounded and this never triggers.

\caveat{\textbf{Utility accuracy ceiling (known limitation).} At high rate the two
utilities fail in opposite directions through the same $\rmax=100$ truncation:
$\Ustd$ can run to $+\infty$ while $\Uda$ silently collapses toward $0$. A rate
guard ($f_{\max}$, default $100$) keeps the operating point in the accurate regime;
a dedicated look-up table for the $\Uda$ path was not built. This is a limitation of
the fixed-$\rmax$ path, cured on the standard path by the adaptive choice of
$\rmax$.}

% ============================================================================
%  sections/part3b_deep_layer.tex  --  Part III, deep layer (proofs / subspace /
%  divergence / kernel fixes / numerics).  \input-ready fragment; no preamble.
%  Uses ONLY the macros defined in preamble.tex.  Domain-neutral (math/ML only).
% ============================================================================

\section{The deep layer: conditioning, subspace, divergence, and numerics}
\label{sec:deeplayer}

Part~III (\S7) \emph{defined} the two acquisition utilities and their entropies;
this section proves \emph{why} those formulas behave as they do. It supplies the
theorems underlying the moment formulas, the exact predictive-conditioning
derivation, the subspace-validity argument, the norm-scaling divergence theorems
that explain the utility's blow-up under unconstrained input optimization, the
two kernel-level fixes, and the four-failure-mode numerical audit of the standard
utility.

Two objects recur, and their ``subtracted term'' is \emph{not} the same (the
single most load-bearing clash; the top-level exposition is in Part~III):
\begin{itemize}[leftmargin=1.4em]
  \item the \textbf{standard} utility
        $\Ustd(x^*)=\Hmarg(x^*)-\E_{\lgr}[\Hnoise(x^*)]$ of \eqref{eq:Ustd},
        which the selection loop maximizes (\S\ref{subsec:numerics} audits
        it); its subtracted term is the expected Poisson \emph{noise} entropy
        $\E_{\lgr}[\Hnoise]$, an average over the point's \emph{own} latent, no
        second input involved;
  \item the \textbf{distribution-aware (DA)} utility
        $\Uda(x^*)=\Hmarg(x^*)-\E_{x\sim\px}[\Hcond(x^*)]$ of \eqref{eq:Uda},
        whose subtracted term is the entropy \emph{after Gaussian-conditioning}
        the GP on a hypothetical observation of $\lat(x_t)$ at a target input
        $x_t$.
\end{itemize}
The two coincide only in the $\corr^2\!=\!1$ deterministic limit, where $\Hcond$
collapses to a single Poisson entropy (\S\ref{subsec:divergence}). The divergence
theorems of \S\ref{subsec:divergence} and the remedies of
\S\ref{subsec:kernelfixes} concern \emph{gradient ascent on the input $x^*$}
(a free input-synthesis setting), which is distinct from selection over a fixed
pool.

% ===========================================================================
\subsection{GP moments and Gaussian conditioning (Proof III)}
\label{subsec:moments}

This subsection derives the marginal and conditional GP moments underlying the
utilities of Part~III.

\subsubsection{The self-kernel identity}
Recall the order-1 arc-cosine kernel of Part~I, \eqref{eq:acoskernel},
\[
  \Kf(x,x') \;=\; \tfrac1\pi\,\Kmag\,\Jang(\kang),\qquad
  \Kmag=\sqrt{\vmag_{x}\,\vmag_{x'}},\quad
  \cos\kang=\frac{x\transpose\Cmat x'+\sigz^2}{\Kmag},\quad
  \Jang(\kang)=\sin\kang+(\pi-\kang)\cos\kang,
\]
with self-magnitude $\vmag_{x}=x\transpose\Cmat x+\sigz^2$.

\begin{lemma}[Self-kernel / prior variance]\label{lem:selfkernel}
For every $x$, $\Kf(x,x)=\lVert x\rVert_{\Cmat}^2+\sigz^2=\vmag_{x}$, where
$\lVert x\rVert_{\Cmat}^2\equiv x\transpose\Cmat x$.
\end{lemma}
\begin{proof}
Set $x'=x$: $\cos\kang=\vmag_x/\vmag_x=1\Rightarrow\kang=0$, and
$\Jang(0)=\sin0+\pi\cos0=\pi$. Hence
$\Kf(x,x)=\tfrac1\pi\,\vmag_x\,\pi=\vmag_x=x\transpose\Cmat x+\sigz^2$.
\end{proof}

This restates \eqref{eq:selfkernel}: \emph{the prior variance at any point equals
its squared $\Cmat$-norm (plus the bias variance).} It is the single identity
from which the entire norm-scaling divergence of \S\ref{subsec:divergence} grows.

\subsubsection{Variational posterior moments}
With $\Mind$ inducing points $\{\ipt_j\}$ and variational posterior
$q(\latt)=\Gauss(\vm,\vV)$, write the inducing gram $\Ktil=\Kf(\Ipt,\Ipt)$, the
cross-kernel vector $\kvec(x)=[\Kf(x,\ipt_1),\dots,\Kf(x,\ipt_\Mind)]\transpose$,
the projection $\pv(x)=\Ktil^{-1}\kvec(x)$, and the (training-fixed) correction
matrix
\[
  W \;=\; \Ktil^{-1}(\vV-\Ktil)\,\Ktil^{-1}.
\]
The marginal moments are
\begin{equation}\label{eq:margmoments}
\begin{aligned}
  \pmean(x^*) &= \kvec(x^*)\transpose\Ktil^{-1}\vm,\\
  \pvar(x^*)  &= \Kf(x^*,x^*)+\kvec(x^*)\transpose W\,\kvec(x^*)\\
              &= \lVert x^*\rVert_{\Cmat}^2+\sigz^2+\kvec(x^*)\transpose W\,\kvec(x^*).
\end{aligned}
\end{equation}
The first term of $\pvar$ is the prior variance (Lemma~\ref{lem:selfkernel}); the
second is the posterior correction (negative when $\vV\prec\Ktil$, i.e.\ data
shrink variance). Between two points,
\[
  \Sigma_q(x^*,x_t)=\Kf(x^*,x_t)+\kvec(x^*)\transpose W\,\kvec(x_t),
  \qquad \Sigma_q(x^*,x^*)=\pvar(x^*),\ \ \Sigma_q(x_t,x_t)=\pvar(x_t),
\]
and the squared posterior correlation --- the Pearson correlation of the joint
Gaussian $[\lat(x^*),\lat(x_t)]$ under $q$ --- is
\begin{equation}\label{eq:corrdef}
  \corr^2(x^*,x_t)=\frac{\Sigma_q(x^*,x_t)^2}{\pvar(x^*)\,\pvar(x_t)}\in[0,1].
\end{equation}

\subsubsection{Conditioning on an observed latent}
Observing $\lat_t$ at $x_t$ and applying the standard Gaussian conditioning
identity to $[\lat(x^*),\lat(x_t)]$ gives the conditional moments:
\begin{equation}\label{eq:condmoments}
  \boxed{\;
  \begin{aligned}
  \pmean_{\mathrm{cond}}(x^*)&=\pmean(x^*)+\frac{\Sigma_q(x^*,x_t)}{\pvar(x_t)}\bigl(\lat_t-\pmean(x_t)\bigr),\\[2pt]
  \pvar_{\mathrm{cond}}(x^*)&=\pvar(x^*)-\frac{\Sigma_q(x^*,x_t)^2}{\pvar(x_t)}=\pvar(x^*)\bigl(1-\corr^2\bigr).
  \end{aligned}
  \;}
\end{equation}

\begin{proposition}[$\corr^2$ is the fractional variance reduction]\label{prop:fracvar}
The conditional variance at $x^*$ is the marginal variance scaled by
$(1-\corr^2)$; equivalently, conditioning on $\lat(x_t)$ removes a fraction
$\corr^2$ of the variance at $x^*$.
\end{proposition}
\begin{proof}
Substitute the correlation definition \eqref{eq:corrdef} into the variance line
of \eqref{eq:condmoments}:
\[
  \pvar(x^*)-\frac{\Sigma_q(x^*,x_t)^2}{\pvar(x_t)}
  =\pvar(x^*)\Bigl[1-\frac{\Sigma_q(x^*,x_t)^2}{\pvar(x^*)\,\pvar(x_t)}\Bigr]
  =\pvar(x^*)(1-\corr^2).
\]
\end{proof}

The conditional moments are obtained by conditioning on the \emph{full joint}
covariance of $[x_t;x^*]$ and applying exactly \eqref{eq:condmoments} (with a
floor $\pvar_{\mathrm{cond}}\ge10^{-8}$). By the theorem of
\S\ref{subsec:predictive} this joint route is the \emph{exact} augmented
predictive, not the sparse approximation.

\subsubsection{Log-rate transform and the DA utility in three scalars}
Under $\lgr=\gain\lat+\latz$ and $R\sim\Poi(e^{\lgr})$, the GP moments push to
$\lgr$-space as $\mug=\gain\pmean+\latz$, $\varg=\gain^2\pvar$ (marginal) and
$\mu_{g,\mathrm{cond}}=\gain\,\pmean_{\mathrm{cond}}+\latz$,
$\sigma^{2}_{g,\mathrm{cond}}=\gain^2\pvar_{\mathrm{cond}}=\varg(1-\corr^2)$
(the variance reduction carries through). Hence, writing
$\Hmarg(\mug,\varg)$ for the marginal Poisson--log-normal entropy \eqref{eq:Hmarg},
\[
  \Uda(x^*)=\Hmarg(\mug,\varg)-\Hmarg\!\left(\mu_{g,\mathrm{cond}},\,\sigma^{2}_{g,\mathrm{cond}}\right).
\]
In deterministic mode ($\lat_t=\pmean(x_t)$) the conditional mean is unchanged,
$\mu_{g,\mathrm{cond}}=\mug$, and \eqref{eq:Uda} reduces to the deterministic
distribution-aware utility \eqref{eq:da-deterministic} of \S\ref{subsubsec:det},
a function of exactly three scalars:
\begin{equation}\label{eq:daboxed}
  \boxed{\;\Uda(x^*)=\Hmarg(\mug,\varg)-\Hmarg\!\left(\mug,\ \varg(1-\corr^2)\right)\;}
\end{equation}
the operating point $\mug=\gain\pmean(x^*)+\latz$, the marginal variance
$\varg=\gain^2\pvar(x^*)$, and the fractional reduction $\corr^2$. Conditioning
slides $(\mug,\varg)\!\to\!(\mug,\varg(1-\corr^2))$: same mean, reduced variance;
the utility is the entropy gap. This is the object whose divergence
\S\ref{subsec:divergence} analyses.

% ===========================================================================
\subsection{The predictive distribution conditioned on a new observation}
\label{subsec:predictive}

The DA utility integrates the entropy of $p(\lat(x^*)\mid\lat(x),\mathcal{D})$:
the predictive over the latent at a query $x^*$ after a \emph{hypothetical} new
latent observation at $x$. We derive it two ways --- sparse-approximated and
exact/augmented --- and prove the two agree in mean.

Define the \emph{prior conditional covariance} (the residual not captured by the
inducing points)
\[
  \ccond(x,x') \;\equiv\; \Kf(x,x')-\pv(x)\transpose\Ktil\,\pv(x')
                        \;=\; \Kf(x,x')-\kvec(x)\transpose\Ktil^{-1}\kvec(x'),
\]
whose diagonal is the residual / aleatoric variance
$\varr(x)\equiv\ccond(x,x)=\Kf(x,x)-\kvec(x)\transpose\Ktil^{-1}\kvec(x)\ge0$.

\subsubsection{The rank-1 posterior update}
\begin{lemma}[Rank-1 update]\label{lem:rank1}
Let the new point $x$ have projection $\pv=\pv(x)$ and residual variance
$\varr=\varr(x)$, so the conditional likelihood is
$p(\lat(x)\mid\latt)=\Gauss(\pv\transpose\latt,\varr)$. Conditioning the
variational posterior $\Gauss(\vm,\vV)$ on the observed $\lat(x)$ yields
$\Gauss(\vm',\vV')$ with
\begin{equation}\label{eq:rank1}
  \vm' \;=\; \vm+\frac{\vV\pv}{\varr+\pv\transpose\vV\pv}\bigl(\lat(x)-\pv\transpose\vm\bigr),
  \qquad
  \vV' \;=\; \vV-\frac{\vV\pv\,\pv\transpose\vV}{\varr+\pv\transpose\vV\pv}.
\end{equation}
\end{lemma}
\begin{proof}
Under the variational approximation the joint of $(\lat(x),\latt)$ is Gaussian:
\[
  \begin{bmatrix}\lat(x)\\ \latt\end{bmatrix}
  \sim\Gauss\!\left(
  \begin{bmatrix}\pv\transpose\vm\\ \vm\end{bmatrix},\
  \begin{bmatrix}\varr+\pv\transpose\vV\pv & \pv\transpose\vV\\[2pt]
                 \vV\pv & \vV\end{bmatrix}\right).
\]
The Gaussian conditioning identity on $\latt$ given $\lat(x)$, with scalar
lower-right-to-upper-left Schur term $\Sigma_{11}=\varr+\pv\transpose\vV\pv$ and
cross-block $\vV\pv$, gives \eqref{eq:rank1} directly.
\end{proof}

\subsubsection{Sparse vs.\ augmented predictive}
Let $\pv^*=\Ktil^{-1}\kvec(x^*)$ and $\varr(x^*)=\Kf(x^*,x^*)-\kvec(x^*)\transpose\Ktil^{-1}\kvec(x^*)$.

\emph{Sparse route.} Assuming conditional independence $p(\lat(x^*)\mid\lat(x),\latt)\approx p(\lat(x^*)\mid\latt)$ and integrating the conditional prior against the updated posterior of Lemma~\ref{lem:rank1},
\[
  \pmean_{\mathrm{pred}}={\pv^*}\transpose\vm',\qquad
  \pvar_{\mathrm{pred}}=\varr(x^*)+{\pv^*}\transpose\vV'\pv^*
                       =\Kf(x^*,x^*)-{\pv^*}\transpose(\Ktil-\vV')\pv^*.
\]

\emph{Augmented (exact) route.} Keep the direct $\lat(x)$--$\lat(x^*)$
correlation. With augmented latent $\latt_{\mathrm{aug}}=[\lat(x);\latt]$ and
augmented gram / cross-kernel
\[
  \widetilde{K}_{x}=\begin{bmatrix}\Kf(x,x) & \kvec(x)\transpose\\ \kvec(x) & \Ktil\end{bmatrix},
  \qquad
  \kvec_{\mathrm{aug}}(x^*)=\begin{bmatrix}\Kf(x,x^*)\\ \kvec(x^*)\end{bmatrix},
\]
the non-approximated conditional prior is
$p(\lat(x^*)\mid\latt_{\mathrm{aug}})=\Gauss(\pv_{\mathrm{aug}}\transpose\latt_{\mathrm{aug}},S^*)$,
$\pv_{\mathrm{aug}}=\widetilde{K}_x^{-1}\kvec_{\mathrm{aug}}(x^*)$,
$S^*=\Kf(x^*,x^*)-\pv_{\mathrm{aug}}\transpose\widetilde{K}_x\,\pv_{\mathrm{aug}}$.
Since $\lat(x)$ is \emph{observed} (variance $0$) while $\latt$ follows the
updated posterior $\Gauss(\vm',\vV')$, the augmented posterior covariance is
$V''=\bigl[\begin{smallmatrix}0&0\transpose\\0&\vV'\end{smallmatrix}\bigr]$, and
\[
  \pmean_{\mathrm{pred}}=\pv_{\mathrm{aug}}\transpose[\lat(x);\vm'],\qquad
  \pvar_{\mathrm{pred}}=S^*+\pv_{\mathrm{aug}}\transpose V''\pv_{\mathrm{aug}}
                       =\Kf(x^*,x^*)-\pv_{\mathrm{aug}}\transpose(\widetilde{K}_x-V'')\,\pv_{\mathrm{aug}}.
\]

\begin{remark}[Augmented route fixes the sparse self-variance artifact]
When $x=x^*$, $\pv_{\mathrm{aug}}$ aligns with the first column of
$\widetilde{K}_x$ and, because the $(1,1)$ block of $V''$ is $0$, the predictive
variance is forced to $0$. A self-consistent predictive \emph{must} have zero
variance at an already-observed point; the sparse route leaves a spurious nonzero
self-variance there. This is the same pathology as a separately-assembled
cross-covariance that can violate Cauchy--Schwarz (Part~II \S4.7); conditioning
on the full joint covariance sidesteps both.
\end{remark}

\subsubsection{Equivalence of predictive means}
\begin{theorem}[Equivalence of predictive means]\label{thm:predmean-equiv}
The augmented-integration predictive mean equals the mean obtained by standard
Gaussian conditioning on the joint predictive $[\lat(x),\lat(x^*)]$ --- exactly,
with no cross-covariance approximation.
\end{theorem}
\begin{proof}
\textbf{Method 1 (joint conditioning).} The joint moments given $\mathcal{D}$ are
$\pmean_x=\pv\transpose\vm$, $\pmean_{x^*}={\pv^*}\transpose\vm$,
\[
  \Sigma_{11}=\varr+\pv\transpose\vV\pv,\qquad
  \Sigma_{12}=\underbrace{\bigl(\Kf(x,x^*)-\pv\transpose\Ktil\,\pv^*\bigr)}_{=\ \ccond}+{\pv^*}\transpose\vV\pv,
\]
where the residual cross term is exactly $\ccond=\ccond(x,x^*)$. The Gaussian
conditioning identity gives the target
\begin{equation}\label{eq:method1}
  \E[\lat(x^*)\mid\lat(x),\mathcal{D}]
  =\pmean_{x^*}+\frac{\Sigma_{12}}{\Sigma_{11}}\bigl(\lat(x)-\pmean_x\bigr)
  ={\pv^*}\transpose\vm+\frac{\Sigma_{12}}{\Sigma_{11}}\bigl(\lat(x)-\pv\transpose\vm\bigr).
\end{equation}

\textbf{Method 2 (augmented integration).} Block-invert $\widetilde{K}_x$. Writing
the residual-correlation coefficient $\vartheta\equiv\ccond/\varr$ (denoted
$\vartheta$ rather than $\alpha$ to avoid clashing with the localization weight
$\loc_i$),
\[
  \pv_{\mathrm{aug}}=\begin{bmatrix}\vartheta\\ \pv^*-\vartheta\,\pv\end{bmatrix},
  \qquad \vartheta=\frac{\Kf(x,x^*)-\pv\transpose\Ktil\,\pv^*}{\varr}=\frac{\ccond}{\varr}.
\]
Then, with $\vm'$ from \eqref{eq:rank1},
\[
  \pmean_{\mathrm{aug}}=\pv_{\mathrm{aug}}\transpose[\lat(x);\vm']
  =\vartheta\,\lat(x)+({\pv^*}\transpose-\vartheta\,\pv\transpose)\vm'
  =\vartheta\,\lat(x)+({\pv^*}\transpose-\vartheta\,\pv\transpose)\!\left[\vm+\frac{\vV\pv}{\Sigma_{11}}\bigl(\lat(x)-\pv\transpose\vm\bigr)\right].
\]
Group by the prior mean ${\pv^*}\transpose\vm$ and the innovation
$(\lat(x)-\pv\transpose\vm)$, adding and subtracting $\vartheta\,\pv\transpose\vm$:
\[
  \pmean_{\mathrm{aug}}={\pv^*}\transpose\vm
   +\underbrace{\left[\vartheta+\frac{{\pv^*}\transpose\vV\pv-\vartheta\,\pv\transpose\vV\pv}{\Sigma_{11}}\right]}_{\displaystyle \mathcal{C}}
   \bigl(\lat(x)-\pv\transpose\vm\bigr).
\]
Put $\mathcal{C}$ over the common denominator $\Sigma_{11}=\varr+\pv\transpose\vV\pv$
and use $\vartheta\,\varr=\ccond$:
\[
  \mathcal{C}=\frac{\vartheta\,\varr+\vartheta\,\pv\transpose\vV\pv+{\pv^*}\transpose\vV\pv-\vartheta\,\pv\transpose\vV\pv}{\Sigma_{11}}
   =\frac{\vartheta\,\varr+{\pv^*}\transpose\vV\pv}{\Sigma_{11}}
   =\frac{\ccond+{\pv^*}\transpose\vV\pv}{\Sigma_{11}}
   =\frac{\Sigma_{12}}{\Sigma_{11}}.
\]
Hence $\pmean_{\mathrm{aug}}={\pv^*}\transpose\vm+(\Sigma_{12}/\Sigma_{11})(\lat(x)-\pv\transpose\vm)$,
identical to \eqref{eq:method1}.
\end{proof}

\begin{remark}
Theorem~\ref{thm:predmean-equiv} certifies that conditioning on the full joint
covariance (\S\ref{subsec:moments}) computes the \emph{exact} augmented
predictive mean --- it does not incur the sparse cross-covariance error.
Setting the residual cross term $\ccond=0$ would enforce the sparse factorization
\[
  p(\lat^*,\lat\mid\latt)\approx p(\lat^*\mid\latt)\,p(\lat\mid\latt);
\]
the inducing (epistemic) cross-covariance ${\pv^*}\transpose\vV\pv$ can still be
nonzero, because changing belief about an inducing point moves both $\lat(x^*)$
and $\lat(x)$.
\end{remark}

% ===========================================================================
\subsection{Subspace theory}
\label{subsec:subspace}

\emph{Why} the utility may be optimized in a low-dimensional subspace of the
input space ($d=11664$ grid points in full; $\sim$900--2356 in the
localized-support region) instead of the full input space. Setting: optimize
$x^*$ to maximize $\Uda(x^*)$ by gradient ascent, constrained to
$x_{\mathrm{loc}}=\bar{x}+\text{(basis)}\cdot\subc$ over the subspace coordinate
$\subc$; the offset $\bar{x}$ is the dataset mean over the localized-support
coordinates, so $\subc=0$ maps to the mean input.

\subsubsection{The kernel touches $x$ only through the $\Cmat$-eigenspace}
Let $\Cmat=\sum_i\gamma_i\,e_i e_i\transpose$ be the eigendecomposition of the
structured PSD matrix of Part~I (\S3.3), with $\gamma_1\ge\dots\ge\gamma_d\ge0$
and orthonormal $\Cmat$-eigenvectors $e_i$.

\begin{proposition}[$\Cmat$-eigenspace validity]\label{prop:subspace}
Every kernel evaluation depends on $x$ only through the quadratic forms
$x\transpose\Cmat x=\sum_i\gamma_i(e_i\transpose x)^2$ and $x\transpose\Cmat x'$;
and $\nabla_x\Kf$ lies in the column space of $\Cmat$, with its component along
$e_i$ scaled by $\gamma_i$. Consequently, gradient ascent already concentrates
movement along the top $\Cmat$-eigenvectors, and truncating to the leading
$\Cmat$-eigen-subspace is (at full rank) an \emph{exact reparameterization} and
(at reduced rank) drops only near-zero-gradient directions.
\end{proposition}
\begin{proof}
The kernel enters $x$ solely via $x\transpose\Cmat x+\sigz^2=\vmag_x$ and
$x\transpose\Cmat x'+\sigz^2$ (through $\Kmag$ and $\cos\kang$); in the
$\Cmat$-eigenbasis $x\transpose\Cmat x=\sum_i\gamma_i(e_i\transpose x)^2$, so a
direction $e_i$ with $\gamma_i\approx0$ contributes $\approx0$ to \emph{every}
kernel value. The input gradient of the arc-cosine kernel (Part~I \S3.5,
\eqref{eq:gradxK}) is
\[
  \nabla_x\Kf(x,x')=\tfrac1\pi\Bigl[(\pi-\kang)\,\Cmat x'+\sin\kang\,\sqrt{\vmag_{x'}/\vmag_x}\;\Cmat x\Bigr],
\]
every term of which is $\Cmat\cdot(\text{vector})$; projecting onto $e_i$ multiplies
by $\gamma_i$. Thus movement along small-$\gamma_i$ directions is intrinsically
suppressed. Removing those directions (a) does not change the optimum (they carried
near-zero gradient) and (b) removes flat directions, giving a better-conditioned,
lower-effective-dimension problem.
\end{proof}

So the $\Cmat$-eigenspace projection is a \emph{convergence aid / exact
reparameterization}, not a constraint on the input distribution (analogous to the
eigenvalue threshold on $\Ktil$ in Part~II \S5.1).

\subsubsection{Three bases and the Szeg\H{o}/Fourier equivalence}
All three write $x_{\mathrm{loc}}=\bar{x}+\text{(basis)}\cdot\subc$ and share one
optimizer/diagnostics; they differ only in the basis.
\begin{itemize}[leftmargin=1.4em]
  \item \textbf{PCA.} Columns $\Phi_K=[\phi_1,\dots,\phi_K]$ = top-$K$ eigenvectors
        of the data covariance $\Sigma_{\mathrm{data}}=\tfrac1N X_c\transpose X_c$
        (eigenvalues $\nu_1\ge\nu_2\ge\dots$). The offset $\bar{x}$ is
        \emph{mathematically intrinsic} (centering before SVD); truncation is a
        \emph{genuine data-driven constraint} (excludes low-data-variance
        directions). At variance threshold $0.80$: $K=197$ of the $2356$
        localized-support coordinates. Gradient projects as
        $\nabla_\subc\Uda=\Phi_K\transpose\nabla_x\Uda$.
        PCA captures \emph{second-order statistics only} (mean $+$ covariance);
        it is blind to phase structure, sparsity, and non-Gaussian marginals.
  \item \textbf{$\Cmat$-eigenspace.} Columns $E_K=[e_1,\dots,e_K]$ = top-$K$
        $\Cmat$-eigenvectors. The offset is \emph{not} intrinsic (the
        eigendecomposition passes through the origin) --- see the caveat below.
        Truncation is the reparameterization of Proposition~\ref{prop:subspace}.
        The spectrum is extremely skewed (top eigenvalue $\sim44\%$ of the total):
        at relative threshold $10^{-3}$, $K=28$ of $2356$ captures $99.7\%$ of the
        eigenvalue mass. (Float32 noise floor: eigenvalues below $\sim10^{-7}\gamma_1$
        are numerical noise; thresholds below $\sim10^{-6}$ pull in garbage
        eigenvectors.)
  \item \textbf{Combined.} Project $\Cmat$ into the PCA subspace and re-diagonalize:
        $\Cmat_{\mathrm{pca}}=\Phi_K\transpose\Cmat\,\Phi_K=Q\,\Gamma_Q Q\transpose$,
        keep the columns $Q$ above the eigen-threshold, basis $=\Phi_K Q$.
        Directions that are \emph{both} data-typical (PCA) and kernel-visible
        ($\Cmat$-eigen). At $(0.80,10^{-3})$: $2356\to K_{\mathrm{pca}}=197\to
        K_{\mathrm{comb}}=28$, achieving near-identical utility to PCA with
        $\sim7\times$ fewer dimensions.
\end{itemize}

\begin{proposition}[PCA $\approx$ Fourier under stationarity (Szeg\H{o})]
\label{prop:szego}
If the input process is second-order stationary on a regular grid, its covariance
is (block-)Toeplitz, $\Sigma[i,j]=f(|i-j|)$. By Szeg\H{o}'s theorem, as $d\to\infty$
the eigenvectors converge to Fourier modes and the eigenvalues to the power
spectral density $S(\omega)=\sum_k f(k)e^{-i\omega k}$. The input distribution has
$S(\omega)\sim1/|\omega|^2$, so $\Sigma_{\mathrm{data}}$'s eigenvectors approach
the Fourier basis and PCA-truncation approximates keeping the $K$
lowest-frequency Fourier modes.
\end{proposition}
The equivalence is weakened by the irregular masked region, the non-stationarity
the localization weight $\loc$ injects (grid points near the center $\rfc$ have
higher variance, breaking Toeplitz structure), finite $N\approx10^4$ in
$d\approx900$, and residual non-stationarity of the input distribution; at
$\sim30\times30$ grid patches it is reasonable.

\subsubsection{Subspace $\ne$ distributional constraint}
All three are \emph{subspace} constraints, not \emph{distributional} ones. None
guarantees $x^*$ is a plausible sample of $\px$: an input with the right power
spectrum but wrong phase (correlated noise) satisfies all three. A true
distributional constraint needs higher-order structure (phases, sparsity)
--- i.e.\ a generative model, which is the motivation for the guided-diffusion
synthesis of Part~IV. The subspace approach is a useful \emph{diagnostic} (does
restricting to the right power spectrum change utility? does the GP utility
operate primarily at second order?), not a full fix.

\caveat{The $\Cmat$-eigenspace offset is not mathematically motivated. Unlike PCA,
where the mean offset is intrinsic to centering, the $\Cmat$-eigendecomposition
passes through the origin, so setting $\bar{x}=\text{mean}$ has no principled
basis --- it is forced by \emph{limiter precision}: the non-capturable
eigenvectors collapse to $0$ instead of to the mean, so the offset is set to the
mean to compensate. This changes the utility landscape (the kernel sees
$\Cmat(\bar{x}+E_K\subc)$, not $\Cmat(E_K\subc)$) and should be investigated. It
is an acknowledged modeling wrinkle, not a resolved result.}

% ===========================================================================
\subsection{Divergence theorems (Proof I): norm scaling and the utility blow-up}
\label{subsec:divergence}

The arc-cosine kernel is positively homogeneous; scaling an input along a fixed
direction inflates the posterior variance as the square of the scale while
preserving the conditioning benefit, so unconstrained gradient ascent on the DA
utility diverges. It is proved once here and reused by \S\ref{subsec:kernelfixes}.
Throughout, $x^*=\scal\,x_t$ with scale $\scal>0$, and $q\equiv x_t\transpose\Cmat x_t$.

\begin{theorem}[Cross-kernel proportionality, $\sigz^2=0$]\label{thm:divergence}
Let $\sigz^2=0$ and $x^*=\scal\,x_t$, $\scal>0$. Then $\Kf(x^*,z)=\scal\,\Kf(x_t,z)$
for every $z$, and hence $\kvec(x^*)=\scal\,\kvec(x_t)$.
\end{theorem}
\begin{proof}
With $\sigz^2=0$, $\vmag_{x^*}=\scal^2 q$ and $\vmag_{x_t}=q$. For any $z$ with
$\vmag_z=z\transpose\Cmat z$,
\[
  \cos\kang(x^*,z)=\frac{(\scal x_t)\transpose\Cmat z}{\sqrt{\scal^2 q\cdot\vmag_z}}
   =\frac{\scal\,x_t\transpose\Cmat z}{\scal\sqrt{q\,\vmag_z}}
   =\frac{x_t\transpose\Cmat z}{\sqrt{q\,\vmag_z}}=\cos\kang(x_t,z),
\]
so $\kang(x^*,z)=\kang(x_t,z)$ and $\Jang(\kang(x^*,z))=\Jang(\kang(x_t,z))$. Then
$\Kf(x^*,z)=\tfrac1\pi\sqrt{\scal^2 q\,\vmag_z}\,\Jang(\kang)=\scal\cdot\tfrac1\pi\sqrt{q\,\vmag_z}\,\Jang(\kang)=\scal\,\Kf(x_t,z)$.
Applying this to each inducing point $z=\ipt_j$ gives $\kvec(x^*)=\scal\,\kvec(x_t)$.
\end{proof}

\begin{corollary}[Perfect posterior correlation]\label{cor:rho1}
Under Theorem~\ref{thm:divergence}, $\corr^2(x^*,x_t)=1$.
\end{corollary}
\begin{proof}
From $\kvec(x^*)=\scal\,\kvec(x_t)$ and (with $x^*\parallel x_t$, $\kang=0$,
$\Jang(0)=\pi$) $\Kf(x^*,x_t)=\scal\,\Kf(x_t,x_t)=\scal\,\vmag_{x_t}$, substitute
into $\Sigma_q$ with $W=\Ktil^{-1}(\vV-\Ktil)\Ktil^{-1}$:
\[
  \Sigma_q(x^*,x_t)=\scal\,\vmag_{x_t}+\scal\,\kvec(x_t)\transpose W\kvec(x_t)=\scal\,\Sigma_q(x_t,x_t),
  \quad
  \Sigma_q(x^*,x^*)=\scal^2\,\Sigma_q(x_t,x_t),
\]
so $\corr^2=[\scal\,\Sigma_q(x_t,x_t)]^2/[\scal^2\Sigma_q(x_t,x_t)\cdot\Sigma_q(x_t,x_t)]=1$.
\end{proof}

\begin{corollary}[Moment scaling]\label{cor:momentscale}
Under Theorem~\ref{thm:divergence}, $\pmean(\scal x_t)=\scal\,\pmean(x_t)$ and
$\pvar(\scal x_t)=\scal^2\,\pvar(x_t)$.
\end{corollary}
\begin{proof}
Substitute $\kvec(x^*)=\scal\,\kvec(x_t)$ into $\pmean=\kvec\transpose\Ktil^{-1}\vm$
(linear in $\kvec$) and $\pvar=\Kf(x^*,x^*)+\kvec\transpose W\kvec$ (quadratic in
$\kvec$, with $\Kf(x^*,x^*)=\scal^2\vmag_{x_t}$).
\end{proof}

By Corollary~\ref{cor:rho1}, $\corr^2=1\Rightarrow\pvar_{\mathrm{cond}}(x^*)=0$:
conditioning removes \emph{all} the variance. Yet by
Corollary~\ref{cor:momentscale} the marginal moments still grow
($\mug\propto\scal$, $\varg\propto\scal^2\Rightarrow\Hmarg\to\infty$), so $\Uda$
still diverges. In the exact parallel case ($\kang=0,\sigz^2=0$),
$\pmean_{\mathrm{cond}}=\pmean$, $\pvar_{\mathrm{cond}}=0$, and $\Hcond$ collapses
to a single Poisson entropy $H_{\Poi}(e^{\mug})$, $\mug=\gain\,\pmean(x^*)+\latz$;
then \eqref{eq:daboxed} becomes
$\Uda(\scal x_t)=\Hmarg(\mug,\varg)-H_{\Poi}(e^{\mug})$ with
$\mug=\gain\scal\,\pmean(x_t)+\latz$, $\varg=\gain^2\scal^2\,\pvar(x_t)$.

\subsubsection{Growth rate of the diverging utility}
\begin{proposition}[Gaussian observation model: $\log\scal$ growth]\label{prop:gauss-logc}
If observations were Gaussian, $r(x)\sim\Gauss(\lat(x),\tau^2)$ with fixed $\tau^2$,
then for $\corr^2=1$, $\pvar=\scal^2\pvar(x_t)$, the DA utility grows as $\log\scal$.
\end{proposition}
\begin{proof}
$\Hmarg=\tfrac12\log(2\pi e(\pvar+\tau^2))$ and, since $\pvar_{\mathrm{cond}}=0$,
$\Hcond=\tfrac12\log(2\pi e\,\tau^2)$, so
$\Uda=\tfrac12\log(1+\pvar/\tau^2)$. With $\pvar=\scal^2\pvar(x_t)$ and large
$\scal$, $\Uda\approx\tfrac12\log(\scal^2\pvar(x_t)/\tau^2)=\log\scal+\text{const}$.
\end{proof}
Unbounded but slow: $d\Uda/d\scal\sim1/\scal$, so ascent slows but never stops.
Norm scaling is \emph{not} unique to Poisson --- the distinction is the \emph{rate}.

\begin{proposition}[Poisson-GP: $\scal^2$ heuristic bound]\label{prop:poisson-c2}
For the Poisson-GP case, under the heuristic $\Hmarg\ge H_{\Poi}(\E[f])$ with
$\E[f]=e^{\mug+\varg/2}$ and the large-rate Poisson entropy
$H_{\Poi}(\lambda)\approx\tfrac12\log(2\pi e\lambda)$,
\[
  \Hmarg\ge\tfrac12\log(2\pi e)+\tfrac12(\mug+\varg/2),
  \qquad\text{dominant term }\tfrac14\varg=\tfrac14\gain^2\scal^2\,\pvar(x_t).
\]
Since $\Hcond=H_{\Poi}(e^{\mug})$ depends only on $\mug$ (grows as $\scal$, not
$\scal^2$),
\[
  \Uda=\Hmarg-\Hcond\ \gtrsim\ \tfrac14\varg\ \propto\ \scal^2
\]
(faster than the Gaussian $\log\scal$).
\end{proposition}
\caveat{The bound $\Hmarg\ge H_{\Poi}(\E[f])$ is \emph{plausible but not proved}.
A Poisson mixture over log-normal rates is overdispersed (``more spread''), but
the Poisson is \emph{not} the max-entropy law on $\mathbb{N}_0$ under a mean
constraint alone (the geometric is); it is max-entropy only under $\mathrm{Var}=\text{mean}$,
so the naive argument does not directly apply. Numerically, fitting
$\Uda\sim\scal^\alpha$ over the Laplace-valid range
$\scal\in\{2,5,10\}$ (all $\sum p(r)>0.95$) gives $\alpha\approx1.9$, consistent
with $\scal^2$. The Laplace approximation truncates at $\rmax$ and breaks once
$\varg$ puts mass above $\rmax$; the safety score
$z_{\mathrm{safe}}=(\log\rmax-\mug)/\sqrt{\varg}$ falls below $2$ (unreliable) at
$\scal\approx10$ in-model, limiting the verifiable range.}

\subsubsection{The $\sigz^2>0$ damping and the alignment peak at $\scal=1$}
With $\sigz^2>0$ the results hold only approximately:
\[
  \vmag_{\scal x_t}=\scal^2 q+\sigz^2\neq\scal^2\vmag_{x_t},\qquad
  \cos\kang(\scal x_t,z)=\frac{\scal\,x_t\transpose\Cmat z+\sigz^2}{\sqrt{(\scal^2 q+\sigz^2)\,\vmag_z}}\neq\cos\kang(x_t,z).
\]
The fractional error $|\Kf(\scal x,z)-\scal\Kf(x,z)|/|\scal\Kf(x,z)|$ converges to a
\emph{constant} $O(\sigz^2/\vmag_{x_t})$ as $\scal\to\infty$ (not to zero: $\sigz^2$
shifts the angle by an amount $\propto\sigz^2/\vmag_{x_t}$, independent of $\scal$).
In a fitted model instance ($\vmag_{x_t}=q+\sigz^2\approx669$, $\sigz^2\approx0.98$, so
$\sigz^2/\vmag_{x_t}\approx0.0015$): the measured fractional error in
$\Kf(\scal x_t,\ipt_j)\to\approx0.6\%$ for large $\scal$, $\corr^2\to\approx0.987$,
and the residual variance ratio $\approx1.3\%$ (conditioning stays very effective,
not perfect). Crucially, the cosine similarity
\[
  \cos\kang(\scal)=\frac{\scal q+\sigz^2}{\sqrt{(\scal^2 q+\sigz^2)(q+\sigz^2)}}
\]
is \emph{strictly maximized at $\scal=1$} when $\sigz^2>0$ (and $\equiv1$ for all
$\scal$ when $\sigz^2=0$). So $\sigz^2$ breaks scale invariance: the kernel
\emph{wants} to match the target in direction \emph{and} magnitude --- but the
magnitude still scales as $\Kf(x^*,x^*)=\scal^2 q+\sigz^2=O(\scal^2)$, so the
bounded correlation reward cannot compete with the unbounded variance reward, and
ascent still runs $\scal\to\infty$.

\begin{remark}[Root cause and remedy]
The mechanism is a decoupling of norm from direction: $\Kf(\scal x,y)=\scal\Kf(x,y)$
at $\sigz^2=0$ lets ascent raise variance (hence entropy and utility) via the norm
degree of freedom without sacrificing directional alignment. Any practical
gradient-based input optimization with this kernel needs a norm constraint or a
normalized kernel (\S\ref{subsec:kernelfixes}). The blow-up is strongly
angle-dependent: effective at small angles ($\kang<0.2$, strong
conditioning), negligible at large angles ($\kang>0.6$).
\end{remark}

Table~\ref{tab:divevidence} shows the fitted-model evidence: the DA-optimized
input finds a \emph{genuine} loophole ($\Uda=0.800$ at a small angle
$\kang=0.114$, variance ratio $0.060\Rightarrow\corr^2\approx0.94$), while the
standard-optimized input diverges to $\Ustd=\infty$ by numerical
overflow at $\mug=1816$ --- two different phenomena
(\S\ref{subsec:numerics} separates them).

\begin{table}[t]\centering\small
\caption{Fitted model ($\Mind=50$, $\Ntr=50$; $\gain=0.0265$,
$\latz=-1.686$). The small gain $\gain$ compresses the raw GP mean into the useful
range, so the DA loophole ($\Uda=0.800$) is real, not a numerical artifact; the
standard-optimized point overflows float32.}
\label{tab:divevidence}
\begin{tabular}{lrrrrrrr}
\toprule
Input & $\pmean$ & $\pvar$ & $\mug$ & $\varg$ & $\kang$(rad) & $\Hmarg$ & $\Uda$\\
\midrule
$x_t$ (target) & $7.86$ & $70.9$ & $-1.48$ & $0.050$ & $0$ & $0.593$ & $0.010$\\
$x^*_{\mathrm{DA}}$ & $101$ & $3240$ & $0.99$ & $2.28$ & $0.114$ & $2.836$ & $0.800$\\
$x^*_{\mathrm{Std}}$ & $68497$ & $2.85\!\times\!10^{8}$ & $1816$ & $2.0\!\times\!10^{5}$ & $0.866$ & $\sim0$ & $\infty$\\
\bottomrule
\end{tabular}
\end{table}

% ===========================================================================
\subsection{Kernel solutions (Proof II)}
\label{subsec:kernelfixes}

Two kernel-level remedies bound the variance term that
\S\ref{subsec:divergence} showed to be the culprit. Decompose the utility (Poisson
link, $\rate=e^{\lat}$) schematically as
\[
  \Ustd(x^*)\ \approx\ \underbrace{\tfrac12\pvar(x^*)}_{\text{variance term, }O(\scal^2)}\ +\ \underbrace{I(\corr^2)}_{\text{correlation term, }\le\text{const}} ,
\]
where the correlation term is bounded (maximized at $\scal=1$ once $\sigz^2>0$) but
the variance term is unbounded; to fix the divergence, bound the variance term.

\subsubsection{Normalized arc-cosine kernel (project onto the unit sphere)}
\begin{proposition}[Scale-invariant normalized kernel]\label{prop:normkernel}
Define $\bar{\Kf}(x,x')=\Kf(x,x')/\sqrt{\Kf(x,x)\,\Kf(x',x')}$. For the order-1
arc-cosine kernel, $\bar{\Kf}(x,x')=\Jang(\kang)/\pi$, and for any $\scal>0$,
$\bar{\Kf}(\scal x,\scal x)=1=\bar{\Kf}(x,x)$ --- the prior variance is strictly
constant, independent of $\lVert x\rVert$.
\end{proposition}
\begin{proof}
By Lemma~\ref{lem:selfkernel} $\Kf(x,x)=\vmag_x$, so
$\bar{\Kf}(x,x')=[\tfrac1\pi\sqrt{\vmag_x\vmag_{x'}}\,\Jang(\kang)]/\sqrt{\vmag_x\vmag_{x'}}=\Jang(\kang)/\pi$,
which depends only on $\kang$. In particular $\bar{\Kf}(x,x)=\Jang(0)/\pi=1$; and
$\bar{\Kf}(\scal x,\scal x)=\Kf(\scal x,\scal x)/\sqrt{\Kf(\scal x,\scal x)^2}=1$.
\end{proof}
With $\bar{\Kf}(x,x)=1$ and cross terms $\bar{\Kf}(x,\ipt_j)=\Jang(\kang_j)/\pi\in[0,1]$
bounded, the whole posterior-correction term is bounded, so $\bar{\pvar}(x)$ is
bounded regardless of $\lVert x\rVert_{\Cmat}$ and $\Uda$ cannot diverge by
scaling: the only way up is to improve $\corr^2$ (small angle to $x_t$; with
$\sigz^2>0$ retained inside $\kang$, minimized only when $x\to x_t$ in both
direction and magnitude).
\caveat{The normalized kernel discards all norm information, but the input norm
carries genuine predictive signal (the target depends on input magnitude, not
direction alone). Empirically, training with $\bar{\Kf}$ drops test correlation by
$\sim25\%$ --- the norm encodes real structure the model needs. A fix for the
input-synthesis divergence is not a free improvement to the fitted model.}

\subsubsection{Saturation kernel (arc-sin)}
The arc-cosine kernel is the covariance of infinite-width ReLU networks, which do
\emph{not} saturate; a bounded activation gives a self-magnitude that saturates
instead of growing without bound. The arc-sin (probit-activation) kernel models
saturation:
\[
  \Kf_{\mathrm{sat}}(x,x')=\satf\cdot\tfrac2\pi\,\arcsin\!\left(\frac{x\transpose\Cmat x'+\sigz^2}{\sqrt{(1+x\transpose\Cmat x+\sigz^2)(1+{x'}\transpose\Cmat x'+\sigz^2)}}\right).
\]
\begin{proposition}[Bounded self-variance]\label{prop:satkernel}
$\displaystyle\lim_{\lVert x\rVert\to\infty}\Kf_{\mathrm{sat}}(x,x)=\satf\cdot\tfrac2\pi\arcsin(1)=\satf$.
\end{proposition}
\begin{proof}
As $\lVert x\rVert\to\infty$ the ratio inside $\arcsin$ tends to
$(x\transpose\Cmat x)/(1+x\transpose\Cmat x)\to1$, and $\arcsin(1)=\pi/2$, so
$\Kf_{\mathrm{sat}}(x,x)\to\satf\cdot(2/\pi)(\pi/2)=\satf$.
\end{proof}
Unlike $\Kf(x,x)\propto\lVert x\rVert^2$, the variance saturates to the constant
$\satf$. This fixes the divergence \emph{without} normalizing the kernel: $\Hmarg$
cannot grow indefinitely with norm, so to raise utility the optimizer must lower
$\Hcond$ by finding $x^*$ highly correlated with $x_t$; inputs of unbounded norm
yield diminishing returns, embedding the constraint that input magnitude beyond a
point yields no more signal variance.

% ===========================================================================
\subsection{Numerical analysis of the standard utility}
\label{subsec:numerics}

This subsection is entirely about the standard utility \eqref{eq:Ustd},
$\Ustd=\Hmarg-\E_{\lgr}[\Hnoise]=I(R;\lgr\mid x^*,\mathcal{D})$. The count index
$r$ below indexes the count observation ($\spk$ of Part~I); the count random
variable is $R\sim\Poi(e^{\lgr})$.

\subsubsection{The closed-form decomposition}
With $\lat(x^*)\mid\mathcal{D}\sim\Gauss(\pmean,\pvar)$, $\lgr=\gain\lat+\latz$, so
$\lgr\mid\mathcal{D}\sim\Gauss(\mug,\varg)$, and $R\mid\lgr\sim\Poi(e^{\lgr})$:
\[
  \Hmarg=-\sum_{r=0}^{\infty}p(r)\log p(r),\qquad p(r)=\int\Poi(r;e^{\lgr})\,\Gauss(\lgr;\mug,\varg)\,d\lgr,
\]
the marginal (Poisson--log-normal) entropy of \eqref{eq:Hmarg} (no closed form;
Laplace-approximated below). The subtracted expected-noise entropy \emph{is}
closed-form via the log-normal MGF ($\E[e^{\lgr}]=e^{\mug+\varg/2}$,
$\E[\lgr e^{\lgr}]=e^{\mug+\varg/2}(\mug+\varg)$):
\begin{equation}\label{eq:Ehnoise}
  \E_{\lgr}[\Hnoise]=\underbrace{-e^{\mug+\varg/2}(\mug+\varg-1)}_{\text{exact (MGF), closed form}}\ +\ \underbrace{\sum_{r=0}^{\infty}p(r)\log r!}_{\text{needs }p(r)\text{, Laplace}} ,
\end{equation}
so $\Ustd\approx-\sum_{r=0}^{\rmax}p(r)\log p(r)-\bigl(-e^{\mug+\varg/2}(\mug+\varg-1)+\sum_{r=0}^{\rmax}p(r)\log r!\bigr)$;
both sums truncate at the same $\rmax$.

\subsubsection{Laplace approximation and the log-space Lambert-W fix}
The saddle point of $p(r)$ is at the mode $\gmode$ of
$\Poi(r;e^{\lgr})\,\Gauss(\lgr;\mug,\varg)$, solving $r=e^{\gmode}+(\gmode-\mug)/\varg$:
\begin{equation}\label{eq:lambertmode}
  \gmode(r)=r\varg+\mug-\LW\!\bigl(\varg\,e^{\,r\varg+\mug}\bigr),
\end{equation}
$\LW$ the principal Lambert branch. \textbf{The overflow failure mode:} forming
$\varg e^{\,r\varg+\mug}$ literally overflows float32 whenever
$r\varg+\mug\gtrsim88$ ($e^{88}\sim10^{38}$); masking the overflowing $r$ lets the
truncated $\sum p(r)$ exceed $1$ by large factors. (This is the historic
``$\exp(\mu)$ diverges.'') \textbf{The log-space fix} never forms $e^y$:
set $y=\log\varg+r\varg+\mug$ (finite for finite inputs) and solve
\[
  w+\log w=y\quad\Longleftrightarrow\quad w=\LW(e^y)
\]
by a fixed $10$ Newton iterations directly on the log-form (the backward
derivative uses $dW/dy=W/(1+W)$, in which $e^y$ cancels), then
$\gmode(r)=r\varg+\mug-w$. Since $w\approx y-\log y$ for large $y$,
$\gmode(r)\approx-\log\varg+\log y\sim\log r$, so $e^{\gmode(r)}\approx r$ (exactly
right: the Poisson likelihood is maximized when rate $=$ count) and stays in
float32 range for any $r$ --- so $\exp(\gmode)$ is no longer an overflow risk
post-fix.

\subsubsection{Four independent residual failure modes}
Distinct thresholds and remedies (they coincide only because all worsen at large
$\mug/\varg$).

\paragraph{(a) Laplace approximation error in $\log p(r)$.} A local quadratic
model of the integrand; inaccurate when it departs from a quadratic-log shape.
Two regimes: $\varg\to0$ (prior degenerates, saddle singular --- at $\varg<10^{-6}$
one uses the \emph{exact} Poisson $\log p=\mug r-e^{\mug}-\log r!$ instead);
$\varg$ large (flat prior; the asymmetric Poisson factor skews the integrand).
\emph{Independent of $\rmax$} --- each $p(r)$ has its own bias, observable as
$\sum p(r)\neq1$ even with no missing mass (up to $1.015$; that is Laplace error,
not truncation). Measured (vs numerical quadrature): error $<1\%$ for $\varg\lesssim1$
at every $\mug$, growing to $5$--$20\%$ at the peak by $\varg\sim50$ (worst where
the predictive peak slides onto $r=0$). The driver is $\varg$, not $\mug$; float32
vs float64 shifts $\log p(r)$ by $\le3.7\times10^{-4}$ nats. \emph{Range gate:}
the count is bounded in the regime of interest ($\lesssim100$ per observation
window), so the sensible regime is $\varg\lesssim6$, within which (a) is a
few-percent effect.

\paragraph{(b) Truncation at $\rmax$.} If the predictive puts non-negligible mass
above $\rmax$, both $-p(r)\log p(r)$ and $p(r)\log r!$ are clipped, biasing
\emph{both} $\Hmarg$ and the closed-sum term low. The moments are
\[
  \E[R]=e^{\mug+\varg/2},\qquad
  \mathrm{Var}(R)=e^{\mug+\varg/2}+e^{2\mug+\varg}(e^{\varg}-1),
\]
the second (log-normal-of-rate) term dominating once $\varg\gtrsim1$. As $\varg$
grows the predictive is strongly right-skewed and its centres separate: the mean
$e^{\mug+\varg/2}$ climbs, the median stays $\approx e^{\mug}$, the mode drifts
\emph{down} toward $0$ ($\approx e^{\mug-\varg}$) --- so ``$p(r)$ peaks at
$e^{\mug}$'' is only the small-$\varg$ limit. A sufficient benign-truncation
condition is $\rmax\gtrsim\E[R]+k\sqrt{\mathrm{Var}(R)}$. An adaptive rule is
stricter (a $k$-$\sigma$ upper tail in $\lgr$-space, exponentiated to a rate, plus
a Poisson std, plus a margin):
\[
  \rmax^{\mathrm{adapt}}=e^{\mug+k\sigma_g}+5\sqrt{e^{\mug+k\sigma_g}}+10,\quad k=3,
\]
clamped to $[200,10000]$, raising an error (rather than silently over-truncating)
if the needed $\rmax$ exceeds the ceiling; a further guard triggers if the
required bound exceeds $e^{80}$, so the overflow clamp cannot hide failures. The
fixed default $\rmax=100$ is benign for confident predictions but \emph{not safe
in general}: it fails whenever $\mug\gtrsim\log100\approx4.6$ \emph{or} $\varg$
grows large --- exactly the high-uncertainty regime active learning explores.

\paragraph{(c) Float32 overflow of the exact closed-form term.} The exact term
$-e^{\mug+\varg/2}(\mug+\varg-1)$ overflows float32 when
$\mug+\tfrac12\varg\gtrsim88$, sending $\E_{\lgr}[\Hnoise]\to-\infty$ and
$\Ustd\to+\infty$. This is the \emph{only} residual ``$\exp$ can diverge'' issue,
and it lives in an \emph{exact analytical} term, not the Laplace approximation ---
distinct from the mode-finding overflow above. In practice truncation (b) bites
long first ($\sigma_g\sim1\Rightarrow$ truncation near $\mug\sim1$ at $\rmax=100$, vs
$\mug\sim88$ here).

\paragraph{(d) Catastrophic cancellation.} Inherent to the decomposition (survives
even with exact Laplace and no truncation). \emph{Inside} $\E_{\lgr}[\Hnoise]$
(large, earliest): both $-e^{\mug+\varg/2}(\mug+\varg-1)$ and $\sum p(r)\log r!$
grow $\sim e^{\mug}$, while the true $\E_{\lgr}[\Hnoise]$ grows only \emph{linearly}
($\sim\tfrac12\mug$); at $\mug=10$ each term $\sim2\times10^5$ vs a true value
$\sim6$ ($\sim4.5$ digits cancel), at $\mug=15$ $\sim6.7$ digits cancel ---
essentially noise in float32's $\sim7$-digit mantissa (and truncating the second
sum makes this \emph{worse}). \emph{Between} $\Hmarg$ and $\E_{\lgr}[\Hnoise]$
(smaller, second-level): by Jensen $\Hmarg\ge\E_{\lgr}[\Hnoise]$ while both grow
linearly in $\mug$, a second cancellation costing fewer digits. Unfixable without
rewriting the decomposition so cancelling pieces are computed together (e.g.\
summing $-p(r)\log p(r)-p(r)\log r!$ as a single term before any large analytical
exponential). Flagged, out of scope.

\subsubsection{``Utility diverges under gradient ascent,'' re-examined}
The claim is \emph{two} distinct things. (i) \textbf{Kernel-driven input blow-up:}
$\Kf(\scal x,\scal x)=\scal^2\Kf(x,x)\Rightarrow$ ascending $\Ustd$ with no norm
constraint sends $\pvar\propto\lVert x\rVert^2\to\infty$, hence $\varg\to\infty$
(the mechanism of \S\ref{subsec:divergence}--\ref{subsec:kernelfixes}). (ii) The
\textbf{computed} utility then misbehaves via the four numerical modes, typically
in order: truncation collapse of $\Hmarg$ (b), then Laplace error at large $\varg$
(a), then cancellation inside $\E_{\lgr}[\Hnoise]$ (d) as $\mug$ grows, finally
float32 overflow (c) at extreme scales. The kernel makes the inputs unbounded;
then numerics make the computed utility ill-behaved --- not the same thing, and
the four numerical effects are not the same thing either. Numerical fixes alone
will not cure it: either constrain the optimizer (norm projection, or the
normalized kernel of Proposition~\ref{prop:normkernel}) or rewrite the closed-form
term in an asymptotically stable form.

\caveat{Utility-accuracy ceiling (known limitation). The standard utility can run
to $+\infty$ and the DA utility can silently collapse, both via the $\rmax=100$
truncation, at high rate; a bound on the maximal rate $f_{\max}$ keeps the
operating point in the accurate regime. This is a known limitation, not a resolved
result.}


\part{From selection to synthesis: guided diffusion}
% ============================================================================
%  Part IV -- From selection to synthesis: guided diffusion  (section fragment)
%  \input-ready; uses ONLY macros from preamble.tex. No preamble here.
%  Cross-part refs used: \eqref{eq:gradxK} (Part I kernel input-gradient),
%  \eqref{eq:Ustd} \eqref{eq:Uda} (Part III utilities),
%  \ref{subsec:numerics} (Part III utility numerics / truncation ceiling).
%  Notation clash discipline (preamble): the diffusion schedule \bsched/\aret/\abar
%  and timestep \dt are DISTINCT from the prior half-width \bwid of the localized
%  region and the locality mask \loc; \dt is a scalar diffusion-time index, NOT a
%  grid coordinate -- the prior is over the 2-D grid only.
% ============================================================================

\section{Guided diffusion for input synthesis}\label{sec:diffusion}

Parts I--III chose the next input by \emph{selecting} the highest-utility member
of a fixed pool. This part develops the natural extension: \emph{synthesize} a new
input, drawn from the input (data) distribution $\px$, that maximizes the same
Gaussian-process (GP) utility for a given target model. A denoising diffusion model
supplies a learned prior over $\px$, and the GP utility gradient---the very object
of Part~III---steers generation toward high-utility inputs. The two ingredients
divide the labor cleanly: the utility gradient (smooth, low-pass;
\S\ref{sec:diff-motiv}) fixes \emph{where} in input space to move, and the
diffusion score supplies the high-frequency structure of $\px$ that a utility
gradient can never produce on its own.

% ----------------------------------------------------------------------------
\subsection{Motivation: synthesize rather than select}\label{sec:diff-motiv}

\paragraph{Why direct gradient ascent on the utility fails.}
Suppose we optimize an input $x^{*}$ directly to maximize a GP utility $U(x^{*})$
by gradient ascent. The gradient flows through the arc-cosine kernel's
input-gradient $\nabla_{x}\Kf$ (Part~I, \eqref{eq:gradxK}), which carries the
structured (grid) prior $\Cmat$ and hence its Gaussian smoothing block
$\Csm$---a convolution with a Gaussian kernel, i.e.\ a \emph{low-pass filter}. Two
consequences follow, independent of the inducing points, the training data, or
which utility is used:
\begin{enumerate}[leftmargin=1.4em,itemsep=1pt,topsep=2pt]
  \item the utility gradient is \emph{inherently smooth}, so ascent iterates lack
        the high-frequency structure (sharp transitions, fine texture, local
        contrast) that characterizes draws from $\px$; and
  \item ascent pushes grid points outside the bounded input box. The arc-cosine
        self-kernel is positively homogeneous, $\Kf(x,x)=x\transpose
        \Cmat x+\sigz^{2}=\lVert x\rVert_{\Cmat}^{2}+\sigz^{2}$, so larger norms
        are rewarded and the utility diverges under input scaling (the
        norm-scaling divergence of Part~III).
\end{enumerate}
Such iterates are not admissible inputs: to lie in the support of $\px$ they must
be indistinguishable from genuine draws from the input distribution.

\paragraph{Why a subspace constraint is not enough.}
Restricting the optimization to a low-dimensional subspace (PCA space, or the
$\Cmat$-eigenspace of Part~III) does not cure this. PCA space fixes only the
\emph{second-order} statistics; under approximate stationarity of $\px$ its
eigenvectors tend to Fourier modes (Szeg\H{o}'s theorem), so PCA-space ascent is
essentially Fourier low-pass filtering. Neither PCA space nor the
$\Cmat$-eigenspace enforces the higher-order structure (phase coherence, sparse
gradients) that places a sample in the support of $\px$: a subspace constraint is a
convergence aid, not a $\px$-fidelity constraint.

\paragraph{The decomposition, and why diffusion.}
The synthesis task has two objectives that pull apart cleanly:
(i)~\emph{informativeness}---$x^{*}$ maximizes the utility (reduces output
uncertainty); and (ii)~\emph{fidelity}---$x^{*}$ is a plausible draw from the input
distribution $\px$. Gradient ascent handles (i) but not (ii). A diffusion model
trained on samples from $\px$ learns the \emph{score}
$\nabla_{x}\log p_{\dt}(x)$ at every noise level $\dt$; at low noise this encodes
the full statistical structure of $\px$. Guided generation adds the two:
\begin{equation}
\underbrace{\score(\xt,\dt)}_{\text{high-freq.\ structure}}
\;+\;
\underbrace{\gw\,\nabla_{\xt}U(\xhat)}_{\text{low-freq.\ utility bias}} .
\label{eq:diff-decomp}
\end{equation}

\begin{remark}[Division of labor]
The smoothness that ruined direct optimization becomes acceptable under
\eqref{eq:diff-decomp}: the utility only has to \emph{bias} generation toward
informative regions, not prescribe individual grid points. The score model supplies
all orders of the statistics of $\px$ (trained once, reused across targets and
utilities); the utility gradient supplies the target-specific, low-frequency
direction. They act at complementary spatial-frequency scales.
\end{remark}

\paragraph{Connection to the selection setting.}
In the fixed-pool active-learning loop, one \emph{selects} the pool input of
highest expected information gain about the latent function. Diffusion-guided
synthesis is the extension: \emph{manufacture} an input of even higher utility that
still lies in the support of $\px$---closing the gap between ``best available pool
input'' and ``best attainable informative input.'' This part is framed by three
questions: (1)~can we generate inputs that are simultaneously plausible under $\px$
and high-utility?  (2)~how does that input change as the GP sharpens from
$\Ntr=50$ to $300$ training inputs?  (3)~how much utility does the $\px$-fidelity
constraint cost---the gap between the unconstrained pure-utility optimum and the
diffusion-guided one, as a function of guidance strength $\gw$?

% ----------------------------------------------------------------------------
\subsection{Diffusion fundamentals}\label{sec:diff-fund}

\subsubsection{Forward (noising) process}
Let $\xz\in\Reals^{d}$ be a clean input ($d=64\times64=4096$ here). Fix a variance
schedule $\bsched\in(0,1)$, $\dt=1,\dots,T$, and define the per-step signal
retention $\aret=1-\bsched$ and its cumulative product
$\abar=\prod_{s=1}^{\dt}\alpha_{s}=\prod_{s=1}^{\dt}(1-\beta_{s})$. The forward
process adds Gaussian noise one step at a time,
\begin{equation}
q(\xt\mid x_{\dt-1})=\Gauss\!\big(\xt;\ \sqrt{1-\bsched}\,x_{\dt-1},\ \bsched\,I\big),
\label{eq:forward-step}
\end{equation}
and admits a closed-form marginal at arbitrary $\dt$ (the key property---no
intermediate steps are needed):
\begin{equation}
q(\xt\mid\xz)=\Gauss\!\big(\xt;\ \sqrt{\abar}\,\xz,\ (1-\abar)\,I\big).
\label{eq:forward-diff}
\end{equation}
Equivalently, by the reparameterization
\begin{equation}
\xt=\sqrt{\abar}\,\xz+\sqrt{1-\abar}\,\epsilon,\qquad \epsilon\sim\Gauss(0,I).
\label{eq:reparam}
\end{equation}
The schedule is chosen so that $\bar\alpha_{T}\approx0$, hence
$q(x_{T}\mid\xz)\approx\Gauss(0,I)$: at the end of the forward process all
information about $\xz$ is destroyed.

\subsubsection{Score $\leftrightarrow$ noise, and the training objective}
Differentiating the log of the forward marginal \eqref{eq:forward-diff} gives the
score in closed form, which identifies it with the added noise:
\begin{equation}
\nabla_{\xt}\log q(\xt\mid\xz)
   =-\frac{\xt-\sqrt{\abar}\,\xz}{1-\abar}
   =-\frac{\epsilon}{\sqrt{1-\abar}}
\quad\Longrightarrow\quad
\nabla_{\xt}\log p_{\dt}(\xt)\approx\score(\xt,\dt)=-\frac{\epsth(\xt,\dt)}{\sqrt{1-\abar}}.
\label{eq:score-noise}
\end{equation}
A network trained to predict the noise therefore approximates the score. Training
minimizes the simplified denoising-score-matching (DSM) $L_{2}$ loss
\begin{equation}
\mathcal{L}_{\mathrm{DSM}}
  =\E_{\dt\sim U(1,T)}\,\E_{\xz\sim\px}\,\E_{\epsilon\sim\Gauss(0,I)}
   \Big[\big\lVert\epsth(\xt,\dt)-\epsilon\big\rVert^{2}\Big],
\label{eq:dsm}
\end{equation}
with $\xt$ formed by \eqref{eq:reparam}. Each mini-batch draws a fresh
$(\xz,\dt,\epsilon)$; the network learns all noise levels though it sees each
infrequently.

\subsubsection{Reverse (denoising) process and the Tweedie estimate}
The generative reverse process is Gaussian,
$p_{\theta}(x_{\dt-1}\mid\xt)=\Gauss(x_{\dt-1};\mu_{\theta}(\xt,\dt),\sigma_{\dt}^{2}I)$,
with mean read off from the noise predictor
\begin{equation}
\mu_{\theta}(\xt,\dt)=\frac{1}{\sqrt{1-\bsched}}
   \Big[\,\xt-\frac{\bsched}{\sqrt{1-\abar}}\,\epsth(\xt,\dt)\Big],
\label{eq:reverse-mean}
\end{equation}
and fixed variance $\sigma_{\dt}^{2}$, either the simple choice
$\sigma_{\dt}^{2}=\bsched$ or the posterior form
$\tilde\beta_{\dt}=\bsched(1-\bar\alpha_{\dt-1})/(1-\abar)$; in practice
$\sigma_{\dt}^{2}=\bsched$ works well. Sampling runs $\dt=T,\dots,1$ from
$x_{T}\sim\Gauss(0,I)$, drawing $x_{\dt-1}=\mu_{\theta}+\sigma_{\dt}z$ with
$z\sim\Gauss(0,I)$ ($z=0$ at $\dt=1$). At any $\dt$ the noise prediction yields the
\emph{Tweedie estimate} of the clean input---the model's posterior mean
$\E[\xz\mid\xt]$:
\begin{equation}
\xhat=\frac{\xt-\sqrt{1-\abar}\,\epsth(\xt,\dt)}{\sqrt{\abar}}.
\label{eq:tweedie}
\end{equation}
It is coarse at high noise and increasingly detailed as $\dt\downarrow$. The
Tweedie estimate \eqref{eq:tweedie} is the workhorse of guidance
(\S\ref{sec:diff-guide}): it maps a noisy intermediate to a clean input on which
the GP utility can be evaluated, and it is a differentiable function of $\xt$.

\begin{remark}[$T-1$ start]
At $\dt=T$ the schedules used here have $\bar\alpha_{T}\!\sim\!10^{-3}$, so
$1/\sqrt{\bar\alpha_{T}}\!\sim\!31$ and the first reverse step amplifies any
prediction error $\sim\!31\times$, causing divergence. Sampling therefore starts
the loop at $\dt=T-1$ (where $1/\sqrt{\alpha}\!\approx\!2$, stable).
\end{remark}

\subsubsection{The U-Net noise predictor and DDIM acceleration}
Because $\epsth(\xt,\dt)$ maps a noisy input to a same-size tensor, it is
parameterized by an encoder--decoder with \emph{skip connections} (a U-Net): a
plain bottleneck would discard exactly the fine detail that distinguishes a draw
from $\px$ from noise, whereas skip connections carry per-resolution detail
directly to the decoder. The scalar timestep $\dt$ enters through a sinusoidal
embedding (Transformer-style),
\[
\mathrm{emb}(\dt)_{2k}=\sin\!\big(\dt/10000^{2k/d_{\mathrm{emb}}}\big),
\qquad
\mathrm{emb}(\dt)_{2k+1}=\cos\!\big(\dt/10000^{2k/d_{\mathrm{emb}}}\big),
\]
passed through a small MLP and \emph{added} (broadcast over space) to each block's
feature map---a global ``how much to activate at this noise level'' modulation.

Full DDPM sampling needs $T-1=999$ sequential U-Net passes. \emph{DDIM} replaces
the stochastic reverse chain by a deterministic one over a subsequence of $S$
timesteps $\tau_{S}>\dots>\tau_{1}$ (e.g.\ $S=50$). Going from $\dt=\tau_{i}$ to
$t'=\tau_{i-1}$,
\begin{equation}
x_{t'}=\sqrt{\bar\alpha_{t'}}\,\xhat
   +\sqrt{1-\bar\alpha_{t'}-\sigma^{2}}\;
     \frac{\xt-\sqrt{\abar}\,\xhat}{\sqrt{1-\abar}}
   +\sigma\,z,
\label{eq:ddim}
\end{equation}
with $\xhat$ from \eqref{eq:tweedie}. Here $\sigma=0$ gives a fully deterministic,
reproducible sampler; the DDPM choice of $\sigma$ recovers stochastic sampling.
Quality degrades below $S\!\approx\!20$; $S=50$ is the default.

% ----------------------------------------------------------------------------
\subsection{Guidance: steering the reverse process with the GP utility gradient}
\label{sec:diff-guide}

\subsubsection{Classifier-style guidance}
Model ``high utility'' as a label $y$ with $p(y\mid\xz)\propto\exp\!\big(\gw\,
U(\xz)\big)$, $\gw>0$. Bayes' rule splits the \emph{conditional} score into a
$\px$-fidelity term and a guidance term:
\begin{equation}
\nabla_{\xt}\log p(\xt\mid y)
  =\underbrace{\score(\xt,\dt)}_{\text{unconditional score, \eqref{eq:score-noise}}}
  +\;\gw\,\nabla_{\xt}\Ustd(\xhat(\xt,\dt)).
\label{eq:guidance}
\end{equation}
The utility is defined on \emph{clean} inputs, so it is evaluated on the Tweedie
estimate \eqref{eq:tweedie} (the standard Dhariwal--Nichol trick). Any of the
Part~III utilities may play the role of $U$ in \eqref{eq:guidance}: the standard
utility $\Ustd$ (\eqref{eq:Ustd}), shown here, or the distribution-aware utility
$\Uda$ (\eqref{eq:Uda}). Define the \emph{guidance gradient}
\begin{equation}
\ggrad=\nabla_{\xt}\Ustd\big(\xhat(\xt,\dt)\big).
\label{eq:ggrad}
\end{equation}

\subsubsection{The guided reverse update (two equivalent forms)}
Converting the guided score \eqref{eq:guidance} back to a noise prediction, and
substituting into the Tweedie map \eqref{eq:tweedie}, yields two algebraically
identical updates:
\begin{align}
\epsilon_{\mathrm{guided}} &= \epsth(\xt,\dt)-\gw\,\sqrt{1-\abar}\,\ggrad,
   \label{eq:guided-noise}\\[2pt]
\xhat^{\mathrm{guided}} &= \xhat+\frac{\gw\,(1-\abar)}{\sqrt{\abar}}\,\ggrad.
   \label{eq:guided-tweedie}
\end{align}
Form \eqref{eq:guided-noise} feeds the DDPM mean \eqref{eq:reverse-mean} in place
of $\epsth$; form \eqref{eq:guided-tweedie} feeds the DDIM update \eqref{eq:ddim}
in place of $\xhat$. The \emph{ramp factor} $(1-\abar)/\sqrt{\abar}$ in
\eqref{eq:guided-tweedie} is large at high noise ($\dt\!\approx\!T$, $\sim\!158$),
giving bold utility steering while global structure is still being decided, and
small at low noise ($\dt\!\approx\!1$, $\sim\!0.3$), preserving fine detail.

\subsubsection{The GP link: what the guidance signal is}\label{sec:diff-gplink}
The guidance gradient \eqref{eq:ggrad} is the gradient of the GP information
utility of Part~III, back-propagated through Tweedie. Both utilities are
differentiable in the input, and the gradient flows through the arc-cosine
kernel's input-gradient $\nabla_{x}\Kf$ (Part~I, \eqref{eq:gradxK}). Because
$\nabla_{x}\Kf$ carries the factor $\Cmat$---and thus the Gaussian spatial
smoother $\Csm$---the utility gradient is exactly the smooth, low-pass signal of
\S\ref{sec:diff-motiv}. This is precisely why guidance works: the utility supplies
a \emph{low-frequency} informative direction, and the score model supplies the
\emph{high-frequency} structure of $\px$ the smooth utility gradient can never
produce.

\paragraph{The locality mask $\Rightarrow$ guidance touches only the
localized-region grid points.}
The structured prior applies a Gaussian locality envelope to each grid point $j$,
\begin{equation}
\loc_{j}=\exp\!\Big(-\frac{\lVert\pixc_{j}-\rfc\rVert^{2}}{4\,\bwid^{2}}\Big),
\label{eq:rfmask}
\end{equation}
where $\pixc_{j}$ is grid point $j$'s coordinate, $\rfc$ the center of the
localized region, and $\bwid$ the \emph{kernel's localized-region half-width}.
Grid points with $\loc_{j}<10^{-3}$ are masked out (zero contribution), so
$\nabla_{x}U$ is nonzero \emph{only at localized-region grid points}. The mask area
scales as $\bwid^{2}$: at the default $\bwid=0.1$ the localized region covers
$\approx21\%$ of the $64\times64$ grid ($\approx869$ grid points), versus $5.4\%$
at $\bwid=0.05$ and $1.7\%$ at $\bwid=0.03$. The utility therefore guides only the
localized region, and \emph{the diffusion model freely generates every remaining
grid point} for $\px$-fidelity---an advantage over the L$_{2}$ approach
(\S\ref{sec:diff-approaches}), whose non-localized-region grid points are frozen at
a start input. Because $\bwid$ is a learnable kernel parameter, it must be frozen
during GP training for synthesis, or the intended mask size drifts.

\caveat{The symbol $\bwid$ in \eqref{eq:rfmask} is the \emph{kernel}
localized-region half-width (Part~I), \textbf{not} the diffusion noise-schedule
variance $\bsched$ of \eqref{eq:forward-step}; likewise the locality envelope
$\loc_{j}$ is the kernel locality mask, not the diffusion signal-retention $\aret$.
These collisions ($\beta$, $\alpha$) are the most dangerous in the whole document;
the macros keep them apart.}

\subsubsection{Full versus approximate guidance gradient}\label{sec:diff-approx}
By the chain rule $\ggrad$ factors through the utility gradient in clean-input
space and the Tweedie Jacobian,
\begin{equation}
\frac{\partial\xhat}{\partial\xt}
   =\frac{1}{\sqrt{\abar}}\Big(I-\sqrt{1-\abar}\,\frac{\partial\epsth}{\partial\xt}\Big),
\label{eq:tweedie-jac}
\end{equation}
where $\partial\epsth/\partial\xt$ is the $d\times d$ U-Net Jacobian, which is
never formed explicitly.
\begin{itemize}[leftmargin=1.4em,itemsep=1pt,topsep=2pt]
  \item \textbf{Full gradient:} retains $\partial\epsth/\partial\xt$ in
        \eqref{eq:tweedie-jac}; one U-Net forward \emph{and} one backward pass,
        $\approx2\times$ the cost of a plain step.
  \item \textbf{Approximate gradient} (default): treat the U-Net Jacobian as
        negligible. Then $\xhat$ depends on $\xt$ only through the direct
        $\xt/\sqrt{\abar}$ term, so
        \begin{equation}
        \ggrad^{\mathrm{approx}}=\frac{1}{\sqrt{\abar}}\,\nabla_{\xhat}\Ustd(\xhat).
        \label{eq:approx-grad}
        \end{equation}
        No U-Net backward pass ($\approx\tfrac12$ the cost). It is justified when
        $\epsth$ is a ``local correction'' whose Jacobian is roughly zero-mean
        (does not systematically align with or oppose the utility gradient), which
        holds best at \emph{low} noise (small $\sqrt{1-\abar}$); a hybrid
        (approximate for early $\dt>200$, full for late $\dt<200$) captures most
        of the benefit cheaply.
\end{itemize}

\subsubsection{The four approaches (A--D) and the adopted method}
\label{sec:diff-approaches}
Four guidance paradigms were catalogued; all attempt to add a fidelity term
counteracting the utility gradient's drift off the data manifold $\mathcal{M}$ (the
support of $\px$). Let $x_{c}$ denote a current input being optimized.

\begin{description}[leftmargin=1.6em,itemsep=2pt,topsep=2pt]
  \item[A --- L$_2$ penalty toward the Tweedie estimate.] Corrupt $x_{c}$ at a
        \emph{fixed} noise level $\dt^{*}$, predict $\hat\epsilon$, form $\xhat$
        (\eqref{eq:tweedie}), hold it fixed (no gradient through it), and penalize
        the distance to it:
        \begin{equation}
        \mathcal{L}_{\mathrm{A}}=-\Uda(x_{c})+\frac{\greg}{2}\,
           \big\lVert x_{c}-\xhat\big\rVert^{2},
        \qquad
        x_{c}\leftarrow x_{c}+\eta\,\nabla\Uda+\eta\,\greg\,(\xhat-x_{c}).
        \label{eq:approachA}
        \end{equation}
        Here $\greg$ is the approach-A regularization weight. Since
        $x_{c}-\xhat\propto(\hat\epsilon-\epsilon)$, \eqref{eq:approachA} is a
        Score-Distillation-Sampling gradient with a different time weighting.
        \emph{Weaknesses:} no signal near $\mathcal{M}$ (for $x_{c}$ already on
        $\mathcal{M}$, $\xhat\!\approx\!x_{c}$ so $\hat\epsilon-\epsilon$ is just
        reconstruction noise---the penalty is reactive, not preventive); the
        Tweedie estimate of an amplified input is itself amplified, so norm growth
        is not prevented; $\dt^{*}$ is arbitrary; a single fixed noise sample is
        biased.
  \item[B --- direct score.] Use $\score(\xt,\dt^{*})=-\hat\epsilon/\sqrt{1-
        \bar\alpha_{\dt^{*}}}$ directly, chain-ruled to $x_{c}$-space. Uses
        $\hat\epsilon$ alone (not $\hat\epsilon-\epsilon$), so it has signal on
        $\mathcal{M}$ but high variance, and estimates the score of the
        \emph{blurred} $p_{\dt^{*}}$, not $p_{0}$.
  \item[C --- score without added noise.] Feed $x_{c}$ at $\dt=1$ directly.
        Appealing ($p_{1}\!\approx\!p_{0}$) but fails: a utility-distorted $x_{c}$
        is far out-of-distribution for a network expecting ``a draw from $\px$ plus
        $\sigma\!\approx\!0.01$ noise,'' and dividing by
        $\sqrt{1-\bar\alpha_{1}}\!\approx\!0.01$ amplifies the unreliable
        prediction $\sim\!100\times$.
  \item[D --- full guided reverse diffusion (\emph{adopted}).] Do not optimize an
        existing input---\emph{generate} from noise, nudging each denoising step
        by the utility gradient via \eqref{eq:guidance}--\eqref{eq:guided-tweedie}.
        The U-Net is \emph{always} in-distribution ($\xt$ is genuinely at noise
        level $\dt$ because it came from the reverse process); generation
        traverses coarse$\to$fine over all noise levels; there is no $\dt^{*}$ to
        choose; the output is a proper sample from the guided distribution. Costs:
        the U-Net forward (with gradients in full mode), $T$ (or $\sim\!50$ DDIM)
        sequential steps, utility evaluated on rough early Tweedie estimates,
        stochastic output (best-of-$K$), and no start-point control.
\end{description}

\begin{remark}[Why the score is a genuine restoring force]
Naive per-coordinate clipping $\Pi_{\mathcal{B}}$ (each $x_{i}$ clamped to
$[x_{\min},x_{\max}]$) acts on each grid point independently and ignores the joint
distribution $\px$: it lands on the boundary of the input box but far from
$\mathcal{M}$. In the guided score \eqref{eq:guidance} the score term is instead a
\emph{structured} restoring force on the joint distribution---if a region saturates
and loses its characteristic spectrum, $p_{\dt}(\xt)$ drops sharply and the score
rearranges grid points to restore the correlations of $\px$ while accommodating the
utility bias. The limitation is out-of-distribution failure: $\epsth$ was trained
only on the forward trajectory (samples from $\px$ plus Gaussian noise), so too
large a $\gw$ drives $\xt$ where $\epsth$ extrapolates and the restoring force
becomes unreliable. Thus $\gw$ must be large enough to explore high-utility regions
yet small enough to keep $\xt$ where the score approximation is valid.
\end{remark}

\subsubsection{The rate guard}
At high noise the Tweedie estimate is coarse and can produce extreme grid points,
hence extreme rates that fall in the regime where the truncated utility is
inaccurate (Part~III, \S\ref{subsec:numerics}). The generator therefore zeroes the
guidance gradient whenever the predicted log-rate is too high:
\begin{equation}
\ggrad\leftarrow 0
\quad\text{whenever}\quad
\exp(\mug)\ge f_{\max},
\qquad
\mug=\gain\,\pmean+\latz,
\label{eq:fmax}
\end{equation}
with $f_{\max}=100$ and $\mug$ the predicted log-rate mean (gain $\gain$ times the
posterior latent mean $\pmean$ plus the offset $\latz$). This guard activates mostly in
the first $\sim\!10$ steps (high $\dt$) and keeps generation in the low-rate regime
where the (truncated) utility matches a Monte-Carlo evaluation. Because the utility
landscape over random seeds is sparse (most seeds land at $U\!\approx\!0$),
generation is run best-of-$K$: retain the highest-utility input over $K$ random
seeds.

% ----------------------------------------------------------------------------
\subsection{Open issues}\label{sec:diff-status}

Two remaining questions are mathematical rather than procedural.

\paragraph{Utility-accuracy ceiling at high rate.}
Both utilities inherit the Poisson-sum truncation at $\rmax$ analyzed in Part~III
(\S\ref{subsec:numerics}): above rate $\approx\!50$ the standard utility $\Ustd$
overflows to $+\infty$ (its analytic noise-entropy term is unbounded), while the
distribution-aware utility $\Uda$ stays finite but is silently wrong (its marginal
entropy collapses toward $0$). Guidance is therefore reliable only in the low-rate
regime, rate $\approx\!0.01$--$1$, where both utilities agree with a Monte-Carlo
evaluation to within $\sim\!3\%$; the rate guard \eqref{eq:fmax} is exactly what
confines generation to that regime. Extending accurate utility evaluation above the
truncation bound---so that guidance remains valid at high rate---is open.

\paragraph{Single-sample distribution-aware estimator.}
In the guided reverse path the distribution-aware utility \eqref{eq:Uda} is
evaluated with a single conditioning sample ($n_{\mathrm{mc}}=1$, the current
input) rather than a Monte-Carlo average over $\px$; enlarging $n_{\mathrm{mc}}$ is
a straightforward but untested refinement.


\appendix
% ============================ APPENDIX ============================

\section{Notation reference}\label{app:notation}

One symbol denotes one concept throughout. Entries marked \textbf{[clash]} are
symbols that collide across the literature and are disambiguated here.

\renewcommand{\arraystretch}{1.15}
\begin{longtable}{@{}p{0.42\textwidth} p{0.52\textwidth}@{}}
\toprule
\textbf{Symbol} & \textbf{Concept} \\
\midrule
\endhead
\multicolumn{2}{@{}l}{\textit{Model, kernel, structured prior}}\\
\midrule
$x$, $x^{*}$ & input (a real-valued signal on a 2-D grid); candidate input \\
$\pixc_i$, $\rfc$ & grid coordinate; center $(\varepsilon_{0x},\varepsilon_{0y})$ of the localized region \\
$\lat(x)$ & latent GP function \textbf{[clash: not $\latz$]} \\
$\latz$ & rate offset \\
$\rate(x)=e^{\gain\lat+\latz}$ & rate (Poisson intensity) \textbf{[clash: $\rate$ is the rate, never the latent]} \\
$\lgr=\gain\lat+\latz$ & log-rate (so $\rate=e^{\lgr}$) \\
$\gain$ & gain \textbf{[clash: not the projection $\amat$]} \\
$\spk$ ($\equiv r$) & count observation (synonyms $r,n$) \\
$\Kf(x,x')$ & arc-cosine prior kernel function \\
$\Ktil=\Kf(\Ipt,\Ipt)$ & inducing gram \textbf{[clash: never a bare $K$ for this]} \\
$\Kmag=\sqrt{\vmag_x\vmag_{x'}}$ & kernel magnitude \textbf{[clash: not $\Mind$]} \\
$\kang$, $\Jang(\kang)$ & kernel angle; angular term \textbf{[$\Jang$ is never a Jacobian]} \\
$\vmag_x=x\transpose\Cmat x+\sigz^{2}$ & kernel self-magnitude \\
$\Cmat$, $\Csm$ & structured (grid) prior covariance; RBF smoothing factor \\
$\loc_i$ & localization weight \textbf{[clash: not diffusion $\aret$]} \\
$\bwid$ & half-width of the localized region \textbf{[clash: not diffusion $\bsched$]} \\
$\smoo$ & smoothness SD \textbf{[clash: not correlation $\corr$]} \\
$\sigz$ & bias std, $\sigz=e^{\mathrm{raw\_sigma\_0}}$ \\
$\Amp$ & amplitude --- kernel hyperparameter, $\Amp=\mathrm{softplus}(\cdot)$, bounds $(0,1000]$, optimized \\
\midrule
\multicolumn{2}{@{}l}{\textit{Variational inference and eigenspace}}\\
\midrule
$\ipt$, $\Ipt$ & inducing point; inducing-point matrix \\
$\latt$ ($\equiv u$) & inducing latent values, $q(\latt)=\Gauss(\vm,\vV)$ \\
$\Mind$, $\Ntr$ & number of inducing points; number of training inputs ($\Mind{=}\Ntr$ in the loop) \\
$\kvec(x)$, $\pv(x)=\Ktil^{-1}\kvec(x)$ & cross-covariance vector; projection vector \\
$\vm$, $\vV$ & variational mean and covariance (inducing space) \\
$\pmean(x)$, $\pvar(x)$ & posterior mean and variance of the latent \\
$\Lchol$ & Cholesky factor of $\Ktil$ \textbf{[clash: $\ELBO$ is the ELBO]} \\
$\eb$, $\evals$, $\nb$ & eigenbasis of $\Ktil$; eigenvalues (diagonal); kept dimension \textbf{[$\evals$ not $\lat$]} \\
$\mb$, $\Vb$, $\Ktilb$ & eigenspace mean, dense covariance, diagonal inducing gram \\
$\amat=\Kf\Ktil^{-1}$ & projection matrix \\
$\fbar=e^{\gain\pmean+\frac12\gain^{2}\pvar+\latz}$ & expected rate \\
$\ELBO$ & evidence lower bound \textbf{[clash: script $\mathcal L$, not the Cholesky factor]} \\
\midrule
\multicolumn{2}{@{}l}{\textit{Training-step helpers}}\\
\midrule
$\gE=\gain\sum_i\kvec_i(\spk_i-\fbar_i)$ & E-step gradient helper \textbf{[clash: not the log-rate $\lgr$]} \\
$\GE=\gain^{2}\sum_i\fbar_i\kvec_i\kvec_i\transpose$ & E-step Fisher curvature (PSD) \\
\midrule
\multicolumn{2}{@{}l}{\textit{Utility and information}}\\
\midrule
$\Ustd$, $\Uda$ & standard and distribution-aware utility \\
$\Hmarg$, $\Hnoise$, $\Hcond$ & marginal, aleatoric-noise, and conditional entropy \textbf{[clash: never a bare $H$]} \\
$\px$ & input (data) distribution \\
$\mug=\gain\pmean+\latz$, $\varg=\gain^{2}\pvar$ & log-rate mean and variance \\
$\corr(x,x^{*})$ & posterior latent correlation \textbf{[clash: not smoothness $\smoo$; $\corr^{2}$ = fractional variance reduction]} \\
$\ccond$ & prior conditional covariance (the ``correction'' term) \\
$\scal$ & input scaling factor, $x^{*}=\scal\,x_t$ \textbf{[clash: not $\ccond$]} \\
$\varr$, $\subc$ & residual/aleatoric variance; subspace coordinate \textbf{[$\subc$ not an inducing point]} \\
$\LW$, $\gmode$, $\rmax$ & Lambert-$W$ (principal); Laplace mode; count truncation \\
$\satf$ & saturation-kernel amplitude \\
\midrule
\multicolumn{2}{@{}l}{\textit{Diffusion}}\\
\midrule
$\dt$ & diffusion timestep \\
$\xz$, $\xt$, $\xhat$ & clean, noisy, and Tweedie-estimated input \\
$\bsched$ & noise-schedule variance \textbf{[clash: not the prior width $\bwid$]} \\
$\aret$, $\abar$ & signal retention; cumulative retention \textbf{[clash: not the localization weight $\loc$]} \\
$\epsth$, $\score$ & U-Net noise predictor; score function \\
$\gw$, $\ggrad$ & guidance scale; guidance gradient ($=$ the utility gradient) \\
$\greg$ & regularization weight (input-optimization approach) \\
\bottomrule
\end{longtable}


\end{document}
